% Leçon 16 — Éléments socioculturels : maladies endémiques en Afrique et en Asie — Chinois Tle A — Cours signé OnBuch+
% Programme officiel MINESEC. Compiler avec : tectonic ZH16-socio-maladies-endemiques-en-afrique-et-en-asie.tex
\newcommand{\DOCMATIERE}{Chinois}
\newcommand{\DOCNIVEAU}{Tle A}
\newcommand{\DOCMODULE}{Module 2 : Environnement et santé}
\newcommand{\DOCLECON}{ZH16}
\newcommand{\DOCTITRE}{Leçon 16 — Éléments socioculturels : maladies endémiques en Afrique et en Asie}
% OnBuch+ — préambule « cours signés » ENRICHI (compilé avec tectonic / XeLaTeX)
% Système visuel « Pop » d'OnBuch+ : fond crème (jamais blanc), bordures encre
% épaisses, ombres « dures » décalées, titres Archivo Black en capitales.
% Version MATHÉMATIQUES : graphes (pgfplots/tikz), boîtes pédagogiques
% (définition, propriété, méthode, attention, le savais-tu, exercice résolu,
% exercice, TP, bilan), page de garde et signature OnBuch+ ancrée en bas.
%
% Macros attendues dans main.tex AVANT \input{preamble} :
%   \DOCMATIERE  \DOCNIVEAU  \DOCMODULE  \DOCLECON (numéro)  \DOCTITRE
\documentclass[11pt]{article}

\usepackage{fontspec}
\usepackage[french]{babel} % français = langue principale ; le chinois via \zh{..} / textebox (xeCJK)
\usepackage{xeCJK}
\usepackage{pifont} % \ding{51} \ding{55} utilisés par les modèles
\usepackage[a4paper,margin=2cm,top=3.4cm,bottom=2.8cm,headheight=1.4cm,footskip=1.3cm]{geometry}
\usepackage{amsmath,amssymb,amsfonts}
\usepackage{mathtools} % \xmapsto, \coloneqq... souvent utilisés par les modèles
\usepackage{cancel} % simplifications barrées (bilans de forces)
% \abs{z} (module, valeur absolue) et \norm{u} (norme), utilisés par les modèles
\providecommand{\abs}{}\let\abs\relax\DeclarePairedDelimiter{\abs}{\lvert}{\rvert}
\providecommand{\norm}{}\let\norm\relax\DeclarePairedDelimiter{\norm}{\lVert}{\rVert}
\usepackage[table]{xcolor} % [table] : fournit aussi \rowcolors (alternance de lignes)
\usepackage{pagecolor}
\usepackage{tikz}
\usetikzlibrary{arrows.meta,positioning,calc,shapes.geometric,shapes.misc,shapes.arrows,decorations.pathmorphing,decorations.pathreplacing,decorations.markings,patterns,shadows,mindmap,trees,fit,backgrounds,matrix,intersections,angles,quotes}
\usepackage{pgfplots}
\usepackage[version=4]{mhchem} % équations nucléaires \ce{^{235}_{92}U}
\usepackage[european,siunitx]{circuitikz} % schémas électriques
\pgfplotsset{compat=1.18}
\usepgfplotslibrary{fillbetween,groupplots} % aires entre courbes (calcul intégral)
\usepackage{enumitem}
\usepackage{fancyhdr}
% [most] charge toutes les bibliothèques tcolorbox (listings, minted, raster,
% documentation...) et gonfle inutilement la mémoire TeX sur un document long
% (~300 boîtes) : on ne charge que ce qui est réellement utilisé.
\usepackage[skins,breakable]{tcolorbox}
\usepackage{graphicx}
\usepackage{colortbl}
\usepackage{booktabs}
\usepackage{tabularx}
\usepackage{diagbox} % case d'en-tête diagonale des tableaux à double entrée
% Colonnes extensibles centrée / gauche / droite (C, L, R), souvent utilisées par les modèles
\newcolumntype{C}{>{\centering\arraybackslash}X}
\newcolumntype{L}{>{\raggedright\arraybackslash}X}
\newcolumntype{R}{>{\raggedleft\arraybackslash}X}
\usepackage{array}
\usepackage{multicol}
\usepackage{multirow}
\usepackage{siunitx}
% parse-numbers=false : accepte aussi \SI{2,5 \times 10^{-3}}{...}
\sisetup{locale=FR,output-decimal-marker={,},group-minimum-digits=4,parse-numbers=false}
\DeclareSIUnit\ohm{\ensuremath{\symup{\Omega}}} % U+2126 absent de la police
\DeclareSIUnit\division{div} % réglages d'oscilloscope (ms/div, V/div)
\DeclareSIUnit\tr{tr} % tours (vitesse de rotation, ex. \SI{1500}{\tr\per\minute})
\DeclareSIUnit\molar{M} % concentration molaire (ex. \SI{3}{\molar}, \SI{1}{\micro\molar})
\DeclareSIUnit\an{an} % année (le modèle confond parfois avec \year, un compteur TeX) — ex. \SI{5}{\kilo\meter\per\an}
\DeclareSIUnit\FCFA{FCFA} % franc CFA (ex. \SI{500}{\FCFA\per\kilo\gram})
\DeclareSIUnit\degreeBrix{\textdegree Brix} % teneur en sucre d'un sirop/jus (réfractométrie)
\DeclareSIUnit\month{mois} % le modèle utilise parfois \month, un compteur TeX, comme unité de durée
\usepackage{qrcode}
% \degree : commande fréquemment utilisée par les modèles pour un angle ou une
% température en dehors de \SI{}{} (non fournie par les paquets chargés).
\providecommand{\degree}{^{\circ}}
% \qed (fin de démonstration), utilisé par les modèles sans amsthm
\providecommand{\qed}{\hfill$\square$}
% \barre : double barre des valeurs interdites dans un tableau de variations (tkz-tab)
\providecommand{\barre}{\ensuremath{\|}}
% environnement proof (amsthm non chargé) utilisé par les modèles
\ifdefined\proof\else\newenvironment{proof}[1][Démonstration]{\par\noindent\textit{#1.}\ }{\qed\par}\fi
\usepackage{lastpage}
\usepackage{hyperref}
\hypersetup{hidelinks}

% ============================================================
% Palette « Pop » OnBuch+ — mêmes hex que la classe P (plus_ui.dart)
% ============================================================
\definecolor{poporange}{HTML}{F59321}
\definecolor{popdark}{HTML}{E07A0C}
\definecolor{popcream}{HTML}{FFF6E8}
\definecolor{popink}{HTML}{1C1714}
\definecolor{poppurple}{HTML}{7A5AE0}
\definecolor{popgreen}{HTML}{2E9E6A}
\definecolor{poppink}{HTML}{E0457B}
\definecolor{popblue}{HTML}{2F7DE1}
\definecolor{popgold}{HTML}{C98A12}
\definecolor{popmuted}{HTML}{8A7F76}
\definecolor{popline}{HTML}{EADFD2}
\definecolor{popgray}{HTML}{8A7F76}
\colorlet{popgrey}{popgray}
% Alias tolérants (le modèle nomme parfois les couleurs différemment)
\colorlet{popwhite}{white}
\colorlet{popblack}{popink}
\colorlet{popred}{poppink}
\colorlet{popyellow}{popgold}
% Teintes claires (fonds de tableaux/encadrés) — le modèle nomme parfois
% les couleurs avec un suffixe L (ex. popblueL) sans les avoir définies.
\colorlet{poporangeL}{poporange!20}
\colorlet{popdarkL}{popdark!20}
\colorlet{poppurpleL}{poppurple!20}
\colorlet{popgreenL}{popgreen!20}
\colorlet{poppinkL}{poppink!20}
\colorlet{popblueL}{popblue!20}
\colorlet{popgoldL}{popgold!20}
\colorlet{popredL}{popred!20}
\colorlet{popgrayL}{popgray!20}
% Teintes claires pour fonds de tableaux / zones
\colorlet{popblueL}{popblue!10!white}
\colorlet{poporangeL}{poporange!14!white}
\colorlet{popgreenL}{popgreen!12!white}
\colorlet{poppurpleL}{poppurple!12!white}
\colorlet{poppinkL}{poppink!12!white}
\colorlet{popgoldL}{popgold!14!white}

\pagecolor{popcream}

% ============================================================
% Typographie
% ============================================================
\defaultfontfeatures{Ligatures=TeX}
\setmainfont{PlusJakartaSans-Medium.ttf}[Path=../../../fonts/,BoldFont=PlusJakartaSans-Bold.ttf]
% Symboles, API et emojis absents de la police principale : repli sur DejaVu Sans (caractères actifs)
\newfontface\symfont{DejaVuSans.ttf}[Path=../../../fonts/]
\makeatletter
\newcommand\symactive[1]{\catcode#1=13 \begingroup\lccode`\~=#1 \lowercase{\endgroup\def~}{{\symfont\char#1}}}
\newcommand\symblank[1]{\catcode#1=13 \begingroup\lccode`\~=#1 \lowercase{\endgroup\def~}{}}
\newcommand\symhyph[1]{\catcode#1=13 \begingroup\lccode`\~=#1 \lowercase{\endgroup\def~}{\mbox{-}}}
\newcount\symc
\newcommand\symrange[2]{\symc=#1 \@whilenum\symc<#2 \do{\expandafter\symactive\expandafter{\the\symc}\advance\symc by 1 }}
\newcommand\symblankrange[2]{\symc=#1 \@whilenum\symc<#2 \do{\expandafter\symblank\expandafter{\the\symc}\advance\symc by 1 }}
\makeatother
\symrange{"0250}{"02FF}\symactive{"2713}\symactive{"2714}\symactive{"2717}\symactive{"2718}\symactive{"2610}\symactive{"2611}\symactive{"2612}
\symrange{"2600}{"27BF}
\catcode"2705=13 \begingroup\lccode`\~="2705 \lowercase{\endgroup\def~}{{\symfont\char"2713}}\symblank{"26BD}\symblank{"26A1}
\symrange{"2460}{"257F}\symactive{"223C}\symactive{"2248}\symactive{"2260}
\symhyph{"2011}\symhyph{"2E1A}\symblankrange{"1F300}{"1FAFF}\symblank{"FE0F}
% Caractères chinois : WenQuanYi Zen Hei (sous-ensemble GB2312 + pinyin), gras/italique simulés
\setCJKmainfont{WenQuanYiZenHei-subset.ttf}[Path=../../../fonts/,AutoFakeBold=3,AutoFakeSlant=0.2]
\xeCJKsetup{CJKspace=false}
% Pinyin : voyelles à caron (ǎ ǐ ǒ ǔ ǖ ǘ ǚ ǜ) absentes de la police principale -> caractères actifs
\newfontface\pyfont{WenQuanYiZenHei-subset.ttf}[Path=../../../fonts/]
\newcommand\pyactive[1]{\catcode#1=13 \begingroup\lccode`\~=#1 \lowercase{\endgroup\def~}{{\pyfont\char#1}}}
\pyactive{"01CD}\pyactive{"01CE}\pyactive{"01CF}\pyactive{"01D0}\pyactive{"01D1}\pyactive{"01D2}\pyactive{"01D3}\pyactive{"01D4}\pyactive{"01D5}\pyactive{"01D6}\pyactive{"01D7}\pyactive{"01D8}\pyactive{"01D9}\pyactive{"01DA}\pyactive{"01DB}\pyactive{"01DC}
% Maths en sans-serif (Fira Math) avec chiffres et lettres droites de Plus
% Jakarta, pour que les formules se fondent dans le texte.
\usepackage[math-style=ISO,bold-style=ISO]{unicode-math}
\setmathfont{FiraMath-Regular.otf}[Path=../../../fonts/]
\setmathfont{PlusJakartaSans-Medium.ttf}[Path=../../../fonts/,range={up/num,up/{Latin,latin},\mathup}]
\usepackage{icomma} % virgule décimale sans espace en mode math
% Caractères mathématiques Unicode absents des polices de texte (Plus Jakarta,
% Archivo Black) : rendus par leur équivalent mathématique, en texte comme en
% formule — sinon ils s'affichent en carré vide (ex. « Arithmétique dans ℤ »).
\usepackage{newunicodechar}
\newunicodechar{ℝ}{\ensuremath{\mathbb{R}}}
\newunicodechar{ℤ}{\ensuremath{\mathbb{Z}}}
\newunicodechar{ℂ}{\ensuremath{\mathbb{C}}}
\newunicodechar{ℕ}{\ensuremath{\mathbb{N}}}
\newunicodechar{ℚ}{\ensuremath{\mathbb{Q}}}
\newunicodechar{𝔻}{\ensuremath{\mathbb{D}}}
\newunicodechar{α}{\ensuremath{\alpha}}
\newunicodechar{β}{\ensuremath{\beta}}
\newunicodechar{γ}{\ensuremath{\gamma}}
\newunicodechar{δ}{\ensuremath{\delta}}
\newunicodechar{ε}{\ensuremath{\varepsilon}}
\newunicodechar{θ}{\ensuremath{\theta}}
\newunicodechar{λ}{\ensuremath{\lambda}}
\newunicodechar{μ}{\ensuremath{\mu}}
\newunicodechar{σ}{\ensuremath{\sigma}}
\newunicodechar{τ}{\ensuremath{\tau}}
\newunicodechar{φ}{\ensuremath{\varphi}}
\newunicodechar{ψ}{\ensuremath{\psi}}
\newunicodechar{ω}{\ensuremath{\omega}}
\newunicodechar{Σ}{\ensuremath{\Sigma}}
% majuscules produites par \MakeUppercase dans les titres (α-aminés -> Α)
\newunicodechar{Α}{\ensuremath{\alpha}}
\newunicodechar{Β}{\ensuremath{\beta}}
\newunicodechar{Γ}{\ensuremath{\Gamma}}
\newunicodechar{Δ}{\ensuremath{\Delta}}
\newunicodechar{Ε}{\ensuremath{\varepsilon}}
\newunicodechar{Θ}{\ensuremath{\Theta}}
\newunicodechar{Λ}{\ensuremath{\Lambda}}
\newunicodechar{Μ}{\ensuremath{\mu}}
\newunicodechar{Τ}{\ensuremath{\tau}}
\newunicodechar{Φ}{\ensuremath{\Phi}}
\newunicodechar{Ψ}{\ensuremath{\Psi}}
\newunicodechar{Ω}{\ensuremath{\Omega}}
\newunicodechar{↦}{\ensuremath{\mapsto}}
\newunicodechar{∈}{\ensuremath{\in}}
\newunicodechar{∉}{\ensuremath{\notin}}
\newunicodechar{∧}{\ensuremath{\wedge}}
\newunicodechar{⇔}{\ensuremath{\Leftrightarrow}}
\newunicodechar{⟺}{\ensuremath{\Longleftrightarrow}}
\newunicodechar{⇒}{\ensuremath{\Rightarrow}}
\newunicodechar{⟹}{\ensuremath{\Longrightarrow}}
\newunicodechar{∘}{\ensuremath{\circ}}
\newunicodechar{∪}{\ensuremath{\cup}}
\newunicodechar{∩}{\ensuremath{\cap}}
\newunicodechar{≡}{\ensuremath{\equiv}}
\newunicodechar{⊂}{\ensuremath{\subset}}
\newunicodechar{⊥}{\ensuremath{\perp}}
\newunicodechar{∃}{\ensuremath{\exists}}
\newunicodechar{∀}{\ensuremath{\forall}}
\newunicodechar{∛}{\ensuremath{\sqrt[3]{\,}}}
\newunicodechar{∜}{\ensuremath{\sqrt[4]{\,}}}
\newunicodechar{⃗}{\ensuremath{{}^{\to}}}
\newunicodechar{ⁿ}{\ifmmode^{n}\else\textsuperscript{\ensuremath{n}}\fi}
\newunicodechar{ˣ}{\ifmmode^{x}\else\textsuperscript{\ensuremath{x}}\fi}
\newunicodechar{ᵇ}{\ifmmode^{b}\else\textsuperscript{\ensuremath{b}}\fi}
\newunicodechar{ʳ}{\ifmmode^{r}\else\textsuperscript{\ensuremath{r}}\fi}
\newunicodechar{ᵘ}{\ifmmode^{u}\else\textsuperscript{\ensuremath{u}}\fi}
\newunicodechar{ʸ}{\ifmmode^{y}\else\textsuperscript{\ensuremath{y}}\fi}
\newunicodechar{ᵃ}{\ifmmode^{a}\else\textsuperscript{\ensuremath{a}}\fi}
\newunicodechar{ᵉ}{\ifmmode^{e}\else\textsuperscript{\ensuremath{e}}\fi}
\newunicodechar{ᵏ}{\ifmmode^{k}\else\textsuperscript{\ensuremath{k}}\fi}
\newunicodechar{ᵖ}{\ifmmode^{p}\else\textsuperscript{\ensuremath{p}}\fi}
\newunicodechar{ᴺ}{\ifmmode^{N}\else\textsuperscript{\ensuremath{N}}\fi}
\newunicodechar{ᶜ}{\ifmmode^{c}\else\textsuperscript{\ensuremath{c}}\fi}
\newunicodechar{ʰ}{\ifmmode^{h}\else\textsuperscript{\ensuremath{h}}\fi}
\newunicodechar{ᵠ}{\ifmmode^{q}\else\textsuperscript{\ensuremath{q}}\fi}
\newunicodechar{ₙ}{\ifmmode_{n}\else\textsubscript{\ensuremath{n}}\fi}
\newunicodechar{ₐ}{\ifmmode_{a}\else\textsubscript{\ensuremath{a}}\fi}
\newunicodechar{ᵢ}{\ifmmode_{i}\else\textsubscript{\ensuremath{i}}\fi}
\newunicodechar{ᵣ}{\ifmmode_{r}\else\textsubscript{\ensuremath{r}}\fi}
\newunicodechar{ₖ}{\ifmmode_{k}\else\textsubscript{\ensuremath{k}}\fi}
\newunicodechar{ᵦ}{\ifmmode_{\beta}\else\textsubscript{\ensuremath{\beta}}\fi}
\newunicodechar{ₓ}{\ifmmode_{x}\else\textsubscript{\ensuremath{x}}\fi}
\newfontfamily\popdisplay{ArchivoBlack-Regular.ttf}[Path=../../../fonts/]
\newfontfamily\poplogo{SpaceGrotesk-Bold.ttf}[Path=../../../fonts/]
\newfontfamily\popsemi{PlusJakartaSans-SemiBold.ttf}[Path=../../../fonts/]
\newfontfamily\popxbold{PlusJakartaSans-ExtraBold.ttf}[Path=../../../fonts/]
\newfontfamily\popxboldls{PlusJakartaSans-ExtraBold.ttf}[Path=../../../fonts/,LetterSpace=3.0]

\setlist[enumerate]{leftmargin=1.7em,itemsep=5pt,topsep=4pt}
\setlist[itemize]{leftmargin=1.7em,itemsep=4pt,topsep=4pt,label={\color{poporange}\textbullet}}
\setlength{\parindent}{0pt}
\setlength{\parskip}{5pt}
\renewcommand{\arraystretch}{1.25}

% pgfplots : style Pop par défaut
\pgfplotsset{
  every axis/.append style={axis line style={popink,line width=1pt},tick style={popink},
    grid style={popline,line width=0.5pt},label style={font=\small},tick label style={font=\footnotesize}},
  popcurve/.style={poporange,line width=1.8pt,no markers,smooth},
}
% Même style disponible dans un \draw TikZ simple (hors pgfplots)
\tikzset{popcurve/.style={poporange,line width=1.8pt}}
\tikzset{popgrid/.style={popline,very thin}}
\colorlet{popcurve}{poporange} % le nom du style est parfois utilisé comme couleur

% ============================================================
% Pill / squiggle
% ============================================================
\newcommand{\poppill}[3]{%
  \tikz[baseline=-0.55ex]\node[draw=popink,line width=1.2pt,rounded corners=2.6pt,%
    fill=#2,inner xsep=6.5pt,inner ysep=3pt]{%
    \popxboldls\fontsize{7}{7}\selectfont\color{#3}\MakeUppercase{#1}};%
}
\newcommand{\popsquiggle}[1]{%
  \tikz[baseline=-0.4ex]{
    \draw[#1,line width=2.6pt,line cap=round]
      plot [smooth] coordinates {(0,0.09) (0.34,0.01) (0.68,0.21) (1.02,0.01) (1.36,0.10)};
  }%
}
% Mot-clé surligné (marqueur orange sous le texte)
\newcommand{\cle}[1]{{\popxbold\color{popdark}#1}}

% ============================================================
% En-tête / pied
% ============================================================
\pagestyle{fancy}
\fancyhf{}
\renewcommand{\headrulewidth}{0pt}
\renewcommand{\footrulewidth}{0pt}
\lhead{\raisebox{-0.15ex}{\poplogo\fontsize{13}{13}\selectfont\color{popink}OnBuch\color{poporange}+}}
\rhead{\poppill{\DOCMATIERE\ \textbullet\ \DOCNIVEAU\ \textbullet\ Leçon \DOCLECON}{white}{popink}}
\lfoot{\popxbold\fontsize{7.6}{7.6}\selectfont\color{popmuted}\MakeUppercase{Cours signé OnBuch+}\ \textbullet\ {\popsemi\DOCTITRE}}
\rfoot{\tikz[baseline=-0.6ex]\node[rounded corners=8.5pt,fill=popink,minimum height=17pt,minimum width=17pt,inner xsep=5pt,inner sep=0]{\popxbold\fontsize{8}{8}\selectfont\color{popcream}\thepage\,/\,\pageref*{LastPage}};}

% ============================================================
% Titres
% ============================================================
\newcommand{\chaptitre}[2]{%
  \begin{tcolorbox}[enhanced,colback=poporange,colframe=popink,boxrule=2.6pt,arc=11pt,
      drop shadow=popink,left=15pt,right=15pt,top=12pt,bottom=11pt,before skip=3pt,after skip=18pt]
    \poppill{#2}{popink}{popcream}\par\vspace{9pt}
    {\raggedright\hyphenpenalty=10000\exhyphenpenalty=10000\popdisplay\fontsize{22}{24}\selectfont\color{popcream}\MakeUppercase{#1}\par}
  \end{tcolorbox}
}
% Section numérotée : pastille orange avec le numéro + titre Archivo
\newcounter{popsec}
\newcommand{\coursec}[1]{%
  \stepcounter{popsec}\setcounter{popsub}{0}%
  \par\vspace{16pt}\noindent
  \begin{minipage}{\linewidth}
  \tikz[baseline=-0.7ex]\node[draw=popink,line width=1.6pt,fill=poporange,rounded corners=4pt,minimum size=22pt,inner sep=0]{\popdisplay\fontsize{12}{12}\selectfont\color{popcream}\thepopsec};%
  \hspace{8pt}{\popdisplay\fontsize{15}{17}\selectfont\color{popink}\MakeUppercase{#1}}\par
  \vspace{4pt}\noindent{\color{poporange}\rule{36pt}{3.6pt}}
  \end{minipage}\par\nopagebreak\vspace{8pt}
}
\newcounter{popsub}
\newcommand{\courssub}[1]{%
  \stepcounter{popsub}%
  \par\vspace{9pt}\noindent{\popxbold\fontsize{11.5}{13}\selectfont\color{popink}{\color{poporange}\thepopsec.\thepopsub}\hspace{6pt}#1}\par\nopagebreak\vspace{4pt}
}

% ============================================================
% Boîtes pédagogiques — même recette : fond blanc, bordure épaisse colorée,
% ombre dure encre, badge plein. Titre optionnel [#1].
% ============================================================
\tcbset{popbase/.style={breakable,enhanced,colback=white,boxrule=2.3pt,arc=9pt,
  shadow={3pt}{-3pt}{0pt}{popink},left=14pt,right=14pt,top=11pt,bottom=11pt,before skip=11pt,after skip=15pt,
  fontupper=\popsemi\fontsize{10.6}{14.5}\selectfont\color{popink},
  fontlower=\popsemi\fontsize{10.6}{14.5}\selectfont\color{popink}}}
% Coche dessinée (le glyphe ✓ n'existe pas dans les polices du document)
\newcommand{\popcheck}[1]{\tikz[baseline=-0.5ex]\draw[#1,line width=1.6pt,line cap=round,line join=round](-0.13,0.0)--(-0.04,-0.09)--(0.13,0.1);}
\providecommand{\coche}{\popcheck{popgreen}}
\providecommand{\resultat}[1]{\par\smallskip\noindent\fcolorbox{popgreen}{popgreenL}{\parbox{\dimexpr\linewidth-2\fboxsep-2\fboxrule\relax}{\textbf{Résultat.} #1}}\par\smallskip}
\newcommand{\popboxhead}[3]{\poppill{#1}{#2}{white}\ifx\relax#3\relax\else\hspace{6pt}{\popxbold\fontsize{10.6}{12}\selectfont\color{popink}#3}\fi\par\vspace{7pt}}

\newtcolorbox{aretenir}[1][]{popbase,colframe=popgreen,before upper={\popboxhead{À retenir}{popgreen}{#1}}}
% Texte en chinois (document de lecture, dialogue, extrait) : cadre
% d'étude. Un titre optionnel (source, auteur) s'affiche dans l'en-tête.
\newtcolorbox{textebox}[1][]{popbase,colframe=popblue,before upper={\popboxhead{文本 · Texte}{popblue}{#1}},after upper={}}
% Marques ✓ / ✗ (forme correcte / fautive) souvent utilisées par les modèles
\providecommand{\cross}{\textcolor{poppink}{\ensuremath{\times}}}
\providecommand{\xmark}{\cross}\providecommand{\crossmark}{\cross}\providecommand{\cmark}{\ensuremath{\checkmark}}
% Ligne de réponse / justification à compléter (inventée par les modèles)
\providecommand{\justification}[1]{\par\noindent\textit{Justification~:} \dotfill\par}
% \textipa (package tipa non chargé) : la police ne couvre pas l'API -> simple passage
\providecommand{\textipa}[1]{#1}
% Soulignement (ulem non chargé) : \uline, \ul
\providecommand{\uline}[1]{\underline{#1}}\providecommand{\ul}[1]{\underline{#1}}
% \glossaire{\item ...} inventée par les modèles : liste de glossaire
\providecommand{\glossaire}[1]{\begin{itemize}#1\end{itemize}}
% \alert (beamer) inventée par les modèles : texte mis en évidence
\providecommand{\alert}[1]{\textbf{\textcolor{poppink}{#1}}}
% variantes ASCII d'ordinaux écrites en macro
\providecommand{\iere}{\textsuperscript{re}}\providecommand{\ieme}{\textsuperscript{e}}\providecommand{\ere}{\textsuperscript{re}}\providecommand{\eme}{\textsuperscript{e}}\providecommand{\er}{\textsuperscript{er}}\providecommand{\ie}{\textsuperscript{e}}
% Ordinaux français écrits en macro par les modèles : 1\ière, 3\ième
\newcommand{\ière}{\textsuperscript{re}}\newcommand{\ième}{\textsuperscript{e}}
% \exemple (inventée par les modèles) : étiquette « Exemple : »
\providecommand{\exemple}{\textbf{Exemple~:}\ }
% \square utilisable aussi hors mode math (cases à cocher)
\AtBeginDocument{\let\squareMath\square\renewcommand{\square}{\ensuremath{\squareMath}}}
% Mot ou phrase en chinois à l'intérieur d'un paragraphe français
\newcommand{\zh}[1]{\ifmmode\text{#1}\else{#1}\fi}
\let\eng\zh \let\esp\zh \let\all\zh % alias tolérés
% flèches d'intonation inventées par les modèles
\providecommand{\textsearrow}{}\renewcommand{\textsearrow}{\ensuremath{\searrow}}\providecommand{\textnearrow}{}\renewcommand{\textnearrow}{\ensuremath{\nearrow}}
% \fr{..} : mot français (inventé par les modèles)
\providecommand{\fr}[1]{\textit{#1}}
% zhbox : encadré de texte chinois inventé par les modèles
\newenvironment{zhbox}{\begin{quote}}{\end{quote}}
% \vs : « versus » inventé par les modèles
\providecommand{\vs}{\textit{vs}\ }
% \blank : case à compléter inventée par les modèles
\providecommand{\blank}[1][]{\underline{\hspace{1.6em}}}
\newtcolorbox{exemplebox}[1][]{popbase,colframe=poporange,before upper={\popboxhead{Exemple}{poporange}{#1}}}
\newtcolorbox{definition}[1][]{popbase,colframe=popblue,before upper={\popboxhead{Définition}{popblue}{#1}}}
\newtcolorbox{propriete}[1][]{popbase,colframe=poppurple,before upper={\popboxhead{Propriété}{poppurple}{#1}}}
\newtcolorbox{methode}[1][]{popbase,colframe=popgold,before upper={\popboxhead{Méthode}{popgold}{#1}}}
\newtcolorbox{attention}[1][]{popbase,colframe=poppink,before upper={\popboxhead{Attention}{poppink}{#1}}}
\newtcolorbox{experience}[1][]{popbase,colframe=popblue,colback=popblueL,before upper={\popboxhead{Expérience}{popblue}{#1}}}
% « Le savais-tu ? » : carte violette pleine (contexte, culture, Cameroun)
\newtcolorbox{savaistu}[1][]{popbase,colback=poppurple,colframe=popink,
  fontupper=\popsemi\fontsize{10.4}{14.2}\selectfont\color{white},
  before upper={\poppill{Le savais-tu\,?}{white}{poppurple}\ifx\relax#1\relax\else\hspace{6pt}{\popxbold\color{white}#1}\fi\par\vspace{7pt}}}
% Exercice résolu : énoncé en haut, \tcblower puis la solution sur fond crème
\newcounter{popexo}\newcounter{popexr}
\newtcolorbox{exoresolu}[1][]{popbase,colframe=poporange,colbacklower=popcream,
  segmentation style={popink,line width=1.2pt,dashed},
  before upper={\stepcounter{popexr}\popboxhead{Exercice résolu \thepopexr}{poporange}{#1}},
  before lower={\poppill{Solution}{popink}{popcream}\par\vspace{6pt}}}
% Exercice (non résolu en place) — la correction est dans la partie Corrigés
\newtcolorbox{exercice}[1][]{popbase,colframe=popink,
  before upper={\stepcounter{popexo}\popboxhead{Exercice \thepopexo}{popink}{#1}}}
% Corrigé
\newtcolorbox{corrige}[1][]{popbase,colframe=popgreen,colback=popgreenL,
  before upper={\popboxhead{Corrigé}{popgreen}{#1}}}
% Objectifs / prérequis / bilan
\newtcolorbox{objectifs}{popbase,colframe=popink,colback=poporangeL,
  before upper={\popboxhead{Objectifs de la leçon}{popink}{}}}
\newtcolorbox{prerequis}{popbase,colframe=popink,colback=popblueL,
  before upper={\popboxhead{Prérequis}{popblue}{}}}
\newtcolorbox{bilan}{popbase,colframe=popink,colback=white,boxrule=2.8pt,
  before upper={\popboxhead{Fiche bilan}{poporange}{L'essentiel à connaître pour l'examen}}}
% Figure encadrée avec légende
\newtcolorbox{popfigure}[1][]{enhanced,breakable=false,colback=white,colframe=popline,boxrule=1.4pt,arc=8pt,
  left=10pt,right=10pt,top=10pt,bottom=8pt,before skip=10pt,after skip=12pt,halign=center,
  lower separated=false,halign lower=center,
  fontlower=\popsemi\fontsize{9}{11.5}\selectfont\color{popmuted},
  #1}
\newcommand{\legende}[1]{\par\vspace{6pt}{\popsemi\fontsize{9}{11.5}\selectfont\color{popmuted}#1}}

% ============================================================
% Signature OnBuch+
% ============================================================
\newcommand{\poplogoinline}[1]{{\poplogo\fontsize{#1}{#1}\selectfont\color{popink}OnBuch\color{poporange}+}}
\newcommand{\signatureonbuch}{%
  \begin{tcolorbox}[enhanced,colback=popink,colframe=popink,boxrule=2.2pt,arc=10pt,
      drop shadow=poporange,left=14pt,right=14pt,top=10pt,bottom=10pt,before skip=0pt,after skip=0pt]
    \tikz[baseline=-0.7ex]\node[circle,fill=popgreen,draw=popcream,line width=1.3pt,minimum size=22pt,inner sep=0]{\popcheck{white}};%
    \hspace{9pt}\begin{minipage}[c]{0.62\linewidth}
      {\popxbold\fontsize{12}{13}\selectfont\color{popcream}Signé\ \textbullet\ {\poplogo OnBuch\color{poporange}+}}\par\vspace{2pt}
      {\raggedright\popsemi\fontsize{8}{10.5}\selectfont\color{popline}\MakeUppercase{Équipe pédagogique OnBuch+}\\\MakeUppercase{Conforme au programme officiel MINESEC}\par}
    \end{minipage}\hfill
    {\poplogo\fontsize{22}{22}\selectfont\color{popcream}OnBuch\color{poporange}+}
  \end{tcolorbox}%
}

% ============================================================
% Page de garde
% #1 titre · #2 matière · #3 niveau · #4 module · #5 n° leçon · #6 durée ·
% #7 description / objectifs (texte ou itemize)
% ============================================================
\newcommand{\pagegarde}[7]{%
  \thispagestyle{empty}
  \newgeometry{a4paper,margin=1.7cm,top=1.6cm,bottom=1.4cm}%
  \begin{tikzpicture}[remember picture,overlay]
    \fill[poporange!18] ([xshift=-3.2cm,yshift=-3.6cm]current page.north east) circle (4.6cm);
    \fill[poppurple!14] ([xshift=2.4cm,yshift=7.6cm]current page.south west) circle (3.8cm);
  \end{tikzpicture}%
  {\poplogo\fontsize{30}{30}\selectfont\color{popink}OnBuch\color{poporange}+}\hfill
  \raisebox{0.6ex}{\poppill{Cours signé \textbullet\ Premium}{popink}{popcream}}\par
  \vspace{1pt}\popsquiggle{poppurple}\par
  \vspace{8pt}
  \begin{tcolorbox}[enhanced,colback=poporange,colframe=popink,boxrule=2.8pt,arc=13pt,
      drop shadow=popink,left=18pt,right=18pt,top=12pt,bottom=13pt,after skip=10pt]
    \poppill{#2 \textbullet\ #3}{popink}{popcream}\hspace{5pt}\poppill{#4}{white}{popink}\par\vspace{8pt}
    {\popdisplay\fontsize{14}{14}\selectfont\color{popink}LEÇON #5}\par\vspace{4pt}
    {\raggedright\hyphenpenalty=10000\exhyphenpenalty=10000\popdisplay\fontsize{25}{27}\selectfont\color{popcream}\MakeUppercase{#1}\par}
    \vspace{8pt}
    {\popxbold\fontsize{9.5}{9.5}\selectfont\color{popink}\MakeUppercase{Durée conseillée : #6}}
  \end{tcolorbox}
  \begin{tcolorbox}[enhanced,colback=white,colframe=popink,boxrule=2.2pt,arc=10pt,
      drop shadow=popink,left=16pt,right=16pt,top=10pt,bottom=10pt,after skip=8pt]
    \poppill{Dans cette leçon}{poppurple}{white}\par\vspace{5pt}
    \popsemi\fontsize{9.3}{12.6}\selectfont\color{popink}#7
  \end{tcolorbox}
  \noindent\begin{minipage}[t]{0.487\linewidth}
    \begin{tcolorbox}[enhanced,colback=popgreen,colframe=popink,boxrule=2.2pt,arc=10pt,
        drop shadow=popink,left=11pt,right=11pt,top=10pt,bottom=10pt,height=3.1cm,valign=center]
      \tcbox[colback=white,colframe=popink,boxrule=1.3pt,arc=5pt,boxsep=0pt,nobeforeafter,
        left=3pt,right=3pt,top=3pt,bottom=3pt,valign=center]{\qrcode[height=1.9cm]{https://play.google.com/store/apps/details?id=com.luvvix.onbuch&hl=fr}}%
      \hspace{8pt}\begin{minipage}[c]{0.5\linewidth}
        {\popdisplay\fontsize{11}{12.5}\selectfont\color{white}\MakeUppercase{Télécharge OnBuch}}\par\vspace{4pt}
        {\raggedright\popsemi\fontsize{8.4}{11}\selectfont\color{white}Gratuit \textbullet\ Google Play\\Tuteur IA, annales, concours, résultats\par}
      \end{minipage}
    \end{tcolorbox}
  \end{minipage}\hfill
  \begin{minipage}[t]{0.487\linewidth}
    \begin{tcolorbox}[enhanced,colback=poppurple,colframe=popink,boxrule=2.2pt,arc=10pt,
        drop shadow=popink,left=11pt,right=11pt,top=10pt,bottom=10pt,height=3.1cm,valign=center]
      \tikz[baseline=-0.7ex]\node[circle,fill=white,draw=popink,line width=1.3pt,minimum size=1.6cm,inner sep=0]{\popdisplay\fontsize{24}{24}\selectfont\color{poppurple}+};%
      \hspace{8pt}\begin{minipage}[c]{0.6\linewidth}
        {\popdisplay\fontsize{11}{12.5}\selectfont\color{white}\MakeUppercase{Passe à OnBuch+}}\par\vspace{4pt}
        {\raggedright\popsemi\fontsize{8.4}{11}\selectfont\color{white}Tous les cours signés, Tuteur IA illimité, fiches PDF — onglet OnBuch+ de l'app\par}
      \end{minipage}
    \end{tcolorbox}
  \end{minipage}
  \vspace{8pt}
  \popcontenu
  \vfill
  \signatureonbuch
  \restoregeometry
  \newpage
}

% Rangée « Ce cours contient » (page de garde)
\newcommand{\poptile}[3]{% #1 couleur #2 gros texte #3 légende
  \begin{tcolorbox}[enhanced,colback=#1,colframe=popink,boxrule=2pt,arc=9pt,shadow={2.5pt}{-2.5pt}{0pt}{popink},
      left=6pt,right=6pt,top=9pt,bottom=9pt,halign=center,valign=center,height=1.75cm,width=\linewidth,nobeforeafter]
    {\popdisplay\fontsize{12.5}{13}\selectfont\color{white}\MakeUppercase{#2}}\par\vspace{3pt}
    {\centering\popsemi\fontsize{7.8}{9.5}\selectfont\color{white}#3\par}
  \end{tcolorbox}}
\newcommand{\popcontenu}{%
  \par\noindent{\popxboldls\fontsize{7.5}{7.5}\selectfont\color{popmuted}CE COURS SIGNÉ CONTIENT}\par\vspace{7pt}
  \noindent\begin{minipage}[t]{0.235\linewidth}\poptile{popblue}{Cours}{complet, illustré, pas à pas}\end{minipage}\hfill
  \begin{minipage}[t]{0.235\linewidth}\poptile{poporange}{Méthodes}{et exercices résolus}\end{minipage}\hfill
  \begin{minipage}[t]{0.235\linewidth}\poptile{poppink}{Exercices}{type examen + corrigés}\end{minipage}\hfill
  \begin{minipage}[t]{0.235\linewidth}\poptile{popgreen}{Fiche bilan}{l'essentiel à retenir}\end{minipage}\par}

% Sommaire
\newtcolorbox{sommaire}{enhanced,colback=white,colframe=popink,boxrule=2.2pt,arc=10pt,
  drop shadow=popink,left=18pt,right=18pt,top=13pt,bottom=13pt,before skip=6pt,after skip=20pt,
  before upper={\poppill{Sommaire}{poppurple}{white}\par\vspace{9pt}\popsemi\fontsize{10.6}{14.5}\selectfont\color{popink}}}

% Signature de fin de cours — poussée tout en bas de la dernière page
\newcommand{\findecours}{%
  \par\vfill\nopagebreak
  \begin{center}\popsquiggle{poporange}\end{center}\vspace{4pt}
  \signatureonbuch
}
\AtBeginDocument{\def\llbracket{\mathopen{[\mkern-3.5mu[}}\def\rrbracket{\mathclose{]\mkern-3.5mu]}}} % intervalles d'entiers (glyphe ⟦ absent de la police)
% Glyphes absents de Fira Math : redéfinis à partir de glyphes disponibles
\AtBeginDocument{%
  \def\setminus{\mathbin{\backslash}}\let\smallsetminus\setminus
  \def\vdots{\mathord{\vbox{\baselineskip3.2pt\lineskiplimit0pt\kern4pt\hbox{.}\hbox{.}\hbox{.}}}}%
  \def\ddots{\mathinner{\mkern1mu\raise6.5pt\hbox{.}\mkern2mu\raise3.6pt\hbox{.}\mkern2mu\raise0.7pt\hbox{.}\mkern1mu}}%
  \def\triangle{\mathord{\scalebox{0.9}{$\symup{\Delta}$}}}%
  \def\nabla{\mathord{\raisebox{\depth}{\scalebox{1}[-1]{$\symup{\Delta}$}}}}%
  \def\checkmark{\popcheck{popdark}}%
  \def\star{\mathbin{\ast}}\def\bigstar{\mathord{\ast}}%
  \def\diamond{\mathbin{\scalebox{0.6}{$\Diamond$}}}%
  \def\bigcup{\mathop{\mathchoice{\raisebox{-0.3em}{\scalebox{1.7}{$\cup$}}}{\scalebox{1.25}{$\cup$}}{\cup}{\cup}}\displaylimits}%
  \def\bigcap{\mathop{\mathchoice{\raisebox{-0.3em}{\scalebox{1.7}{$\cap$}}}{\scalebox{1.25}{$\cap$}}{\cap}{\cap}}\displaylimits}%
  \def\lesseqqgtr{\mathrel{\lessgtr}}\def\bigoplus{\mathop{\oplus}\displaylimits}%
}
\providecommand{\ssi}{si et seulement si} % abréviation fréquente
\AtBeginDocument{\providecommand{\wideparen}{\overparen}} % arc de cercle

%% FIGFIX-BEGIN v5 — correctifs globaux des figures (docs/onbuch-generation/tools/figfix)
\usepackage{adjustbox}\usepackage{etoolbox}\tcbuselibrary{raster}
% toute figure TikZ/pgfplots plus large que la ligne est réduite à la largeur disponible
\BeforeBeginEnvironment{tikzpicture}{\begin{adjustbox}{max width=\linewidth}}
\AfterEndEnvironment{tikzpicture}{\end{adjustbox}}
% légendes pgfplots lisibles par-dessus les courbes
\pgfplotsset{every axis/.append style={legend style={fill=white,fill opacity=0.92,text opacity=1,draw=black!15,inner sep=3pt}}}
% étiquettes d'arêtes (syntaxe edge["..."]) avec fond blanc
\tikzset{every edge quotes/.append style={fill=white,inner sep=1.2pt}}
%% FIGFIX-END
\begin{document}

\pagegarde{Leçon 16 — Éléments socioculturels : maladies endémiques en Afrique et en Asie}{Chinois}{Tle A}{Module 2 : Environnement et santé}{ZH16}{2 h}{Cette leçon explore les maladies endémiques en Afrique et en Asie à travers une approche comparative factuelle et respectueuse Cameroun/Chine. Elle outille l'apprenant en vocabulaire médical spécialisé, structures de comparaison et d'argumentation pour produire un mini-exposé oral structuré.
\vspace{4pt}
\begin{multicols}{2}\small
\begin{itemize}
\item Identifier et écrire le vocabulaire clé des maladies endémiques, symptômes et prévention en caractères chinois et pinyin.
\item Analyser un texte authentique simplifié sur la santé publique en Chine et en Afrique.
\item Comparer factuellement les approches camerounaise et chinoise (médecine traditionnelle, vaccination, politiques publiques).
\item Utiliser les structures de comparaison (比, 跟...一样/不一样) et de cause-conséquence (因为...所以..., 由于...因此...).
\item Exprimer un point de vue personnel argumenté sur une question de santé publique en chinois oral.
\item Réaliser un mini-exposé de 2 minutes mobilisant lexique, faits comparés et connecteurs logiques.
\end{itemize}\end{multicols}}

\begin{sommaire}
\begin{multicols}{2}
\begin{enumerate}
\item 1. Entrée en matière \& Lexique thématique : Maladies, Symptômes, Prévention
\item 2. Texte support : « Lutte contre les épidémies : regards croisés »
\item 3. Faits socioculturels comparés : Cameroun vs Chine (Tableau factuel)
\item 4. Outils linguistiques : Comparer, Argumenter, Nuancer
\item 5. Atelier oral : Jeu de rôle \& Débat guidé « Forum Santé Jeunesse »
\item 6. Production écrite intégrée \& Évaluation formative
\item Activité d'intégration
\item Méthodes et exercices résolus
\item Exercices
\item Corrigés des exercices
\item Fiche bilan
\end{enumerate}
\end{multicols}
\end{sommaire}

\begin{objectifs}
\begin{itemize}
\item Identifier et écrire le vocabulaire clé des maladies endémiques, symptômes et prévention en caractères chinois et pinyin.
\item Analyser un texte authentique simplifié sur la santé publique en Chine et en Afrique.
\item Comparer factuellement les approches camerounaise et chinoise (médecine traditionnelle, vaccination, politiques publiques).
\item Utiliser les structures de comparaison (比, 跟...一样/不一样) et de cause-conséquence (因为...所以..., 由于...因此...).
\item Exprimer un point de vue personnel argumenté sur une question de santé publique en chinois oral.
\item Réaliser un mini-exposé de 2 minutes mobilisant lexique, faits comparés et connecteurs logiques.
\item Reconnaître les radicaux et l'ordre des traits de 10 caractères médicaux fréquents.
\end{itemize}
\end{objectifs}

\begin{prerequis}
\begin{itemize}
\item Vocabulaire de base du corps humain (头, 手, 肚子) et sensations (疼, 热, 冷) vu en Première.
\item Structures de base de la comparaison avec 比 et 跟...一样 (Première, Module 1).
\item Expression de la cause et de la conséquence avec 因为...所以... (Première).
\item Notions de base sur la géographie de la Chine et du Cameroun (capitales, climat).
\item Méthode de prise de notes en tableau comparatif (compétence transversale).
\end{itemize}
\end{prerequis}

\newpage

\coursec{Entrée en matière \& Lexique thématique : Maladies, Symptômes, Prévention}

\courssub{Situation d'accroche : deux continents, une même urgence}

Au centre de santé de Mvog-Ada à Yaoundé, un infirmier explique patiemment à une patiente les symptômes du paludisme : fièvre, frissons, fatigue. Sur l'écran de la salle d'attente passe un reportage sur la lutte de la Chine contre les épidémies : dépistages massifs, vaccination, médecine traditionnelle chinoise en appui. Deux continents, des défis sanitaires différents, mais une même urgence : \cle{comprendre}, \cle{prévenir} et \cle{soigner}.

\textbf{Problématique :} comment les savoirs traditionnels et la médecine moderne s'articulent-ils face aux maladies endémiques au Cameroun et en Asie ? Pour y répondre, il faut d'abord maîtriser le vocabulaire chinois de la santé : c'est l'objectif de cette première section.

\begin{definition}[Qu'est-ce qu'une maladie endémique ?]
Une \cle{maladie endémique} est une maladie présente de façon permanente ou régulière dans une région donnée. En chinois, on distingue :
\begin{itemize}
\item \zh{地方病} \emph{dìfāngbìng} : « maladie du lieu », maladie endémique liée à un territoire (climat, eau, insectes) ;
\item \zh{流行病} \emph{liúxíngbìng} : « maladie qui se répand », maladie épidémique qui circule dans une population.
\end{itemize}
Exemple : \zh{疟疾是喀麦隆的地方病。} \emph{Nüèjí shì Kāmàilóng de dìfāngbìng.} « Le paludisme est une maladie endémique au Cameroun. »
\end{definition}

\begin{popfigure}
\begin{tikzpicture}[node distance=1.5cm, every node/.style={align=center}]
\node[draw=popblue, fill=popblueL, rounded corners, thick, font=\bfseries, align=center] (center){Maladies endémiques\\地方病 \emph{dìfāngbìng}};
\node[below left of=center, draw=popgreen, fill=popgreenL, rounded corners, xshift=-1.2cm, align=center] (af){Afrique / Cameroun\\非洲 / 喀麦隆};
\node[below right of=center, draw=poporange, fill=poporangeL, rounded corners, xshift=1.2cm, align=center] (as){Asie / Chine\\亚洲 / 中国};
\node[below of=af, draw=popgreen, rounded corners, align=center] (m1){Paludisme 疟疾\\Choléra 霍乱};
\node[below of=as, draw=poporange, rounded corners, align=center] (m2){Dengue 登革热\\Tuberculose 结核病};
\draw[->, thick, popblue] (center) -- (af);
\draw[->, thick, popblue] (center) -- (as);
\draw[->, popgreen] (af) -- (m1);
\draw[->, poporange] (as) -- (m2);
\end{tikzpicture}
\legende{Carte mentale d'entrée : deux aires géographiques, des maladies endémiques différentes.}
\end{popfigure}

\courssub{Lexique thématique : le glossaire structuré}

Observons d'abord le vocabulaire en contexte, regroupé en trois familles logiques : les \cle{maladies}, les \cle{symptômes} et la \cle{prévention / le traitement}.

\begin{textebox}[词汇表 — Glossaire du thème]
疟疾 \emph{nüèjí} — paludisme\\
霍乱 \emph{huòluàn} — choléra\\
登革热 \emph{dēnggérè} — dengue\\
结核病 \emph{jiéhébìng} — tuberculose\\
艾滋病 \emph{àizībìng} — SIDA\\
疫苗 \emph{yìmiáo} — vaccin\\
蚊帐 \emph{wénzhàng} — moustiquaire\\
中西医结合 \emph{Zhōng-Xīyī jiéhé} — médecine intégrée chino-occidentale
\end{textebox}

\textbf{Famille 1 : les maladies (疾病 \emph{jíbìng})}

\begin{center}
\begin{tabularx}{\linewidth}{|X|X|X|X|}
\hline
\rowcolor{popblueL} Caractères & Pinyin & Nature & Traduction \\
\hline
疟疾 & nüèjí & nom & paludisme \\
\hline
霍乱 & huòluàn & nom & choléra \\
\hline
登革热 & dēnggérè & nom & dengue \\
\hline
结核病 & jiéhébìng & nom & tuberculose \\
\hline
艾滋病 & àizībìng & nom & SIDA \\
\hline
地方病 & dìfāngbìng & nom & maladie endémique \\
\hline
流行病 & liúxíngbìng & nom & maladie épidémique \\
\hline
\end{tabularx}
\end{center}

\textbf{Famille 2 : les symptômes (症状 \emph{zhèngzhuàng})}

\begin{center}
\begin{tabularx}{\linewidth}{|X|X|X|X|}
\hline
\rowcolor{popblueL} Caractères & Pinyin & Nature & Traduction \\
\hline
发烧 & fāshāo & verbe/nom & avoir de la fièvre \\
\hline
发冷 & fālěng & verbe & avoir froid, frissonner \\
\hline
腹泻 & fùxiè & verbe/nom & diarrhée \\
\hline
咳血 & kéxuè & verbe & cracher du sang \\
\hline
乏力 & fálì & adjectif & fatigué, sans force \\
\hline
头疼 & tóuténg & verbe & avoir mal à la tête \\
\hline
\end{tabularx}
\end{center}

\textbf{Famille 3 : prévention et traitement (预防和治疗 \emph{yùfáng hé zhìliáo})}

\begin{center}
\begin{tabularx}{\linewidth}{|X|X|X|X|}
\hline
\rowcolor{popblueL} Caractères & Pinyin & Nature & Traduction \\
\hline
疫苗 & yìmiáo & nom & vaccin \\
\hline
疫苗接种 & yìmiáo jiēzhòng & nom & vaccination \\
\hline
蚊帐 & wénzhàng & nom & moustiquaire \\
\hline
中药 & Zhōngyào & nom & médecine traditionnelle chinoise \\
\hline
西药 & Xīyào & nom & médecine occidentale \\
\hline
隔离 & gélí & verbe/nom & isolement, quarantaine \\
\hline
预防 & yùfáng & verbe & prévenir \\
\hline
治疗 & zhìliáo & verbe & soigner, traiter \\
\hline
\end{tabularx}
\end{center}

\courssub{Écriture des caractères : les radicaux de la santé}

En chinois, les caractères ne sont pas des dessins arbitraires : ils sont construits autour de \cle{radicaux} (clés) qui indiquent le sens. Le domaine de la santé possède des radicaux très réguliers : les reconnaître, c'est deviner le sens d'un caractère inconnu.

\begin{definition}[Radicaux clés du domaine médical]
\begin{itemize}
\item \zh{疒} \emph{nè} : radical de la \cle{maladie}. Il enveloppe le caractère par le haut et la gauche. Exemples : 病 \emph{bìng} (maladie), 疾 \emph{jí} (maladie), 疫 \emph{yì} (épidémie), 症 \emph{zhèng} (symptôme), 疟 \emph{nüè} (paludisme).
\item \zh{月} (variante de 肉 \emph{ròu}, « chair ») : radical des \cle{organes et parties du corps}. Exemples : 肝 \emph{gān} (foie), 肺 \emph{fèi} (poumons), 肠 \emph{cháng} (intestins).
\item \zh{氵} (eau) : radical des \cle{liquides}. Exemples : 泻 \emph{xiè} (diarrhée), 液 \emph{yè} (liquide).
\item \zh{虫} \emph{chóng} : radical des \cle{insectes, vers et parasites}. Exemples : 蚊 \emph{wén} (moustique), 蛇 \emph{shé} (serpent), 寄生虫 \emph{jìshēngchóng} (parasite).
\end{itemize}
\end{definition}

\begin{propriete}[Ordre des traits : 疫 et 苗]
Règle générale d'écriture : de \textbf{haut en bas}, de \textbf{gauche à droite}, l'\textbf{horizontal avant le vertical}, l'\textbf{extérieur avant l'intérieur}.

\textbf{疫} \emph{yì} (épidémie), 9 traits : on écrit d'abord le radical 疒 (extérieur, 5 traits) : 丶 (point), 一 (horizontal), 丿 (trait gauche), 丶 (point), ㇀ (remontant) ; puis l'intérieur 殳 (4 traits) : 丿, ㇍, ㇇, ㇏. L'extérieur (疒) précède toujours l'intérieur.

\textbf{苗} \emph{miáo} (pousse, jeune plant), 8 traits : d'abord la clé de l'herbe 艹 en haut (3 traits : 一, 丨, 丨), puis 田 « champ » en bas (5 traits : 丨, (trait brisé), 一, 丨, 一). Image mnémotechnique : une \emph{pousse} (艹) qui sort du \emph{champ} (田) — d'où 疫苗, le « vaccin », littéralement « pousse contre l'épidémie ».

Autres caractères cibles de la leçon : 疾 (10 traits, radical 疒 + 矢 « flèche »), 预 \emph{yù} (10 traits, « pré- »), 蚊 \emph{wén} (10 traits, radical 虫 + 文), 帐 \emph{zhàng} (7 traits, radical 巾 « tissu » + 长).
\end{propriete}

\begin{methode}[Mémoriser les caractères de la santé]
\begin{enumerate}
\item Associe le radical 疒 à l'image d'un \textbf{lit d'hôpital} vu de profil : le trait horizontal est le matelas, le trait gauche le pied du lit. Tout caractère « coiffé » de 疒 parle de maladie.
\item Crée des \textbf{cartes mentales par radical} : un centre (疒), des branches (病, 疾, 疫, 症, 疟) avec pinyin et traduction.
\item Apprends les caractères par \textbf{paires de sens} : 疫苗 (vaccin) et 蚊帐 (moustiquaire) sont des mots de prévention ; 发烧 et 腹泻 sont des mots de symptômes.
\item Réécris chaque caractère cible \textbf{cinq fois} en comptant les traits à voix haute.
\end{enumerate}
\end{methode}

\begin{attention}[Pièges fréquents des francophones]
\begin{itemize}
\item Ne confonds pas \zh{病} et \zh{疾}. 病 \emph{bìng} est le nom courant : 我病了 « je suis malade ». 疾 \emph{jí} est plus formel ou littéraire (疾病 « maladies » en terme soutenu ; 疾走 « marcher vite », sens ancien de « rapide »). Dans la conversation, utilise 病.
\item N'écris jamais \zh{气} (air, souffle) à la place du radical \zh{疒} : 病 ne s'écrit pas avec 气.
\item Attention aux tons : \zh{疟疾} \emph{nüèjí} (4\textsuperscript{e} + 2\textsuperscript{e} ton) ; un ton faux peut faire perdre le sens.
\item 登革热 \emph{dēnggérè} est un emprunt phonétique à l'anglais \emph{dengue} : ne cherche pas à traduire chaque caractère séparément.
\item Ne confonds pas \zh{预防} \emph{yùfáng} (prévenir) et \zh{治疗} \emph{zhìliáo} (soigner) : la prévention précède la maladie, le traitement la suit.
\end{itemize}
\end{attention}

\courssub{Collocations utiles : parler de la santé en phrases}

En chinois comme en français, les mots voyagent par groupes stables. Voici les \cle{collocations} essentielles du thème :

\begin{center}
\begin{tabularx}{\linewidth}{|X|X|X|}
\hline
\rowcolor{popblueL} Collocation & Pinyin & Traduction \\
\hline
患上疟疾 & huànshàng nüèjí & contracter le paludisme \\
\hline
接种疫苗 & jiēzhòng yìmiáo & se faire vacciner \\
\hline
传播途径 & chuánbō tújìng & voie de transmission \\
\hline
防疫措施 & fángyì cuòshī & mesures de prévention sanitaire \\
\hline
中西医结合 & Zhōng-Xīyī jiéhé & combiner médecine chinoise et occidentale \\
\hline
\end{tabularx}
\end{center}

\begin{exemplebox}[Le vocabulaire en action]
\zh{我患上了疟疾，发烧又发冷。} \emph{Wǒ huànshàng le nüèjí, fāshāo yòu fālěng.} « J'ai attrapé le paludisme, j'ai de la fièvre et des frissons. »\\[4pt]
\zh{我们要接种疫苗，预防疾病。} \emph{Wǒmen yào jiēzhòng yìmiáo, yùfáng jíbìng.} « Nous devons nous faire vacciner pour prévenir les maladies. »\\[4pt]
\zh{蚊子叮咬是疟疾的传播途径。} \emph{Wénzi dīngyǎo shì nüèjí de chuánbō tújìng.} « La piqûre de moustique est la voie de transmission du paludisme. »\\[4pt]
\zh{中西医结合治疗效果好。} \emph{Zhōng-Xīyī jiéhé zhìliáo xiàoguǒ hǎo.} « Le traitement combiné médecine chinoise et occidentale est efficace. »\\[4pt]
\zh{睡觉时挂蚊帐可以预防疟疾。} \emph{Shuìjiào shí guà wénzhàng kěyǐ yùfáng nüèjí.} « Suspendre une moustiquaire au moment de dormir permet de prévenir le paludisme. »
\end{exemplebox}

Remarque grammaticale : la structure \zh{患上 + maladie} correspond au français « contracter / attraper une maladie » ; \zh{接种 + vaccin} signifie « recevoir une injection » (接 = recevoir, 种 = inoculer). Aucune préposition n'est nécessaire entre le verbe et son objet, contrairement au français « se vacciner \emph{contre} » : on dit directement \zh{接种疫苗} et non « 接种对疫苗 ».

\begin{exoresolu}[Reconnaissance visuelle : caractère → pinyin → sens]
Énoncé : pour chaque caractère ou mot, donne le pinyin avec tons et la traduction française : ① 疟疾 ② 疫苗 ③ 发烧 ④ 蚊帐 ⑤ 隔离.
\tcblower
Solution détaillée :
\begin{enumerate}
\item 疟疾 : \emph{nüèjí} — le paludisme (radical 疒 dans les deux caractères : indice « maladie »).
\item 疫苗 : \emph{yìmiáo} — le vaccin (疫 = épidémie, 苗 = pousse).
\item 发烧 : \emph{fāshāo} — avoir de la fièvre (烧 = brûler : la fièvre « brûle »).
\item 蚊帐 : \emph{wénzhàng} — la moustiquaire (蚊 = moustique, radical 虫 ; 帐 = tissu, radical 巾).
\item 隔离 : \emph{gélí} — l'isolement, la quarantaine.
\end{enumerate}
Méthode : repère d'abord le radical pour deviner la famille de sens, puis vérifie le pinyin.
\end{exoresolu}

\begin{exercice}[Dictée courte : pinyin → caractères]
Écris en caractères chinois les mots dictés suivants, puis traduis-les en français :
\begin{enumerate}
\item \emph{huòluàn}
\item \emph{yùfáng}
\item \emph{fùxiè}
\item \emph{Zhōngyào}
\item \emph{jiēzhòng yìmiáo}
\end{enumerate}
\end{exercice}

\begin{corrige}[Dictée courte]
\begin{enumerate}
\item 霍乱 — le choléra
\item 预防 — prévenir
\item 腹泻 — la diarrhée
\item 中药 — la médecine traditionnelle chinoise
\item 接种疫苗 — se faire vacciner
\end{enumerate}
Vérifie l'ordre des traits : 预 s'écrit de gauche à droite (予 puis 页) ; 霍 s'écrit de haut en bas (雨 puis 隹).
\end{corrige}

\begin{aretenir}[L'essentiel de la section]
\begin{itemize}
\item Une maladie endémique se dit \zh{地方病} \emph{dìfāngbìng} ; une épidémie se dit \zh{流行病} \emph{liúxíngbìng}.
\item Le radical \zh{疒} signale la maladie ; \zh{月/肉} les organes ; \zh{虫} les insectes et parasites ; \zh{氵} les liquides.
\item Collocations clés : 患上疟疾 (contracter le paludisme), 接种疫苗 (se faire vacciner), 传播途径 (voie de transmission), 防疫措施 (mesures de prévention), 中西医结合 (médecine intégrée).
\item Écriture : extérieur avant intérieur (疫, 9 traits), haut avant bas (苗, 8 traits).
\end{itemize}
\end{aretenir}

\coursec{Texte support : « Lutte contre les épidémies : regards croisés »}

Dans la section précédente, nous avons construit le lexique thématique des maladies, des symptômes et de la prévention (\zh{疟疾} \emph{nüèji} « paludisme », \zh{症状} \emph{zhèngzhuàng} « symptôme », \zh{预防} \emph{yùfáng} « prévention »). Nous allons maintenant mettre ce vocabulaire en action dans un texte authentique de type article de presse. L'objectif est double : comprendre un document écrit en chinois et dégager des faits socioculturels comparés entre le Cameroun et la Chine.

\courssub{Stratégie de lecture active}

Avant de lire le texte, adoptons une démarche méthodique. Un bon lecteur ne lit pas un texte chinois mot à mot : il le parcourt d'abord globalement, puis en détail.

\begin{methode}[Stratégie de lecture en quatre étapes]
\begin{enumerate}
\item \textbf{Lire le titre} et anticiper le contenu : de quoi va parler le texte ?
\item \textbf{Surligner les noms propres et les chiffres} : pays, organisations, dates. Ce sont des repères faciles à repérer dans les caractères (\zh{中国}, \zh{2021年}, \zh{世卫组织}).
\item \textbf{Identifier les connecteurs logiques} : \zh{但是} \emph{dànshì} « mais », \zh{因此} \emph{yīncǐ} « c'est pourquoi », \zh{如今} \emph{rújīn} « aujourd'hui ». Ils structurent l'argumentation.
\item \textbf{Résumer chaque paragraphe en une phrase en français}, sans traduire mot à mot.
\end{enumerate}
\end{methode}

\courssub{Lecture du texte}

\begin{textebox}[Texte : 非洲和亚洲的地方病]
文本：\\[4pt]
非洲和亚洲面临着不同的地方病挑战。喀麦隆是疟疾高发国家，疟疾每年导致许多儿童死亡。中国通过大规模使用青蒿素和蚊帐，于2021年获得世卫组织“无疟疾”认证。青蒿素源于中医典籍《肘后备急方》，是中西医结合的典范。如今，由于气候变化和抗药性的出现，两国不但在疫苗研发上努力，而且在疾控合作上加强交流。\\[8pt]
\emph{Fēizhōu hé Yàzhōu miànlínzhe bùtóng de dìfāngbìng tiǎozhàn. Kāmàilóng shì nüèji gāofā guójiā, nüèji měi nián dǎozhì xǔduō értóng sǐwáng. Zhōngguó tōngguò dà guīmó shǐyòng qīnghāosù hé wénzhàng, yú 2021 nián huòdé Shìwèi Zǔzhī “wú nüèji” rènzhèng. Qīnghāosù yuányú Zhōngyī diǎnjí 《Zhǒu Hòu Bèijí Fāng》, shì Zhōng-Xīyī jiéhé de diǎnfàn. Rújīn, yóuyú qìhòu biànhuà hé kàngyàoxìng de chūxiàn, liǎng guó bùdàn zài yìmiáo yánfā shang nǔlì, érqiě zài jíkòng hézuò shang jiāqiáng jiāoliú.}\\[8pt]
Traduction : L'Afrique et l'Asie font face à des défis différents en matière de maladies endémiques. Le Cameroun est un pays à forte incidence du paludisme ; le paludisme provoque chaque année la mort de nombreux enfants. La Chine, grâce à l'utilisation à grande échelle de l'artémisinine et des moustiquaires, a obtenu en 2021 la certification « sans paludisme » de l'OMS. L'artémisinine, issue du traité classique de médecine traditionnelle chinoise \emph{Recettes d'urgence derrière le coude}, est un modèle d'alliance entre médecine chinoise et occidentale. Aujourd'hui, en raison du changement climatique et de l'apparition de résistances, les deux pays non seulement déploient des efforts dans la recherche sur les vaccins, mais renforcent aussi leurs échanges en matière de contrôle des maladies.
\end{textebox}

\courssub{Glossaire du texte}

\begin{center}
\begin{tabularx}{\linewidth}{|X|X|X|X|}
\hline
\rowcolor{popblueL} Caractères & Pinyin & Nature & Traduction \\
\hline
地方病 & dìfāngbìng & nom & maladie endémique \\
\hline
面临 & miànlín & verbe & faire face à \\
\hline
挑战 & tiǎozhàn & nom & défi \\
\hline
疟疾 & nüèji & nom & paludisme \\
\hline
高发 & gāofā & adj./verbe & à forte incidence \\
\hline
导致 & dǎozhì & verbe & entraîner, provoquer \\
\hline
大规模 & dà guīmó & loc. adj. & à grande échelle \\
\hline
青蒿素 & qīnghāosù & nom & artémisinine \\
\hline
蚊帐 & wénzhàng & nom & moustiquaire \\
\hline
获得 & huòdé & verbe & obtenir \\
\hline
认证 & rènzhèng & nom/verbe & certification ; certifier \\
\hline
典籍 & diǎnjí & nom & ouvrage classique, canon \\
\hline
源于 & yuányú & verbe & provenir de, tirer son origine de \\
\hline
典范 & diǎnfàn & nom & modèle, exemple parfait \\
\hline
抗药性 & kàngyàoxìng & nom & résistance (aux médicaments) \\
\hline
疫苗 & yìmiáo & nom & vaccin \\
\hline
研发 & yánfā & verbe & rechercher et développer \\
\hline
疾控 & jíkòng & nom & contrôle des maladies \\
\hline
加强 & jiāqiáng & verbe & renforcer \\
\hline
交流 & jiāoliú & verbe/nom & échanger ; échange \\
\hline
\end{tabularx}
\end{center}

\begin{definition}[Mots-clés du texte]
\cle{青蒿素} \emph{qīnghāosù} : artémisinine, molécule antipaludique découverte par Tu Youyou (\zh{屠呦呦} \emph{Tú Yōuyōu}), Prix Nobel de médecine 2015.\\
\cle{世卫组织} \emph{Shìwèi Zǔzhī} : Organisation mondiale de la Santé (OMS), abréviation de \zh{世界卫生组织} \emph{Shìjiè Wèishēng Zǔzhī}.\\
\cle{认证} \emph{rènzhèng} : certification officielle, reconnaissance formelle par une autorité.\\
\cle{高发} \emph{gāofā} : forte incidence, forte prévalence (d'une maladie dans une zone).\\
\cle{面临} \emph{miànlín} : « faire face à », suivi directement d'un nom : \zh{面临挑战} « faire face à un défi ». Ne pas ajouter de préposition après ce verbe.
\end{definition}

\courssub{Frise chronologique : deux parcours de lutte contre le paludisme}

\begin{popfigure}
\begin{tikzpicture}
\draw[thick, ->, popdark] (0,0) -- (10.5,0) node[right] {\small Temps};
\foreach \x/\annee in {0/2000, 2/2005, 4/2010, 6/2015, 8/2020, 10/2024}
  \draw[popdark] (\x,0.1) -- (\x,-0.1) node[below, font=\small] {\annee};
\node[above, align=center, text=popblue, font=\small] at (2,0.4) {Chine : recherche sur\\ \zh{青蒿素} (années 1970-2000)};
\node[above, align=center, text=popblue, font=\small] at (6,0.4) {Chine : zéro cas\\ autochtone (2017)};
\node[above, align=center, text=popblue, font=\small] at (8.8,1.2) {Chine : certification\\ OMS (2021)};
\node[below, align=center, text=poporange, font=\small] at (3.4,-0.6) {Cameroun : distribution\\ de moustiquaires (MIIL)};
\node[below, align=center, text=poporange, font=\small] at (8,-0.6) {Cameroun : campagne\\ vaccin RTS,S (2024)};
\end{tikzpicture}
\legende{Frise comparée : étapes majeures de la lutte antipaludique en Chine et au Cameroun.}
\end{popfigure}

\courssub{Questions de compréhension}

\textbf{A. Compréhension globale}
\begin{enumerate}
\item Quelle est l'idée principale du texte ? Résumez-la en une phrase en français.
\item Quelles sont les deux traditions médicales mentionnées dans le texte ?
\end{enumerate}

\textbf{B. Compréhension détaillée}
\begin{enumerate}
\item En quelle année la Chine a-t-elle obtenu la certification « sans paludisme » ?
\item Quelles sont les deux mesures utilisées massivement par la Chine contre le paludisme ?
\item Quels sont les deux défis actuels cités à la fin du texte ?
\end{enumerate}

\textbf{C. Compréhension inférentielle}
\begin{enumerate}
\item Pourquoi, selon vous, la Chine a-t-elle réussi à éliminer le paludisme ? Relevez dans le texte les éléments qui permettent de répondre.
\item Pourquoi le texte dit-il que l'artémisinine est un « modèle d'alliance entre médecine chinoise et occidentale » ?
\end{enumerate}

\courssub{Recherche des structures cibles dans le texte}

Le texte contient trois structures grammaticales essentielles du programme de Terminale. Cherchons-les, observons-les, puis énonçons la règle.

\textbf{Observation.} Relevez dans le texte les phrases contenant \zh{通过}, \zh{由于} et \zh{不但……而且……}. Que remarquez-vous sur leur place dans la phrase ?

\begin{propriete}[La structure 不但……而且…… « non seulement… mais encore… »]
\textbf{Forme :} Sujet + \zh{不但} + Verbe/Adjectif 1 + \zh{而且} + Verbe/Adjectif 2.\\
\textbf{Emploi :} elle exprime une \emph{addition graduelle} : le second élément ajoute une information plus forte ou complémentaire au premier. C'est l'équivalent de « non seulement… mais aussi/encore… ».\\
\textbf{Exemple du texte :} \zh{两国不但在疫苗研发上努力，而且在疾控合作上加强交流。} \emph{Liǎng guó bùdàn zài yìmiáo yánfā shang nǔlì, érqiě zài jíkòng hézuò shang jiāqiáng jiāoliú.} « Les deux pays non seulement déploient des efforts dans la recherche vaccinale, mais renforcent aussi leurs échanges en matière de contrôle des maladies. »\\
\textbf{Exception :} si les deux membres ont des sujets différents, \zh{不但} se place \emph{avant} le premier sujet : \zh{不但中国成功了，而且喀麦隆也在进步。} \emph{Bùdàn Zhōngguó chénggōng le, érqiě Kāmàilóng yě zài jìnbù.} « Non seulement la Chine a réussi, mais le Cameroun progresse aussi. »
\end{propriete}

\begin{exemplebox}[Autres exemples de la structure]
\zh{中国不但消除了疟疾，而且帮助非洲国家防疫。} \emph{Zhōngguó bùdàn xiāochúle nüèji, érqiě bāngzhù Fēizhōu guójiā fángyì.} « La Chine a non seulement éliminé le paludisme, mais elle aide aussi les pays africains à prévenir les épidémies. »\\[4pt]
\zh{青蒿素不但有效，而且价格不高。} \emph{Qīnghāosù bùdàn yǒuxiào, érqiě jiàgé bù gāo.} « L'artémisinine est non seulement efficace, mais en plus peu coûteuse. »\\[4pt]
\zh{这种蚊帐不但便宜，而且很耐用。} \emph{Zhè zhǒng wénzhàng bùdàn piányi, érqiě hěn nàiyòng.} « Ce type de moustiquaire est non seulement bon marché, mais aussi très résistant. »
\end{exemplebox}

\begin{propriete}[La structure 由于……因此…… « en raison de… c'est pourquoi… »]
\textbf{Forme :} \zh{由于} + cause, (\zh{因此}) + conséquence.\\
\textbf{Emploi :} structure de cause à conséquence, plus soutenue que \zh{因为……所以……}. On peut omettre \zh{因此}, mais jamais inverser l'ordre cause/conséquence.\\
\textbf{Exemple :} \zh{由于气候变化，疟疾的传播地区扩大了。} \emph{Yóuyú qìhòu biànhuà, nüèji de chuánbō dìqū kuòdà le.} « En raison du changement climatique, les zones de transmission du paludisme se sont étendues. »\\
\textbf{Variante :} \zh{因此} peut aussi se placer après le sujet de la seconde proposition : \zh{由于气候变化，疟疾的传播地区因此扩大了。}
\end{propriete}

\begin{propriete}[Le comparatif avec 比]
\textbf{Forme :} A + \zh{比} + B + Adjectif.\\
\textbf{Exemple :} \zh{喀麦隆的疟疾发病率比中国高。} \emph{Kāmàilóng de nüèji fābìnglǜ bǐ Zhōngguó gāo.} « Le taux d'incidence du paludisme au Cameroun est plus élevé qu'en Chine. »\\
\textbf{Rappel :} dans une comparaison avec \zh{比}, l'adjectif n'est jamais précédé de \zh{很}. Pour nuancer ou renforcer, on utilise \zh{比……高得多} « beaucoup plus élevé que », \zh{比……高一点} « un peu plus élevé que », ou \zh{比……更高} « encore plus élevé que » (\zh{更} est correct, \zh{很} ne l'est pas).
\end{propriete}

\begin{center}
\begin{tabularx}{\linewidth}{|X|X|X|}
\hline
\rowcolor{popblueL} Structure & Exemple (caractères + pinyin) & Traduction \\
\hline
A + 比 + B + Adj. & \zh{疟疾比流感更危险。} \emph{Nüèji bǐ liúgǎn gèng wēixiǎn.} & Le paludisme est plus dangereux que la grippe. \\
\hline
由于 + cause, 因此 + conséq. & \zh{由于抗药性，治疗更困难了。} \emph{Yóuyú kàngyàoxìng, zhìliáo gèng kùnnan le.} & En raison des résistances, le traitement est devenu plus difficile. \\
\hline
S + 不但 + V1, 而且 + V2 & \zh{他不但研究疫苗，而且培训医生。} \emph{Tā bùdàn yánjiū yìmiáo, érqiě péixùn yīshēng.} & Il ne fait pas seulement de la recherche vaccinale, il forme aussi des médecins. \\
\hline
\end{tabularx}
\end{center}

\courssub{Travail prosodique : lecture à voix haute}

La lecture orale du texte exige une attention particulière aux tons et aux liaisons.

\begin{itemize}
\item \textbf{Troisième ton :} dans \zh{疟疾} \emph{nüèji}, le quatrième ton suit le troisième : le 3\textsuperscript{e} ton est alors prononcé « demi-ton » (bas, sans remontée). Dans \zh{所以} \emph{suǒyǐ}, deux 3\textsuperscript{e} tons se suivent : le premier devient un 2\textsuperscript{e} ton (\emph{suóyǐ}).
\item \textbf{Ton neutre :} dans \zh{孩子们} \emph{háizimen}, le suffixe \zh{们} est atone ; dans \zh{得了} \emph{déle}, la particule \zh{了} est toujours au ton neutre.
\item \textbf{Liaisons et pauses :} marquez une courte pause après chaque virgule chinoise \zh{，} et une pause longue après le point \zh{。}. Ne vous arrêtez jamais au milieu d'un mot de deux caractères comme \zh{认证} ou \zh{疫苗}.
\end{itemize}

\begin{experience}[Lecture en écho]
L'enseignant lit le texte phrase par phrase ; les élèves répètent en chœur en imitant exactement les tons. Puis deux élèves lisent à tour de rôle : l'un lit la phrase en chinois, l'autre donne immédiatement la traduction française. Le reste de la classe vérifie avec le texte.
\end{experience}

\courssub{Exercice résolu et pièges}

\begin{exoresolu}[Traduction thème]
Traduisez en chinois : « La Chine a non seulement éliminé le paludisme, mais elle a aussi obtenu la certification de l'OMS. »
\tcblower
\textbf{Solution détaillée.}\\
1. Identifier la structure : « non seulement… mais aussi… » → \zh{不但……而且……}.\\
2. Même sujet pour les deux membres (\zh{中国}) : le sujet se place une seule fois, avant \zh{不但}.\\
3. « Éliminer le paludisme » → \zh{消除了疟疾} ; « obtenir la certification de l'OMS » → \zh{获得世卫组织的认证}.\\
4. Phrase finale : \zh{中国不但消除了疟疾，而且获得了世卫组织的认证。} \emph{Zhōngguó bùdàn xiāochúle nüèji, érqiě huòdéle Shìwèi Zǔzhī de rènzhèng.}
\end{exoresolu}

\begin{attention}[Piège : 获得认证 ou 接受认证 ?]
Les francophones confondent souvent \zh{获得} \emph{huòdé} « obtenir » (action active, le sujet conquiert quelque chose) et \zh{接受} \emph{jiēshòu} « accepter, recevoir » (attitude passive). Dans notre texte, le sujet est la Chine : c'est elle qui a \emph{obtenu} la certification après des années d'efforts. On dit donc \zh{获得认证} et non \zh{接受认证}. Autre erreur fréquente : placer \zh{不但} après le sujet quand les deux sujets sont différents — quand les sujets sont identiques, on dit \zh{他不但……而且……} ; quand ils sont différents, on dit \zh{不但他……而且……}, jamais l'inverse.
\end{attention}

\begin{savaistu}[Une recette vieille de 1 700 ans]
L'artémisinine (\zh{青蒿素}) a été extraite de l'armoise annuelle (\zh{黄花蒿} \emph{huánghuāhāo}) grâce à une recette du traité \zh{《肘后备急方》}, datant de la dynastie des Jin (vers 340 apr. J.-C.) : « Une poignée d'armoise, trempée dans deux litres d'eau ; presser le jus et tout boire. » C'est en remarquant que la recette ne faisait \emph{pas} bouillir la plante que Tu Youyou a compris que la chaleur détruisait la molécule active — l'intuition qui a sauvé des millions de vies, notamment en Afrique.
\end{savaistu}

\begin{exercice}[À vous de jouer]
\begin{enumerate}
\item Complétez avec \zh{不但}, \zh{而且}, \zh{由于} ou \zh{因此} : \zh{（　）气候变化，疟疾的传播地区（　）扩大了。}
\item Traduisez : « Le Cameroun est un pays à forte incidence du paludisme. »
\item Relevez dans le texte trois mots contenant le radical \zh{疒} (clé de la maladie) et expliquez leur sens.
\end{enumerate}
\end{exercice}

\begin{corrige}[Exercice « À vous de jouer »]
1. \zh{由于气候变化，疟疾的传播地区因此扩大了。} \emph{Yóuyú qìhòu biànhuà, nüèji de chuánbō dìqū yīncǐ kuòdà le.} (On peut aussi omettre \zh{因此}.)\\
2. \zh{喀麦隆是疟疾高发国家。} \emph{Kāmàilóng shì nüèji gāofā guójiā.}\\
3. \zh{疟} (dans \zh{疟疾}, paludisme), \zh{病} (dans \zh{地方病}, maladie), \zh{疫} (dans \zh{疫苗}, vaccin). La clé \zh{疒} \emph{nè}, « lit de malade », indique le champ sémantique de la maladie.
\end{corrige}

\begin{aretenir}[L'essentiel de la section]
Le texte « Regards croisés » oppose deux situations : le Cameroun, pays à \cle{forte incidence} (\zh{高发}) du paludisme, et la Chine, certifiée « sans paludisme » par l'OMS en 2021 grâce à l'artémisinine (\zh{青蒿素}) et aux moustiquaires (\zh{蚊帐}). Trois structures à maîtriser : le comparatif \zh{比} (sans \zh{很} devant l'adjectif), la causalité soutenue \zh{由于……因此……} et l'addition graduelle \zh{不但……而且……} (attention à la place de \zh{不但} selon que les sujets sont identiques ou différents). En lecture orale : demi-3\textsuperscript{e} ton devant un autre ton, particules au ton neutre, pauses sur la ponctuation chinoise.
\end{aretenir}

\coursec{Faits socioculturels comparés : Cameroun vs Chine (tableau factuel)}

\courssub{Transition et objectifs de la section}

Dans la section précédente, le texte « Lutte contre les épidémies : regards croisés » nous a permis de découvrir, à travers un dialogue entre deux élèves, comment le Cameroun et la Chine affrontent les maladies endémiques. Nous allons maintenant \cle{organiser ces informations} de manière systématique : d'abord observer les phrases-clés du texte, puis construire un \cle{tableau factuel comparatif}, enfin analyser les \cle{similarités} et les \cle{différences} entre les deux pays. L'objectif est double : enrichir votre culture générale pour le baccalauréat et apprendre à \emph{comparer sans juger}, en chinois comme en français.

\begin{textebox}[Phrases-clés à observer — 关键句]
喀麦隆正在推广疟疾疫苗。\\
\emph{Kāmàilóng zhèngzài tuīguǎng nüèji yìmiáo.}\\
« Le Cameroun est en train de promouvoir le vaccin contre le paludisme. »\\[4pt]
中国已经消除了疟疾。\\
\emph{Zhōngguó yǐjing xiāochú le nüèji.}\\
« La Chine a déjà éliminé le paludisme. »\\[4pt]
两国都重视社区防疫。\\
\emph{Liǎng guó dōu zhòngshì shèqū fángyì.}\\
« Les deux pays attachent tous de l'importance à la prévention communautaire. »\\[4pt]
中医在非洲越来越受欢迎。\\
\emph{Zhōngyī zài Fēizhōu yuè lái yuè shòu huānyíng.}\\
« La médecine chinoise est de plus en plus appréciée en Afrique. »
\end{textebox}

\textbf{Observation guidée.} Avant de remplir le tableau, observons ces quatre phrases :
\begin{itemize}
\item La phrase 1 utilise \zh{正在} (zhèngzài) : action \emph{en cours} → le Cameroun est dans une dynamique de progrès.
\item La phrase 2 utilise \zh{已经……了} (yǐjing... le) : action \emph{accomplie} → la Chine a mené ce combat à son terme.
\item La phrase 3 utilise \zh{都} (dōu, « tous les deux »), placé \emph{après} le sujet \zh{两国} et \emph{avant} le verbe : c'est le marqueur grammatical de la \cle{similitude}.
\item La phrase 4 utilise \zh{越来越} (yuè lái yuè, « de plus en plus ») devant un verbe ou un adjectif : c'est le marqueur d'une \cle{évolution}.
\end{itemize}
Ces quatre petits mots-outils suffisent à construire un discours comparatif nuancé. Retenez-les : ils seront vos outils pour le mini-exposé de la section 5.

\courssub{Vocabulaire socioculturel essentiel}

\begin{definition}[Institutions et concepts de santé]
\begin{itemize}
\item \zh{社区卫生服务中心} (shèqū wèishēng fúwù zhōngxīn) : \cle{centre de service de santé communautaire} (Chine). C'est la porte d'entrée du système de santé chinois en ville : consultation de base, vaccination, suivi des maladies chroniques. Équivalent fonctionnel du « centre de santé intégré » camerounais.
\item \zh{合作医疗} (hézuò yīliáo) : \cle{assurance maladie coopérative}. Ancien système rural chinois (années 1950-1970), ancêtre de l'assurance maladie rurale actuelle. À ne pas confondre avec l'assurance commerciale privée.
\item \zh{传统医学} (chuántǒng yīxué) : \cle{médecine traditionnelle}, terme générique utilisé par l'OMS. En Chine on dit précisément \zh{中医} (Zhōngyī, « médecine chinoise ») ; au Cameroun on parle de « médecine traditionnelle africaine » ou de « pharmacopée ».
\end{itemize}
\end{definition}

\begin{center}
\begin{tabularx}{\linewidth}{|X|X|X|X|}
\hline
\rowcolor{popblueL} Caractères & Pinyin & Nature & Traduction \\
\hline
疟疾 & nüèji & n. & paludisme \\
\hline
疫苗 & yìmiáo & n. & vaccin \\
\hline
防疫站 & fángyì zhàn & n. & station anti-épidémique (Chine) \\
\hline
医疗队 & yīliáo duì & n. & équipe médicale \\
\hline
蚊帐 & wénzhàng & n. & moustiquaire \\
\hline
消除 & xiāochú & v. & éliminer, éradiquer \\
\hline
推广 & tuīguǎng & v. & promouvoir, diffuser \\
\hline
重视 & zhòngshì & v. & attacher de l'importance à \\
\hline
医疗保险 & yīliáo bǎoxiǎn & n. & assurance maladie \\
\hline
相似之处 & xiāngsì zhī chù & loc. n. & points de ressemblance \\
\hline
流行 & liúxíng & v./adj. & être endémique, se répandre \\
\hline
资源 & zīyuán & n. & ressources \\
\hline
\end{tabularx}
\end{center}

\courssub{Le tableau factuel comparatif}

Voici le tableau de référence de la leçon. Il synthétise les faits, sans jugement de valeur, à partir de sources publiques (OMS, MINSANTE, Commission nationale de la santé de Chine).

\begin{center}
\begin{tabularx}{\linewidth}{|X|X|X|}
\hline
\rowcolor{popblueL} Critère & Cameroun & Chine \\
\hline
Maladies cibles & Paludisme endémique (première cause de consultation), choléra, méningite, fièvre jaune & Paludisme éliminé (certification OMS 2021) ; surveillance des hépatites, de la tuberculose et des grippes \\
\hline
Médecine traditionnelle & Pharmacopée reconnue par les textes en vigueur ; pratique \emph{parallèle} au système public (tradipraticiens) & 中医 \emph{intégrée} au système : hôpitaux de médecine chinoise (中医院), enseignement universitaire, remboursement partiel \\
\hline
Couverture vaccinale & Programme Élargi de Vaccination (PEV) ; introduction du vaccin antipaludique RTS,S en 2024 & Calendrier national gratuit pour les enfants ; couverture très élevée grâce au réseau des 防疫站 \\
\hline
Infrastructure de base & Centres de santé intégrés ; accès difficile dans certaines zones rurales et enclavées & 社区卫生服务中心 en ville, cliniques de village en zone rurale ; maillage dense \\
\hline
Rôle de l'État & PNLP (Programme National de Lutte contre le Paludisme) : distribution gratuite de moustiquaires imprégnées ; système décentralisé (régions, districts) & Politique centralisée pilotée par la Commission nationale de la santé ; assurance maladie quasi universelle ; gestion stricte du COVID-19 puis réouverture \\
\hline
\end{tabularx}
\end{center}

\begin{aretenir}[L'essentiel en trois faits]
\begin{enumerate}
\item \textbf{Cameroun} : le paludisme reste endémique, mais l'État agit (PNLP, moustiquaires gratuites, vaccin RTS,S depuis 2024) et la médecine traditionnelle est reconnue par les textes en vigueur.
\item \textbf{Chine} : le paludisme a été éliminé après environ 70 ans d'efforts (de plusieurs dizaines de millions de cas par an dans les années 1940 à zéro cas local) ; la médecine chinoise est \emph{intégrée} aux hôpitaux.
\item \textbf{Point commun} : les deux pays misent sur la \cle{prévention communautaire} et la \cle{vaccination}.
\end{enumerate}
\end{aretenir}

\courssub{Schéma de synthèse : convergences et divergences}

\begin{popfigure}
\begin{tikzpicture}[every node/.style={font=\small}]
\node[draw, rectangle, rounded corners, fill=popblueL, align=center] (cam) at (-3.5,0) {\textbf{Cameroun}};
\node[draw, rectangle, rounded corners, fill=poppinkL, align=center] (chi) at (3.5,0) {\textbf{Chine}};
\node[draw, rectangle, rounded corners, fill=popgreenL, align=center] (com) at (0,-2.6) {\textbf{Points communs}};
\node[align=center, fill=popblue!10, rounded corners, inner sep=4pt] at (-3.5,1.6) {Paludisme endémique\\Médecine trad. parallèle\\Système décentralisé};
\node[align=center, fill=poppink!30, rounded corners, inner sep=4pt] at (3.5,1.6) {Paludisme éliminé (2021)\\Médecine trad. intégrée\\Système centralisé};
\node[align=center, fill=popgreen!15, rounded corners, inner sep=4pt] at (0,-4.1) {Coopération Sud-Sud\\Vaccination prioritaire\\Prévention communautaire};
\draw[->, thick, popblue] (cam) -- (com);
\draw[->, thick, poppink] (chi) -- (com);
\end{tikzpicture}
\legende{Carte de comparaison : spécificités du Cameroun et de la Chine, et zone de convergence.}
\end{popfigure}

\courssub{Analyse des similarités}

Trois convergences majeures relient les deux pays :

\begin{enumerate}
\item \textbf{La prévention communautaire d'abord.} Au Cameroun, les relais communautaires sensibilisent les familles à l'usage des moustiquaires ; en Chine, les 社区卫生服务中心 organisent vaccination et dépistage de proximité. Dans les deux cas, la santé commence \emph{dans le quartier}, pas à l'hôpital.
\item \textbf{La valorisation des savoirs endogènes.} Le Cameroun reconnaît sa pharmacopée ; la Chine a institutionnalisé la médecine chinoise. On peut dire : \zh{非洲传统医药与中医药有很多相似之处。} \emph{Fēizhōu chuántǒng yīyào yǔ Zhōngyīyào yǒu hěn duō xiāngsì zhī chù.} « La médecine traditionnelle africaine et la médecine chinoise ont beaucoup de similitudes. »
\item \textbf{La coopération Sud-Sud.} Depuis plus de cinquante ans, la Chine envoie des équipes médicales au Cameroun (hôpitaux de district, campagnes de soins). \zh{中国派遣医疗队援助喀麦隆已有50多年历史。} \emph{Zhōngguó pàiqiǎn yīliáo duì yuánzhù Kāmàilóng yǐ yǒu wǔshí duō nián lìshǐ.} « La Chine envoie des équipes médicales aider le Cameroun depuis plus de 50 ans. »
\end{enumerate}

\courssub{Analyse des différences}

\begin{itemize}
\item \textbf{Échelle démographique} : la Chine compte plus d'un milliard d'habitants, le Cameroun environ 28 millions. Les outils chinois (maillage territorial dense) ne sont pas transposables tels quels.
\item \textbf{Centralisation vs décentralisation} : en Chine, la décision sanitaire est centralisée (Commission nationale de la santé) ; au Cameroun, l'action passe par les régions et les districts de santé.
\item \textbf{Place de la médecine traditionnelle} : \emph{intégrée} à l'hôpital en Chine (中医院), \emph{parallèle} au Cameroun (tradipraticiens reconnus mais hors du circuit hospitalier classique).
\end{itemize}

\courssub{Outils grammaticaux pour comparer et nuancer}

\begin{propriete}[La concession 虽然……但是……]
\textbf{Structure} : 虽然 + proposition 1 ，但是 + proposition 2.\\
\textbf{Emploi} : « bien que..., cependant... ». Contrairement au français, le chinois emploie \emph{les deux} marqueurs dans la même phrase ; en français, « bien que » et « mais » ne se combinent pas.\\
\textbf{Exemple} : \zh{虽然喀麦隆医疗资源有限，但是防疫意识在提高。} \emph{Suīrán Kāmàilóng yīliáo zīyuán yǒuxiàn, dànshì fángyì yìshi zài tígāo.} « Bien que les ressources médicales du Cameroun soient limitées, la conscience de la prévention progresse. »\\
\textbf{Exception} : à l'oral, on peut omettre 虽然 et garder seulement 但是 ; l'inverse (虽然 sans 但是) est aussi possible, mais jamais l'omission des deux.
\end{propriete}

\begin{propriete}[La comparaison 与……不同]
\textbf{Structure} : 与 + B + 不同 ，+ proposition.\\
\textbf{Emploi} : « différent de... / contrairement à... ». Très utile en introduction de comparaison.\\
\textbf{Exemple} : \zh{与中国不同，喀麦隆的疟疾仍然流行。} \emph{Yǔ Zhōngguó bùtóng, Kāmàilóng de nüèji réngrán liúxíng.} « Contrairement à la Chine, le paludisme reste endémique au Cameroun. »\\
\textbf{Variante} : A 和 B 不一样 (A hé B bù yíyàng) = « A et B ne sont pas pareils » (registre oral).\\
\textbf{Attention à l'ordre} : 与 + B + 不同 se place \emph{avant} la proposition principale, jamais après : on ne dit pas ✱\zh{喀麦隆的疟疾仍然流行与中国不同。}
\end{propriete}

\begin{attention}[Consignes de neutralité : éviter les généralisations]
En expression orale comme écrite, interdiction des jugements de valeur et des « tous / aucun ».
\begin{itemize}
\item FAUX : « Les Chinois ne croient qu'en la médecine chinoise. » → En réalité, la médecine conventionnelle (dite « occidentale ») domine dans les grands hôpitaux chinois ; la médecine chinoise est \emph{complémentaire}.
\item FAUX : « Les Camerounais refusent les vaccins. » → En réalité, la couverture du Programme Élargi de Vaccination atteint des niveaux élevés pour plusieurs antigènes.
\item DITES plutôt : \zh{一部分人更喜欢传统医学。} \emph{Yí bùfen rén gèng xǐhuan chuántǒng yīxué.} « Une partie de la population privilégie la médecine traditionnelle. » L'outil-clé : \cle{一部分} (une partie de), qui nuance toute généralisation.
\end{itemize}
Citez toujours vos sources de façon sobre : OMS, MINSANTE, Commission nationale de la santé (NHC).
\end{attention}

\begin{methode}[Remplir un tableau à double entrée en binôme]
\begin{enumerate}
\item Tracez un tableau à 4 colonnes : \textbf{Critère | Cameroun | Chine | Similitude ou différence}.
\item Choisissez 4 critères : maladies cibles, médecine traditionnelle, vaccination, rôle de l'État.
\item Remplissez chaque case avec \emph{un fait} (pas une opinion), en vous aidant du tableau de référence.
\item Utilisez le code couleurs : \textcolor{popblue}{bleu} = santé publique, \textcolor{popgreen}{vert} = médecine traditionnelle, \textcolor{poporange}{orange} = défis.
\item Formulez ensuite deux phrases en chinois : une avec 虽然……但是……, une avec 与……不同.
\end{enumerate}
\end{methode}

\begin{exoresolu}[Comparer en chinois]
\textbf{Énoncé.} À partir du tableau factuel, complétez les phrases suivantes avec 虽然……但是…… ou 与……不同, puis traduisez.\\
1. \zh{＿＿＿中国不同，喀麦隆的疟疾＿＿＿流行。}\\
2. \zh{＿＿＿喀麦隆的疫苗接种有进步，＿＿＿农村地区的医疗条件还需要改善。}
\tcblower
\textbf{Solution détaillée.}\\
1. La phrase oppose deux pays : il faut 与……不同. Réponse : \zh{与中国不同，喀麦隆的疟疾仍然流行。} \emph{Yǔ Zhōngguó bùtóng, Kāmàilóng de nüèji réngrán liúxíng.} « Contrairement à la Chine, le paludisme reste endémique au Cameroun. »\\
2. La phrase oppose un progrès et une limite dans le même pays : il faut la concession 虽然……但是……. Réponse : \zh{虽然喀麦隆的疫苗接种有进步，但是农村地区的医疗条件还需要改善。} \emph{Suīrán Kāmàilóng de yìmiáo jiēzhòng yǒu jìnbù, dànshì nóngcūn dìqū de yīliáo tiáojiàn hái xūyào gǎishàn.} « Bien que la vaccination progresse au Cameroun, les conditions médicales des zones rurales doivent encore être améliorées. »\\
\textbf{Piège déjoué} : ne traduisez jamais « bien que... mais... » mot à mot en français dans la même phrase ; en chinois, les deux marqueurs cohabitent, en français non.
\end{exoresolu}

\begin{exemplebox}[Deux phrases-modèles pour votre exposé]
\zh{中国派遣医疗队援助喀麦隆已有50多年历史。}\\
\emph{Zhōngguó pàiqiǎn yīliáo duì yuánzhù Kāmàilóng yǐ yǒu wǔshí duō nián lìshǐ.}\\
« La Chine envoie des équipes médicales aider le Cameroun depuis plus de 50 ans. »\\[4pt]
\zh{非洲传统医药与中医药有很多相似之处。}\\
\emph{Fēizhōu chuántǒng yīyào yǔ Zhōngyīyào yǒu hěn duō xiāngsì zhī chù.}\\
« La médecine traditionnelle africaine et la médecine chinoise ont beaucoup de similitudes. »
\end{exemplebox}

\begin{savaistu}[Le paludisme, une victoire chinoise en 70 ans]
Dans les années 1940, la Chine comptait environ 30 millions de cas de paludisme par an. Grâce aux stations anti-épidémiques (防疫站), à la distribution massive de moustiquaires et à l'artémisinine — molécule issue de la pharmacopée chinoise, découverte par l'équipe de Tu Youyou, prix Nobel de médecine 2015 — le pays a été certifié « sans paludisme » par l'OMS en 2021. Cette même artémisinine est aujourd'hui la base des traitements utilisés au Cameroun : un beau pont entre les savoirs chinois et la santé africaine.
\end{savaistu}

\begin{exercice}[À vous de comparer]
En binôme, rédigez trois phrases en chinois (caractères + pinyin) à partir du tableau factuel :\\
1. Une similitude avec \zh{都} (dōu).\\
2. Une différence avec \zh{与……不同}.\\
3. Une nuance avec \zh{虽然……但是……}.\\
Vous les présenterez oralement lors du « Forum Santé Jeunesse » de la section 5.
\end{exercice}

\begin{corrige}[Exercice « À vous de comparer » — propositions de réponses]
1. Similitude : \zh{两国都重视疫苗接种。} \emph{Liǎng guó dōu zhòngshì yìmiáo jiēzhòng.} « Les deux pays attachent de l'importance à la vaccination. »\\
2. Différence : \zh{与中国不同，喀麦隆的传统医学在医院之外。} \emph{Yǔ Zhōngguó bùtóng, Kāmàilóng de chuántǒng yīxué zài yīyuàn zhī wài.} « Contrairement à la Chine, la médecine traditionnelle camerounaise reste en dehors des hôpitaux. »\\
3. Nuance : \zh{虽然喀麦隆的疟疾仍然流行，但是预防工作越来越有效。} \emph{Suīrán Kāmàilóng de nüèji réngrán liúxíng, dànshì yùfáng gōngzuò yuè lái yuè yǒuxiào.} « Bien que le paludisme reste endémique au Cameroun, le travail de prévention devient de plus en plus efficace. »\\
\textbf{Critères de réussite} : place correcte de 都 (après le sujet), de 与……不同 (en tête de phrase) et des deux marqueurs 虽然……但是…… ; pinyin avec tons ; aucun jugement de valeur.
\end{corrige}

\coursec{Outils linguistiques : Comparer, Argumenter, Nuancer}

Dans la section précédente, nous avons dressé un tableau factuel des maladies endémiques au Cameroun et en Chine : paludisme, typhoïde, dengue, bilharziose. Mais un tableau ne suffit pas : pour participer au « Forum Santé Jeunesse », il faut savoir \cle{comparer} ces réalités, \cle{argumenter} et \cle{nuancer} son point de vue en chinois. C'est l'objet de cette section : transformer des faits en discours organisé. Nous procéderons en trois temps : d'abord les structures de comparaison, ensuite les connecteurs de l'argumentation, enfin les marqueurs du point de vue.

\courssub{Comparer : les structures de comparaison}

\textbf{Observation.}\par

Lisez d'abord ces quatre phrases tirées de notre thème. Ne cherchez pas encore la règle : repérez simplement la place de chaque mot.

\begin{exemplebox}[Observer avant d'apprendre la règle]
\zh{中国的疫苗接种率比喀麦隆高。} \emph{Zhōngguó de yìmiáo jiēzhònglǜ bǐ Kāmàilóng gāo.} « Le taux de vaccination de la Chine est plus élevé que celui du Cameroun. »\\
\zh{雅温得的蚊子跟农村的蚊子一样多。} \emph{Yǎwēndé de wénzi gēn nóngcūn de wénzi yíyàng duō.} « Il y a autant de moustiques à Yaoundé qu'à la campagne. »\\
\zh{疟疾是喀麦隆最常见的疾病。} \emph{Nüèji shì Kāmàilóng zuì chángjiàn de jíbìng.} « Le paludisme est la maladie la plus fréquente au Cameroun. »\\
\zh{预防像盾牌一样保护我们。} \emph{Yùfáng xiàng dùnpái yíyàng bǎohù wǒmen.} « La prévention nous protège comme un bouclier. »
\end{exemplebox}

Observez : dans la phrase 1, le mot \zh{比} se place \emph{avant} le terme de comparaison, et l'adjectif \zh{高} vient \emph{après}, sans \zh{很}. C'est l'inverse du français (« plus élevé \emph{que} »). Dans la phrase 2, l'égalité se construit avec \zh{跟……一样}. Dans la phrase 3, le superlatif utilise \zh{最} placé devant l'adjectif. Dans la phrase 4, la comparaison imagée (« comme ») passe par \zh{像……一样}.

\begin{propriete}[La phrase comparative : A 比 B + adjectif]
\textbf{Structure :} A + 比 (bǐ) + B + adjectif (+ complément de degré : 得多, 一点儿, 一些).\\
\textbf{Emploi :} exprimer la supériorité de A sur B. L'adjectif est \textbf{obligatoire} et ne prend \textbf{jamais} \zh{很} après \zh{比} (on dit \zh{比……高得多} « bien plus élevé », jamais *\zh{比……很高}).\\
\textbf{Négation :} A + 没有 (méiyǒu) + B + adjectif (« A n'est pas aussi… que B »), ou A + 比 + B + 不 + adjectif.\\
\textbf{Égalité :} A + 跟/和 + B + 一样 + adjectif ; différence : A + 跟 + B + 不一样.\\
\textbf{Similitude :} A + 像 (xiàng) + B + 一样 + verbe/adjectif (« comme »).\\
\textbf{Superlatif :} 最 (zuì) + adjectif ; comparatif renforcé : 更 (gèng) + adjectif (« encore plus »).
\end{propriete}

\begin{exemplebox}[Déclinaisons sur le thème de la santé]
\zh{杜阿拉比雅温得更潮湿。} \emph{Dùālā bǐ Yǎwēndé gèng cháoshī.} « Douala est encore plus humide que Yaoundé. »\\
\zh{喀麦隆的登革热病例没有亚洲多。} \emph{Kāmàilóng de dēnggérè bìnglì méiyǒu Yàzhōu duō.} « Le Cameroun n'a pas autant de cas de dengue que l'Asie. »\\
\zh{这种药比那种药便宜一点儿。} \emph{Zhè zhǒng yào bǐ nà zhǒng yào piányi yìdiǎnr.} « Ce médicament est un peu moins cher que celui-là. »\\
\zh{洗手跟戴口罩一样重要。} \emph{Xǐshǒu gēn dài kǒuzhào yíyàng zhòngyào.} « Se laver les mains est aussi important que porter un masque. »
\end{exemplebox}

\begin{attention}[Pièges des francophones avec 比]
\begin{itemize}
\item Erreur fréquente : *\zh{我觉得中国比喀麦隆好} — trop vague, presque impoli. Correct : \zh{我觉得中国在疟疾防治方面的经验比喀麦隆丰富。} \emph{Wǒ juéde Zhōngguó zài nüèji fángzhì fāngmiàn de jīngyàn bǐ Kāmàilóng fēngfù.} « Je pense que la Chine a une expérience plus riche que le Cameroun dans la lutte contre le paludisme. » \textbf{Précisez toujours le domaine de comparaison} avec \zh{在……方面}.
\item Ne pas oublier l'adjectif après \zh{比} : *\zh{中国比喀麦隆} ne veut rien dire.
\item Jamais de \zh{很} après \zh{比} : *\zh{比……很高} est fautif ; utilisez \zh{得多} ou \zh{多了}.
\item Attention à l'ordre : le français dit « A est plus grand \emph{que} B », le chinois dit « A \zh{比} B grand ». Ne traduisez jamais « que » par un mot placé après l'adjectif.
\end{itemize}
\end{attention}

\courssub{Argumenter : cause, accumulation, concession, balancement}

Pour convaincre, il ne suffit pas de comparer : il faut \cle{relier} ses idées. Le chinois dispose de couples de connecteurs dont les deux moitiés apparaissent dans la même phrase. Observons-les d'abord en contexte.

\begin{exemplebox}[Observer les connecteurs logiques]
\zh{由于气候潮湿，所以蚊子滋生快。} \emph{Yóuyú qìhòu cháoshī, suǒyǐ wénzi zīshēng kuài.} « En raison du climat humide, les moustiques prolifèrent vite. »\\
\zh{蚊帐不但便宜，而且有效。} \emph{Wénzhàng búdàn piányi, érqiě yǒuxiào.} « La moustiquaire est non seulement bon marché, mais aussi efficace. »\\
\zh{虽然中药有效，但是剂量要控制。} \emph{Suīrán zhōngyào yǒuxiào, dànshì jìliàng yào kòngzhì.} « Bien que la médecine chinoise soit efficace, il faut contrôler la dose. »\\
\zh{一方面要预防，另一方面要治疗。} \emph{Yì fāngmiàn yào yùfáng, lìng yì fāngmiàn yào zhìliáo.} « D'un côté il faut prévenir, de l'autre il faut soigner. »
\end{exemplebox}

\begin{propriete}[Les quatre couples de connecteurs logiques]
\begin{itemize}
\item \textbf{Cause → conséquence :} 由于/因为……所以/因此…… (\emph{yóuyú / yīnwèi … suǒyǐ / yīncǐ …}). En chinois, on peut garder les deux moitiés du couple dans la même phrase, contrairement au français (« parce que… donc » est impossible en français, normal en chinois).
\item \textbf{Accumulation :} 不但……而且…… (\emph{búdàn … érqiě …}) « non seulement… mais aussi… ». Si les deux sujets sont identiques, le sujet se place avant \zh{不但}.
\item \textbf{Concession :} 虽然……但是/可是…… (\emph{suīrán … dànshì / kěshì …}) « bien que… (pourtant)… ». Les deux membres coexistent en chinois.
\item \textbf{Balancement :} 一方面……另一方面…… (\emph{yì fāngmiàn … lìng yì fāngmiàn …}) « d'une part… d'autre part… », idéal pour présenter les deux faces d'un problème de santé publique.
\end{itemize}
\end{propriete}

\begin{center}
\begin{tabularx}{\linewidth}{|X|X|X|}
\hline
\rowcolor{popblueL} Structure & Exemple (caractères + pinyin) & Traduction \\
\hline
A 比 B + Adj & \zh{城市的医生比乡村多。} \emph{Chéngshì de yīshēng bǐ xiāngcūn duō.} & Il y a plus de médecins en ville qu'à la campagne. \\
\hline
A 跟 B 一样 + Adj & \zh{预防跟治疗一样重要。} \emph{Yùfáng gēn zhìliáo yíyàng zhòngyào.} & Prévenir est aussi important que soigner. \\
\hline
因为……所以…… & \zh{因为下雨多，所以疟疾常见。} \emph{Yīnwèi xiàyǔ duō, suǒyǐ nüèji chángjiàn.} & Comme il pleut beaucoup, le paludisme est fréquent. \\
\hline
虽然……但是…… & \zh{虽然药贵，但是必须买。} \emph{Suīrán yào guì, dànshì bìxū mǎi.} & Bien que le médicament soit cher, il faut l'acheter. \\
\hline
不但……而且…… & \zh{他不但研究中医，而且研究西医。} \emph{Tā búdàn yánjiū zhōngyī, érqiě yánjiū xīyī.} & Il étudie non seulement la médecine chinoise, mais aussi la médecine occidentale. \\
\hline
一方面……另一方面…… & \zh{一方面要灭蚊，另一方面要教育。} \emph{Yì fāngmiàn yào miè wén, lìng yì fāngmiàn yào jiàoyù.} & D'une part éliminer les moustiques, d'autre part éduquer. \\
\hline
\end{tabularx}
\end{center}

\begin{popfigure}
\begin{tikzpicture}[mindmap, grow cyclic, every node/.style=concept, concept color=popblue!20, level 1/.append style={level distance=3cm,sibling angle=90}]
\node {Argumenter 论证}
child[concept color=popgreen!30] {node {Comparer 比较}
child {node {比}}
child {node {像}}
child {node {一样}}
}
child[concept color=poporange!30] {node {Cause/Effet 因果}
child {node {因为……所以}}
child {node {由于……因此}}
}
child[concept color=poppink!40] {node {Concession 让步}
child {node {虽然……但是}}
child {node {尽管……还是}}
}
child[concept color=poppurple!30] {node {Opinion 观点}
child {node {我认为}}
child {node {依我看}}
};
\end{tikzpicture}
\legende{Carte mentale des outils pour argumenter en chinois.}
\end{popfigure}

\courssub{Nuancer : exprimer un point de vue organisé}

Comparer et relier ne suffisent pas encore : un bon exposé annonce aussi \emph{qui parle} et \emph{dans quel ordre} il parle. C'est le rôle des marqueurs d'opinion et des connecteurs de plan.

\begin{propriete}[Marqueurs du point de vue et connecteurs de discours]
\textbf{Opinion :} 我认为 (wǒ rènwéi) / 我觉得 (wǒ juéde) / 在我看来 (zài wǒ kànlái) / 依我看 (yī wǒ kàn) + point de vue + 因为 + raison + 所以 + conclusion.\\
\textbf{Plan du discours :} 首先 (shǒuxiān) « d'abord », 其次 (qícì) « ensuite », 最后 (zuìhòu) « enfin », 总之 (zǒngzhī) « en somme ».\\
\textbf{Nuancer :} ajouter \zh{比较} « assez », \zh{可能} « peut-être », \zh{一般来说} « en général » pour éviter les affirmations absolues : \zh{一般来说，中药比较温和。} \emph{Yìbān lái shuō, zhōngyào bǐjiào wēnhé.} « En général, la médecine chinoise est assez douce. »
\end{propriete}

\begin{methode}[Construire un mini-exposé en cinq étapes]
\begin{enumerate}
\item \textbf{Accroche} par un fait : \zh{疟疾是喀麦隆最常见的疾病。}
\item \textbf{Annonce du plan} : \zh{我想谈两个方面。} \emph{Wǒ xiǎng tán liǎng ge fāngmiàn.} « Je veux parler de deux aspects. »
\item \textbf{Argument 1} : \zh{首先，因为气候潮湿，所以蚊子很多。}
\item \textbf{Argument 2} : \zh{其次，虽然医院越来越多，但是农村还不够。}
\item \textbf{Conclusion} : \zh{总之，我建议多用蚊帐，多洗手。} \emph{Zǒngzhī, wǒ jiànyì duō yòng wénzhàng, duō xǐshǒu.} « En somme, je conseille d'utiliser plus de moustiquaires et de se laver les mains souvent. »
\end{enumerate}
\end{methode}

\begin{exoresolu}[Transformation : phrases disjointes → paragraphe cohérent]
Énoncé : reliez ces trois phrases avec les connecteurs appropriés : \zh{气候潮湿。蚊子很多。疟疾常见。}
\tcblower
Solution détaillée : on identifie la cause (\zh{气候潮湿}), l'effet intermédiaire (\zh{蚊子很多}) et la conséquence finale (\zh{疟疾常见}). On applique le couple 因为……所以…… et le connecteur 因此 :\\
\zh{因为气候潮湿，所以蚊子很多，因此疟疾很常见。} \emph{Yīnwèi qìhòu cháoshī, suǒyǐ wénzi hěn duō, yīncǐ nüèji hěn chángjiàn.} « Comme le climat est humide, il y a beaucoup de moustiques, et c'est pourquoi le paludisme est très fréquent. »\\
Vérification : l'ordre chinois Cause → Conséquence est respecté ; le sujet n'est pas répété inutilement.
\end{exoresolu}

\begin{exoresolu}[Thème guidé : français → chinois]
Énoncé : traduisez « Bien que le paludisme soit dangereux, on peut le prévenir avec des moustiquaires. »
\tcblower
Solution : on repère la concession → 虽然……但是…… ; « prévenir » = \zh{预防} ; « moustiquaire » = \zh{蚊帐}.\\
\zh{虽然疟疾很危险，但是用蚊帐可以预防。} \emph{Suīrán nüèji hěn wēixiǎn, dànshì yòng wénzhàng kěyǐ yùfáng.}\\
Piège évité : ne pas traduire « on peut » par un sujet \zh{人们} superflu ; le chinois préfère la tournure sans sujet avec \zh{可以}.
\end{exoresolu}

\begin{exercice}[À vous de jouer]
\begin{enumerate}
\item Transformez en phrase comparative avec \zh{比} : \zh{中国的医院多。喀麦隆的医院少。}
\item Complétez avec 不但……而且…… : \zh{这种药\_\_\_\_便宜，\_\_\_\_有效。}
\item Traduisez : « D'une part il faut des vaccins, d'autre part il faut de l'éducation sanitaire. »
\item Rédigez trois phrases avec 首先, 其次, 总之 sur la prévention de la typhoïde à Douala.
\end{enumerate}
\end{exercice}

\begin{corrige}[Exercice « À vous de jouer »]
\begin{enumerate}
\item \zh{中国的医院比喀麦隆多。} \emph{Zhōngguó de yīyuàn bǐ Kāmàilóng duō.} « Il y a plus d'hôpitaux en Chine qu'au Cameroun. »
\item \zh{这种药不但便宜，而且有效。} \emph{Zhè zhǒng yào búdàn piányi, érqiě yǒuxiào.} « Ce médicament est non seulement bon marché, mais aussi efficace. »
\item \zh{一方面需要疫苗，另一方面需要卫生教育。} \emph{Yì fāngmiàn xūyào yìmiáo, lìng yì fāngmiàn xūyào wèishēng jiàoyù.} « D'une part il faut des vaccins, d'autre part il faut de l'éducation sanitaire. »
\item Modèle : \zh{首先，要喝干净的水。其次，要常洗手。总之，预防比治疗好。} \emph{Shǒuxiān, yào hē gānjìng de shuǐ. Qícì, yào cháng xǐshǒu. Zǒngzhī, yùfáng bǐ zhìliáo hǎo.} « D'abord, boire de l'eau propre. Ensuite, se laver souvent les mains. En somme, prévenir vaut mieux que guérir. »
\end{enumerate}
\end{corrige}

\courssub{Focus prononciation : le troisième ton en séquence}

Le troisième ton est un ton \cle{bas-descendant-remontant}, mais dans la parole rapide il devient souvent un simple ton bas. Règle essentielle : quand deux troisièmes tons se suivent, le premier devient un deuxième ton : \zh{比较} se prononce \emph{bíjiào} (et non bǐjiào) dans la chaîne parlée, de même \zh{所以} → \emph{suóyǐ}, \zh{我想} → \emph{wó xiǎng}. À l'oral, placez l'\cle{accent de phrase} sur le mot-clé de votre argument : dans \zh{我认为预防最重要} \emph{wǒ rènwéi yùfáng zuì zhòngyào} « je pense que la prévention est le plus important », c'est \zh{预防} qui porte la force du message.

\begin{savaistu}[L'ordre logique chinois : Cause → Conséquence]
En chinois, l'ordre logique naturel est presque toujours « Cause → Conséquence » (\zh{由于……因此……}), alors qu'en français on inverse volontiers (« Il y a beaucoup de moustiques \emph{parce que} le climat est humide »). À l'oral, respectez l'ordre chinois : votre exposé gagnera en fluidité et paraîtra plus naturel aux oreilles d'un sinophone. C'est aussi l'ordre du raisonnement dans la tradition argumentative chinoise : on pose d'abord le contexte, puis la conclusion.
\end{savaistu}

\begin{aretenir}[L'essentiel de la section]
\begin{itemize}
\item Comparaison : A 比 B + Adj (sans \zh{很}) ; égalité : A 跟 B 一样 ; superlatif : 最 ; similitude : 像……一样.
\item Argumentation : 因为……所以……, 不但……而且……, 虽然……但是……, 一方面……另一方面…… — les couples se gardent entiers en chinois.
\item Point de vue : 我认为 / 依我看 + 因为 + raison + 所以 + conclusion ; plan : 首先, 其次, 最后, 总之.
\item Nuancez avec 比较, 可能, 一般来说 ; précisez le domaine avec 在……方面.
\end{itemize}
\end{aretenir}

\coursec{Atelier oral : Jeu de rôle \& Débat guidé « Forum Santé Jeunesse »}

Dans la section précédente, nous avons appris à \cle{comparer}, à \cle{argumenter} et à \cle{nuancer} en chinois : \zh{虽然……但是……}, \zh{一方面……另一方面……}, \zh{我认为……}. Ces outils ne valent que s'ils sont \emph{utilisés}. C'est exactement l'objectif de cet atelier : passer de la phrase écrite à la \cle{prise de parole en public}, dans une situation réaliste et motivante — un forum jeunesse Cameroun--Chine consacré à la santé. Vous allez incarner un rôle, défendre un point de vue, réagir à vos camarades : bref, \emph{parler pour convaincre}.

\courssub{Mise en situation : le « Forum Santé Jeunesse »}

\begin{experience}[Contexte de la simulation]
Le lycée organise un \cle{forum jeunesse Cameroun--Chine} sur le thème : « \zh{传统医学和现代医学：对手还是伙伴？} » (\emph{Chuántǒng yīxué hé xiàndài yīxué : duìshǒu háishi huǒbàn ?} — « Médecine traditionnelle et moderne : rivales ou partenaires ? »). La classe est divisée en groupes de six élèves. Chaque groupe forme un plateau de débat avec les rôles suivants :
\begin{itemize}
\item \textbf{Modérateur / 主持人} (\emph{zhǔchírén}) : il ouvre le forum, donne la parole, relance, conclut.
\item \textbf{Expert OMS / 世卫组织专家} (\emph{Shìwèi zǔzhī zhuānjiā}) : il apporte des faits sur les maladies endémiques (paludisme, dengue) et la coopération sanitaire.
\item \textbf{Médecin traditionnel camerounais / 喀麦隆传统医生} (\emph{Kāmàilóng chuántǒng yīshēng}) : il défend les plantes médicinales locales et l'expérience des communautés.
\item \textbf{Médecin chinois (MTC) / 中医师} (\emph{zhōngyīshī}) : il présente l'acupuncture, la pharmacopée chinoise et l'exemple de l'artémisinine (\zh{青蒿素}, \emph{qīnghāosù}).
\item \textbf{Jeune patient / 年轻病人} (\emph{niánqīng bìngrén}) : il raconte son expérience du paludisme et pose des questions concrètes (prix, accès, confiance).
\item \textbf{Journaliste / 记者} (\emph{jìzhě}) : il interviewe les experts et prépare un compte rendu.
\end{itemize}
\end{experience}

\begin{popfigure}
\begin{tikzpicture}
\node[draw, circle, fill=popgold!40, align=center, font=\bfseries] (mod) at (0,0) {主持人\\ Modérateur};
\foreach \angle/\name in {60/{Expert OMS}, 120/{Médecin Trad. CM}, 180/{Médecin MTC}, 240/{Jeune patient}, 300/{Journaliste}}
\node[draw, rectangle, rounded corners, fill=popblue!15, align=center] at (\angle:3.2) {\name};
\draw[popline, thick] (mod) -- (60:3.2) (mod) -- (120:3.2) (mod) -- (180:3.2) (mod) -- (240:3.2) (mod) -- (300:3.2);
\end{tikzpicture}
\legende{Organisation du plateau : le modérateur au centre distribue la parole aux cinq intervenants.}
\end{popfigure}

\courssub{Fiches de rôle : arguments clés et vocabulaire imposé}

Chaque intervenant reçoit une \cle{fiche de rôle} (\zh{角色卡}, \emph{juésè kǎ}). Règle du jeu : chaque intervention doit contenir \textbf{au moins cinq mots du vocabulaire imposé} de la fiche. Voici la fiche du journaliste, à titre de modèle ; l'enseignant prépare des fiches équivalentes pour les autres rôles.

\begin{textebox}[角色卡 — 青年记者 (Fiche de rôle — le jeune journaliste)]
角色：青年记者。\\
\emph{Juésè : qīngnián jìzhě.}\\
Rôle : jeune journaliste.\\
任务：采访专家，提问：\\
\emph{Rènwù : cǎifǎng zhuānjiā, tíwèn :}\\
Tâche : interviewer les experts, poser les questions :\\
“中非医疗合作最大的困难是什么？”\\
\emph{« Zhōng-Fēi yīliáo hézuò zuìdà de kùnnan shì shénme ? »}\\
« Quelle est la plus grande difficulté de la coopération médicale sino-africaine ? »\\
“青蒿素对非洲意味着什么？”\\
\emph{« Qīnghāosù duì Fēizhōu yìwèizhe shénme ? »}\\
« Que signifie l'artémisinine pour l'Afrique ? »\\
关键词：合作、障碍、意义、推广、信任。\\
\emph{Guānjiàncí : hézuò, zhàng'ài, yìyì, tuīguǎng, xìnrèn.}\\
Mots-clés : coopération, obstacle, signification, diffusion, confiance.
\end{textebox}

\begin{center}
\begin{tabularx}{\linewidth}{|X|X|X|X|}
\hline
\rowcolor{popblueL} Caractères & Pinyin & Nature & Traduction \\
\hline
采访 & cǎifǎng & verbe & interviewer \\
\hline
提问 & tíwèn & verbe / nom & poser une question ; question \\
\hline
合作 & hézuò & verbe / nom & coopérer ; coopération \\
\hline
障碍 & zhàng'ài & nom & obstacle \\
\hline
意义 & yìyì & nom & signification, portée \\
\hline
推广 & tuīguǎng & verbe & diffuser, promouvoir \\
\hline
信任 & xìnrèn & verbe / nom & faire confiance ; confiance \\
\hline
负担 & fùdān & nom & fardeau, charge \\
\hline
副作用 & fùzuòyòng & nom & effet secondaire \\
\hline
经验 & jīngyàn & nom & expérience \\
\hline
\end{tabularx}
\end{center}

Suggestions d'arguments par rôle (à distribuer) :
\begin{itemize}
\item \textbf{Expert OMS} : le paludisme reste un fardeau lourd en Afrique ; la prévention (moustiquaires, dépistage) et la coopération internationale sauvent des vies. Mots imposés : \zh{预防} (\emph{yùfáng}, prévenir), \zh{负担} (\emph{fùdān}), \zh{合作} (\emph{hézuò}), \zh{数据} (\emph{shùjù}, données), \zh{疫苗} (\emph{yìmiáo}, vaccin).
\item \textbf{Médecin traditionnel camerounais} : les plantes locales sont accessibles et peu chères ; il faut les étudier scientifiquement. Mots imposés : \zh{草药} (\emph{cǎoyào}, plantes médicinales), \zh{便宜} (\emph{piányi}, bon marché), \zh{经验} (\emph{jīngyàn}), \zh{研究} (\emph{yánjiū}, recherche), \zh{信任} (\emph{xìnrèn}).
\item \textbf{Médecin MTC} : la médecine chinoise traite la cause et pas seulement le symptôme ; l'artémisinine, issue d'une plante, montre que tradition et science peuvent travailler ensemble. Mots imposés : \zh{中医} (\emph{zhōngyī}), \zh{青蒿素} (\emph{qīnghāosù}), \zh{结合} (\emph{jiéhé}, combiner), \zh{副作用} (\emph{fùzuòyòng}), \zh{推广} (\emph{tuīguǎng}).
\item \textbf{Jeune patient} : il veut des soins efficaces, rapides et abordables ; il demande des preuves. Mots imposés : \zh{发烧} (\emph{fāshāo}, avoir de la fièvre), \zh{效果} (\emph{xiàoguǒ}, effet), \zh{价格} (\emph{jiàgé}, prix), \zh{医院} (\emph{yīyuàn}, hôpital), \zh{担心} (\emph{dānxīn}, s'inquiéter).
\end{itemize}

\courssub{La boîte à phrases du débatteur}

Avant de lancer le forum, observons comment un locuteur chinois \emph{organise} sa prise de parole. Lisez d'abord ce mini-discours :

\begin{exemplebox}[Observer avant de parler]
\zh{关于传统医学和现代医学的问题，我认为它们是伙伴，不是对手。原因有二：首先，传统医学有经验；其次，现代医学有科学方法。当然，也有人担心草药的安全，但我认为，只要结合起来，就能更好地保护病人。综上所述，合作比竞争更重要。}\\
\emph{Guānyú chuántǒng yīxué hé xiàndài yīxué de wèntí, wǒ rènwéi tāmen shì huǒbàn, bú shì duìshǒu. Yuányīn yǒu èr : shǒuxiān, chuántǒng yīxué yǒu jīngyàn; qícì, xiàndài yīxué yǒu kēxué fāngfǎ. Dāngrán, yě yǒu rén dānxīn cǎoyào de ānquán, dàn wǒ rènwéi, zhǐyào jiéhé qǐlai, jiù néng gèng hǎo de bǎohù bìngrén. Zōngshàngsuǒshù, hézuò bǐ jìngzhēng gèng zhòngyào.}\\
« Concernant la question de la médecine traditionnelle et de la médecine moderne, je pense qu'elles sont partenaires, pas rivales. Il y a deux raisons : d'abord, la médecine traditionnelle possède l'expérience ; ensuite, la médecine moderne possède la méthode scientifique. Bien sûr, certains s'inquiètent de la sécurité des plantes, mais je pense que, pourvu qu'on les combine, on protège mieux les malades. En résumé, la coopération est plus importante que la compétition. »
\end{exemplebox}

Qu'observe-t-on ? Le discours suit un \emph{plan annoncé} : thème → opinion → deux raisons → concession → réponse → conclusion. C'est exactement la structure ci-dessous.

\begin{propriete}[Structure type d'une prise de position argumentée]
Schéma : \zh{关于} + thème + \zh{的问题，我认为} + opinion. \zh{原因有二：首先……其次……。当然，也有人担心……，但我认为……。综上所述……}\\
« Concernant la question de…, je pense que… Il y a deux raisons : primo…, secundo… Bien sûr, certains s'inquiètent de…, mais je pense que… En résumé… »
\begin{itemize}
\item \zh{关于……的问题} : cadre le sujet (équivalent de « en ce qui concerne »). Ne pas dire \zh{*} « 对于……问题 » sans \zh{的} : la forme correcte est \zh{关于……的问题}.
\item \zh{首先……其次……} : ordonne les arguments (primo / secundo). On peut ajouter \zh{最后} (\emph{zuìhòu}, enfin).
\item \zh{当然，也有人担心……} : concession, très valorisée à l'oral (montre que l'on écoute l'autre).
\item \zh{综上所述} (\emph{Zōngshàngsuǒshù}) : formule de conclusion soutenue, parfaite pour un exposé ; à l'oral familier, on préfère \zh{总之} (\emph{zǒngzhī}, bref).
\end{itemize}
\end{propriete}

\begin{definition}[Réagir et interagir : les formules du débat]
\begin{itemize}
\item \textbf{Exprimer l'accord} : \zh{我赞同……} (\emph{Wǒ zàntóng…}, je suis d'accord avec…) ; \zh{我完全同意您的看法。} (\emph{Wǒ wánquán tóngyì nín de kànfǎ.}, je partage entièrement votre avis).
\item \textbf{Exprimer le désaccord poli} : \zh{我不同意，因为……} (\emph{Wǒ bù tóngyì, yīnwèi…}, je ne suis pas d'accord, parce que…). En chinois, on justifie \emph{toujours} un désaccord : dire seulement \zh{我不同意} paraît brutal.
\item \textbf{Ménager une transition / calmer le jeu} : \zh{这个问题很有意思，不过……} (\emph{Zhège wèntí hěn yǒu yìsi, búguò…}, cette question est très intéressante, cependant…).
\item \textbf{Demander une clarification} : \zh{请问，您说的“结合”具体指什么？} (\emph{Qǐngwèn, nín shuō de “jiéhé” jùtǐ zhǐ shénme ?}, excusez-moi, que voulez-vous dire exactement par « combiner » ?).
\end{itemize}
\end{definition}

\begin{center}
\begin{tabularx}{\linewidth}{|X|X|X|}
\hline
\rowcolor{popblueL} Fonction & Chinois + pinyin & Traduction \\
\hline
Accord & 我赞同您的观点。Wǒ zàntóng nín de guāndiǎn. & Je partage votre point de vue. \\
\hline
Désaccord justifié & 我不同意，因为…… Wǒ bù tóngyì, yīnwèi… & Je ne suis pas d'accord, parce que… \\
\hline
Concession & 虽然您说得对，但是…… Suīrán nín shuō de duì, dànshì… & Bien que vous ayez raison, mais… \\
\hline
Transition & 这个问题很有意思，不过…… Zhège wèntí hěn yǒu yìsi, búguò… & Question intéressante, cependant… \\
\hline
Clarification & 请问，您说的……指什么？ Qǐngwèn, nín shuō de… zhǐ shénme ? & Que signifie précisément… ? \\
\hline
Donner la parole & 请说。 Qǐng shuō. & Je vous en prie, parlez. \\
\hline
Terminer son tour & 我说完了。 Wǒ shuō wán le. & J'ai fini. \\
\hline
\end{tabularx}
\end{center}

\courssub{Déroulement de la séance}

\begin{methode}[Les trois temps du forum]
\begin{enumerate}
\item \textbf{Exposés liminaires (6 min)} : chaque intervenant parle \textbf{une minute} pour présenter son rôle et sa position. Le modérateur chronomètre et coupe poliment : \zh{时间到了，谢谢。} (\emph{Shíjiān dào le, xièxie.}, le temps est écoulé, merci).
\item \textbf{Table ronde modérée (10 min)} : thème « Médecine traditionnelle et moderne : rivales ou partenaires ? ». Le modérateur distribue la parole, impose au moins une \emph{réaction} par intervenant (accord, désaccord ou question de clarification). Le journaliste pose ses deux questions préparées.
\item \textbf{Synthèse collective (5 min)} : au tableau, la classe construit un tableau « Arguments pour / Arguments contre / Points de convergence ». Le journaliste de chaque groupe lit son compte rendu en trois phrases.
\end{enumerate}
\end{methode}

\begin{methode}[La technique « 4-3-2 » pour gagner en fluidité]
\begin{enumerate}
\item Préparez un exposé de \textbf{4 minutes} sur votre rôle, avec des cartes de mots-clés (pas de texte rédigé !).
\item Présentez-le à votre binôme en \textbf{3 minutes} : vous devez garder l'essentiel et parler plus vite, sans bafouiller.
\item Recommencez en \textbf{2 minutes} : seuls les arguments forts survivent ; la prononciation et le débit deviennent naturels.
\item Chronométrez-vous mutuellement, puis changez de rôle et recommencez.
\end{enumerate}
Cette méthode, très utilisée en didactique des langues, transforme un texte « appris » en parole \emph{spontanée}.
\end{methode}

\begin{exemplebox}[Ouverture du débat : modèle d'échange]
主持人：今天我们讨论中非医疗合作。\\
\emph{Zhǔchírén : Jīntiān wǒmen tǎolùn Zhōng-Fēi yīliáo hézuò.}\\
Modérateur : aujourd'hui, nous discutons de la coopération médicale sino-africaine.\\
专家：我认为合作很有必要，因为非洲疾病负担重，中国经验丰富。\\
\emph{Zhuānjiā : Wǒ rènwéi hézuò hěn yǒu bìyào, yīnwèi Fēizhōu jíbìng fùdān zhòng, Zhōngguó jīngyàn fēngfù.}\\
Expert : je pense que la coopération est très nécessaire, car le fardeau des maladies est lourd en Afrique et l'expérience de la Chine est riche.\\
传统医师：虽然西药快，但中药和草药治本，副作用小。\\
\emph{Chuántǒng yīshī : Suīrán xīyào kuài, dàn zhōngyào hé cǎoyào zhì běn, fùzuòyòng xiǎo.}\\
Médecin traditionnel : bien que les médicaments occidentaux soient rapides, la médecine chinoise et les plantes traitent la racine du mal, avec peu d'effets secondaires.
\end{exemplebox}

\begin{attention}[Les quatre pièges de l'oral en public]
\begin{itemize}
\item \textbf{Ne lisez pas vos notes.} Un texte lu tue l'interaction et fait chuter la note « interaction ». N'emportez que des \emph{cartes de mots-clés} : 合作、负担、首先、综上所述…
\item \textbf{Regardez votre interlocuteur}, pas votre feuille ni le plafond. En Chine comme au Cameroun, le regard marque le respect et la conviction.
\item \textbf{Gérez le tour de parole} avec les formules rituelles : \zh{请说} (\emph{Qǐng shuō}, je vous en prie) pour céder la parole, \zh{我说完了} (\emph{Wǒ shuō wán le}, j'ai fini) pour la rendre. Ne coupez jamais la parole sans formule : utilisez \zh{对不起，我想补充一点。} (\emph{Duìbuqǐ, wǒ xiǎng bǔchōng yìdiǎn.}, excusez-moi, je voudrais ajouter un point).
\item \textbf{Attention aux tons dans les mots-clés} : \zh{合作} \emph{hézuò} (2\up{e}-4\up{e} ton), \zh{意义} \emph{yìyì} (4\up{e}-4\up{e} ton). Un ton faux sur un mot répété dix fois devient très gênant : entraînez les mots imposés \emph{avant} le forum.
\end{itemize}
\end{attention}

\courssub{Évaluation et auto-évaluation}

\begin{exoresolu}[Construire une intervention complète]
Énoncé : vous êtes le médecin MTC. Rédigez puis prononcez une intervention de 4 à 5 phrases avec la structure du débat, en utilisant au moins cinq mots imposés : \zh{中医、青蒿素、结合、副作用、推广}.
\tcblower
Solution détaillée :\\
\zh{关于传统医学和现代医学的问题，我认为中医很有价值。原因有二：首先，青蒿素来自传统草药，救了很多人；其次，中药副作用小。当然，也有人担心中药的效果，但我认为，只要把中医和现代医学结合起来，就能更好地帮助病人。综上所述，我们应该推广这种合作。我说完了。}\\
\emph{Guānyú chuántǒng yīxué hé xiàndài yīxué de wèntí, wǒ rènwéi zhōngyī hěn yǒu jiàzhí. Yuányīn yǒu èr : shǒuxiān, qīnghāosù láizì chuántǒng cǎoyào, jiù le hěn duō rén; qícì, zhōngyào fùzuòyòng xiǎo. Dāngrán, yě yǒu rén dānxīn zhōngyào de xiàoguǒ, dàn wǒ rènwéi, zhǐyào bǎ zhōngyī hé xiàndài yīxué jiéhé qǐlai, jiù néng gèng hǎo de bāngzhù bìngrén. Zōngshàngsuǒshù, wǒmen yīnggāi tuīguǎng zhè zhǒng hézuò. Wǒ shuō wán le.}\\
« Concernant la question…, je pense que la MTC a une grande valeur. Deux raisons : d'abord, l'artémisinine vient d'une plante traditionnelle et a sauvé beaucoup de vies ; ensuite, les remèdes chinois ont peu d'effets secondaires. Certes, certains doutent de leur efficacité, mais je pense que, si l'on combine MTC et médecine moderne, on aide mieux les malades. En résumé, nous devons promouvoir cette coopération. J'ai fini. »\\
Vérification : structure complète (关于 → 首先/其次 → 当然 → 综上所述), cinq mots imposés employés, formule de clôture présente.
\end{exoresolu}

\begin{center}
\begin{tabularx}{\linewidth}{|X|X|X|}
\hline
\rowcolor{popblueL} Critère & Ce qui est attendu & Points \\
\hline
Prononciation et tons & Tons corrects sur les mots-clés, débit naturel, pas de lecture & /4 \\
\hline
Vocabulaire & Au moins 5 mots imposés, employés correctement & /4 \\
\hline
Structures & Structure du débat respectée (关于…首先…其次…综上所述) & /4 \\
\hline
Interaction & Réagit aux autres (accord, désaccord, question), gère son tour de parole & /4 \\
\hline
Contenu factuel & Arguments exacts et respectueux (faits Cameroun/Chine vus en section 3) & /4 \\
\hline
\end{tabularx}
\end{center}

\begin{experience}[Enregistrement audio et auto-évaluation différée]
Si l'équipement le permet, chaque groupe \textbf{enregistre} sa table ronde (téléphone, dictaphone). Le soir même, chaque élève réécoute sa prestation et remplit une fiche d'auto-évaluation : (1) trois tons que j'ai ratés ; (2) deux formules du débat que j'ai bien utilisées ; (3) un moment où j'aurais dû réagir ; (4) mon objectif pour le prochain forum. Cette écoute différée est l'un des moyens les plus efficaces de progresser en prononciation : on s'entend enfin « de l'extérieur ».
\end{experience}

\begin{exercice}[À vous de jouer]
\begin{enumerate}
\item En binôme, appliquez la technique 4-3-2 à votre rôle, puis enregistrez la version finale de 2 minutes.
\item Rédigez sur carte vos cinq mots-clés (caractères + pinyin) et deux formules d'interaction que vous vous engagez à utiliser.
\item Pendant le forum, notez dans un tableau les arguments « pour », « contre » et les « points de convergence » ; servez-vous-en pour la synthèse collective au tableau.
\end{enumerate}
\end{exercice}

\begin{corrige}[Exercice n 2 — exemple de carte de mots-clés (rôle : jeune patient)]
Carte recto : \zh{发烧 fāshāo — 效果 xiàoguǒ — 价格 jiàgé — 担心 dānxīn — 医院 yīyuàn}.\\
Carte verso : \zh{我不同意，因为……} / \zh{请问，您说的……指什么？}\\
Critères de réussite : pinyin avec tons corrects (\emph{fāshāo} : 1\up{er}-1\up{er} ; \emph{yīyuàn} : 1\up{er}-4\up{e}), formules complètes et mémorisées, aucune phrase rédigée en entier sur la carte — seulement des déclencheurs.
\end{corrige}

\begin{aretenir}[L'essentiel de l'atelier]
Un bon débatteur en chinois : (1) annonce son plan avec \zh{关于……的问题，我认为……} ; (2) ordonne avec \zh{首先……其次……} ; (3) concède avec \zh{当然，也有人担心……} ; (4) conclut avec \zh{综上所述} (\emph{Zōngshàngsuǒshù}) ; (5) interagit avec \zh{我赞同……} / \zh{我不同意，因为……} ; (6) gère la parole avec \zh{请说} et \zh{我说完了}. La technique 4-3-2 et l'écoute de son propre enregistrement transforment l'effort en progrès mesurable.
\end{aretenir}

\coursec{Production écrite intégrée \& Évaluation formative}

Après le « Forum Santé Jeunesse », où vous avez défendu vos idées à l'oral, il est temps de fixer ces idées par écrit. Comme au baccalauréat, l'expression écrite en chinois exige un plan clair, un vocabulaire précis et des structures de comparaison maîtrisées. Cette dernière séquence transforme tout ce que vous avez appris dans la leçon en une production écrite complète, évaluée selon une grille transparente.

\courssub{La tâche : rédiger un article de blog}

\begin{definition}[Sujet de production écrite]
Rédiger un article de blog de \cle{120 à 150 caractères chinois} intitulé \zh{我对抗疟斗争的看法：喀麦隆与中国的互鉴} (\emph{Wǒ duì kàngnüè dòuzhēng de kànfǎ : Kāmàilóng yǔ Zhōngguó de hùjiàn} — « Mon regard sur la lutte antipaludique : leçons croisées Cameroun-Chine »). Structure imposée : \cle{Titre} ; \cle{Introduction} (faits) ; \cle{Comparaison} (2 points) ; \cle{Avis personnel} (une proposition) ; \cle{Conclusion}.
\end{definition}

Avant d'écrire, observons un modèle rédigé par une élève de Terminale, que nous analyserons ensemble.

\begin{textebox}[Modèle d'article de blog]
我对抗疟斗争的看法\\
\emph{Wǒ duì kàngnüè dòuzhēng de kànfǎ}\\
« Mon regard sur la lutte antipaludique »\\
喀麦隆是疟疾高负担国家，每年很多儿童感染疟疾。2021年，中国获得了世界卫生组织的无疟疾认证。\\
\emph{Kāmàilóng shì nüèjí gāo fùdān guójiā, měi nián hěn duō értóng gǎnrǎn nüèjí. 2021 nián, Zhōngguó huòdé le Shìjiè Wèishēng Zǔzhī de wú nüèjí rènzhèng.}\\
« Le Cameroun est un pays à forte charge de paludisme ; chaque année, de nombreux enfants sont infectés. En 2021, la Chine a obtenu la certification "sans paludisme" de l'OMS. »\\
两国的情况虽然不同，但是可以互相学习。喀麦隆推广蚊帐和疫苗接种，中国不但重视环境管理，而且发展了青蒿素药物。\\
\emph{Liǎng guó de qíngkuàng suīrán bùtóng, dànshì kěyǐ hùxiāng xuéxí. Kāmàilóng tuīguǎng wénzhàng hé yìmiáo jiēzhòng, Zhōngguó bùdàn zhòngshì huánjìng guǎnlǐ, érqiě fāzhǎn le qīnghāosù yàowù.}\\
« Bien que les situations des deux pays soient différentes, elles peuvent apprendre l'une de l'autre. Le Cameroun diffuse les moustiquaires et la vaccination ; la Chine, non seulement accorde de l'importance à la gestion de l'environnement, mais a aussi développé les médicaments à base d'artémisinine. »\\
我认为两国应该加强南南合作，共同研发更有效的疫苗。总之，健康是每个人共同的权利。\\
\emph{Wǒ rènwéi liǎng guó yīnggāi jiāqiáng Nán-Nán hézuò, gòngtóng yánfā gèng yǒuxiào de yìmiáo. Zǒngzhī, jiànkāng shì měi ge rén gòngtóng de quánlì.}\\
« Je pense que les deux pays devraient renforcer la coopération Sud-Sud et développer ensemble des vaccins plus efficaces. En somme, la santé est un droit commun à chacun. »
\end{textebox}

\textbf{Observation guidée.} Relevez dans ce modèle : (1) deux faits en introduction ; (2) deux points de comparaison reliés par \zh{虽然……但是……} et \zh{不但……而且……} ; (3) un avis introduit par \zh{我认为} avec la modalité \zh{应该} ; (4) une conclusion brève avec \zh{总之}. C'est exactement le squelette attendu.

\begin{popfigure}
\begin{tikzpicture}[node distance=0.35cm]
\node[draw=popblue, thick, rounded corners, fill=popblueL, text width=9cm, align=center] (t1) {\textbf{Titre} : 我对抗疟斗争的看法};
\node[draw=popblue, thick, rounded corners, fill=popblueL, text width=9cm, align=center, below=of t1] (t2) {\textbf{Introduction} : 事实引入 (喀麦隆疟疾高发 / 中国获得消除认证)};
\node[draw=popgreen, thick, rounded corners, fill=popgreenL, text width=9cm, align=center, below=of t2] (t3) {\textbf{Corps 1} : 比较 (预防手段：蚊帐 vs 环境管理)};
\node[draw=popgreen, thick, rounded corners, fill=popgreenL, text width=9cm, align=center, below=of t3] (t4) {\textbf{Corps 2} : 比较 (青蒿素的发现 vs 本土草药研究)};
\node[draw=poporange, thick, rounded corners, fill=poporangeL, text width=9cm, align=center, below=of t4] (t5) {\textbf{Conclusion} : 我认为……应该加强南南合作，共同研发疫苗};
\draw[-{Stealth}, thick, popdark] (t1) -- (t2);
\draw[-{Stealth}, thick, popdark] (t2) -- (t3);
\draw[-{Stealth}, thick, popdark] (t3) -- (t4);
\draw[-{Stealth}, thick, popdark] (t4) -- (t5);
\end{tikzpicture}
\legende{Schéma de la structure imposée de l'article de blog}
\end{popfigure}

\begin{methode}[Rédiger en trois étapes]
\begin{enumerate}
\item \cle{Plan détaillé en français} (5 min) : notez vos deux faits, vos deux points de comparaison et votre proposition. Ne rédigez pas encore en chinois.
\item \cle{Rédaction en chinois} (15 min) : phrases courtes, une idée par phrase, connecteurs imposés (\zh{虽然……但是……}, \zh{不但……而且……}, \zh{由于……因此……}).
\item \cle{Relecture croisée avec la grille} (5 min) : échangez votre copie avec un voisin ; utilisez le dictionnaire pour vérifier l'ordre des traits des 10 caractères cibles.
\end{enumerate}
\end{methode}

\courssub{Grille d'évaluation détaillée (total /20)}

\begin{center}
\begin{tabularx}{\linewidth}{|X|X|X|}\hline
\rowcolor{popblueL} \textbf{Critère} & \textbf{Attendus} & \textbf{Points} \\ \hline
Contenu & Faits exacts (certification chinoise 2021, RTS,S au Cameroun), comparaison équilibrée des deux pays & /5 \\ \hline
Vocabulaire & Au moins 10 mots du glossaire, dont \zh{青蒿素}, \zh{疫苗接种}, \zh{社区防疫}, \zh{耐药性} & /5 \\ \hline
Grammaire & \zh{比}, \zh{由于……因此……}, \zh{不但……而且……}, \zh{虽然……但是……} sans erreur & /5 \\ \hline
Caractères & 10 caractères cibles écrits correctement (ordre des traits, radicaux) & /3 \\ \hline
Cohérence & Plan respecté, connecteurs logiques, conclusion liée à l'introduction & /2 \\ \hline
\end{tabularx}
\end{center}

\begin{definition}[Grille d'évaluation (extrait commenté)]
\cle{Vocabulaire 5/5} : utilise plus de 10 mots du glossaire dont \zh{青蒿素} (\emph{qīnghāosù}, artémisinine), \zh{疫苗接种} (\emph{yìmiáo jiēzhòng}, vaccination), \zh{社区防疫} (\emph{shèqū fángyì}, prévention sanitaire communautaire), \zh{耐药性} (\emph{nàiyàoxìng}, résistance aux médicaments). \cle{Grammaire 5/5} : maîtrise \zh{比}, \zh{由于……因此……}, \zh{不但……而且……}, \zh{虽然……但是……} sans erreur. \cle{Caractères 3/3} : les 10 caractères cibles (\zh{疟}, \zh{疾}, \zh{疫}, \zh{苗}, \zh{蚊}, \zh{帐}, \zh{药}, \zh{防}, \zh{健}, \zh{康}) sont écrits correctement, ordre des traits et radicaux respectés.
\end{definition}

\begin{attention}[Piège de traduction : « Le Cameroun a beaucoup de paludisme »]
La traduction mot à mot \zh{喀麦隆有很多疟疾} (\emph{Kāmàilóng yǒu hěn duō nüèjí}) est compréhensible à l'oral mais maladroite à l'écrit : le verbe \zh{有} est trop vague pour un registre journalistique. Préférez des verbes d'état ou de statut : \zh{喀麦隆疟疾高发} (\emph{Kāmàilóng nüèjí gāofā}, « le paludisme est très fréquent au Cameroun ») ou \zh{喀麦隆是疟疾高负担国家} (\emph{Kāmàilóng shì nüèjí gāo fùdān guójiā}, « le Cameroun est un pays à forte charge palustre »). Règle générale : en chinois écrit, évitez \zh{有} + nom abstrait ; cherchez le verbe précis (\zh{高发}, \zh{流行}, \zh{蔓延}).
\end{attention}

\courssub{Exercice de version (français $\rightarrow$ chinois)}

\begin{exoresolu}[Version : la vaccination RTS,S au Cameroun]
Énoncé. Traduisez en chinois : « Depuis janvier 2024, le Cameroun intègre le vaccin RTS,S dans son programme de vaccination de routine. Ce vaccin réduit les cas graves de paludisme chez les jeunes enfants. C'est une nouvelle espérance pour la prévention. » Vocabulaire fourni : \zh{常规接种} (vaccination de routine), \zh{减少} (réduire), \zh{重症} (cas graves), \zh{希望} (espérance).
\tcblower
Solution détaillée. \zh{自2024年1月起，喀麦隆把RTS,S疫苗纳入常规接种计划。这种疫苗可以减少幼儿的重症疟疾病例。这是预防工作的新希望。} (\emph{Zì 2024 nián 1 yuè qǐ, Kāmàilóng bǎ RTS,S yìmiáo nàrù chángguī jiēzhòng jìhuà. Zhè zhǒng yìmiáo kěyǐ jiǎnshǎo yòu'ér de zhòngzhèng nüèjí bìnglì. Zhè shì yùfáng gōngzuò de xīn xīwàng.}) Points de méthode : (1) « depuis » se rend par \zh{自……起} placé en tête de phrase ; (2) « intégrer » se traduit élégamment par la construction en \zh{把} : \zh{把……纳入……} ; (3) « les cas graves chez les jeunes enfants » devient un complément du nom : \zh{幼儿的重症疟疾病例} (le nom \zh{病例}, « cas », rend la phrase pleinement naturelle) ; (4) « c'est » en début de phrase = \zh{这是}, sans sujet répété.
\end{exoresolu}

\courssub{Exercice de thème (chinois $\rightarrow$ français)}

\begin{textebox}[Texte de thème : l'expérience chinoise]
经过七十年的努力，中国于2021年获得了世界卫生组织的无疟疾认证。中国的成功经验包括三个方面：第一，政府长期重视公共卫生；第二，社区防疫工作深入到每个村庄；第三，科研人员发现了青蒿素，为全球抗疟事业作出了巨大贡献。由于疟疾曾经在中国南方广泛流行，因此这一成就尤其令人鼓舞。中国的经验表明，消除疟疾不但需要药物，而且需要环境管理和全民参与。\\
\emph{Jīngguò qīshí nián de nǔlì, Zhōngguó yú 2021 nián huòdé le Shìjiè Wèishēng Zǔzhī de wú nüèjí rènzhèng. Zhōngguó de chénggōng jīngyàn bāokuò sān ge fāngmiàn : dì-yī, zhèngfǔ chángqī zhòngshì gōnggòng wèishēng ; dì-èr, shèqū fángyì gōngzuò shēnrù dào měi ge cūnzhuāng ; dì-sān, kēyán rényuán fāxiàn le qīnghāosù, wèi quánqiú kàngnüè shìyè zuòchū le jùdà gòngxiàn. Yóuyú nüèjí céngjīng zài Zhōngguó nánfāng guǎngfàn liúxíng, yīncǐ zhè yī chéngjiù yóuqí lìng rén gǔwǔ. Zhōngguó de jīngyàn biǎomíng, xiāochú nüèjí bùdàn xūyào yàowù, érqiě xūyào huánjìng guǎnlǐ hé quánmín cānyù.}
\end{textebox}

\begin{center}
\begin{tabularx}{\linewidth}{|X|X|X|X|}\hline
\rowcolor{popblueL} Caractères & Pinyin & Nature & Traduction \\ \hline
认证 & rènzhèng & n. & certification \\ \hline
经验 & jīngyàn & n. & expérience \\ \hline
社区防疫 & shèqū fángyì & loc. n. & prévention sanitaire communautaire \\ \hline
村庄 & cūnzhuāng & n. & village \\ \hline
青蒿素 & qīnghāosù & n. & artémisinine \\ \hline
贡献 & gòngxiàn & n. & contribution \\ \hline
流行 & liúxíng & v. & sévir, être répandu \\ \hline
消除 & xiāochú & v. & éliminer, éradiquer \\ \hline
全民参与 & quánmín cānyù & loc. & participation de toute la population \\ \hline
\end{tabularx}
\end{center}

\begin{exercice}[Thème noté sur 5]
Traduisez le texte ci-dessus en français. Barème indicatif : 1 point par segment correctement rendu (5 segments), moins 0,25 par contresens sur un terme du glossaire.
\end{exercice}

\begin{corrige}[Exercice de thème]
« Après soixante-dix ans d'efforts, la Chine a obtenu en 2021 la certification "sans paludisme" de l'Organisation mondiale de la santé. L'expérience réussie de la Chine comprend trois aspects : premièrement, le gouvernement a accordé une attention durable à la santé publique ; deuxièmement, le travail de prévention communautaire s'est étendu jusqu'à chaque village ; troisièmement, les chercheurs ont découvert l'artémisinine, apportant une immense contribution à la lutte antipaludique mondiale. Comme le paludisme sévissait autrefois largement dans le sud de la Chine, cette réussite est particulièrement encourageante. L'expérience chinoise montre qu'éliminer le paludisme exige non seulement des médicaments, mais aussi la gestion de l'environnement et la participation de toute la population. »
\end{corrige}

\courssub{Auto-correction guidée et réécriture}

\begin{methode}[Le code de correction V-G-C-S]
Corrigez la copie de votre voisin en annotant dans la marge : \cle{V} = erreur de vocabulaire (mot impropre ou absent du glossaire) ; \cle{G} = erreur de grammaire (structure comparative, connecteur, ordre des mots) ; \cle{C} = caractère mal écrit (trait manquant, radical confondu) ; \cle{S} = problème de sens (phrase incompréhensible ou hors sujet). Ensuite, réécrivez intégralement votre article en tenant compte des annotations : la version finale est celle qui sera notée sur 20.
\end{methode}

\begin{attention}[Erreurs fréquentes à repérer lors de la relecture]
\begin{itemize}
\item \zh{比} suivi d'un adjectif sans degré : on dit \zh{中国的人口比喀麦隆多} (\emph{Zhōngguó de rénkǒu bǐ Kāmàilóng duō}), jamais \zh{比喀麦隆很多} ; pour marquer l'écart, ajoutez le complément de degré après l'adjectif : \zh{多得多} (« beaucoup plus »).
\item Connecteurs coupés : \zh{虽然……} exige \zh{但是……} dans la seconde proposition ; ne pas oublier la seconde moitié.
\item Confusion de caractères proches : \zh{疟} (paludisme) et \zh{疾} (maladie) partagent le radical de la maladie \zh{疒} ; \zh{苗} (vaccin, pousse) s'écrit avec la clé de l'herbe \zh{艹} au-dessus de \zh{田}.
\item Traduction de « je pense que » : \zh{我认为} + proposition complète, sans \zh{是} superflu devant un verbe.
\end{itemize}
\end{attention}

\courssub{Prolongement « Bac » : analyse de document}

\begin{exercice}[Sujet type baccalauréat]
Document : le texte de thème ci-dessus, accompagné d'un graphique montrant la couverture vaccinale RTS,S qui progresse dans les régions du Centre et du Littoral entre 2024 et 2025. Consigne : en 100 à 120 caractères, présentez le document, dégagez sa thèse principale et dites, en vous appuyant sur un fait camerounais et un fait chinois, si la coopération Sud-Sud peut accélérer l'élimination du paludisme. Utilisez au moins deux connecteurs étudiés.
\end{exercice}

\begin{savaistu}[Le vaccin RTS,S/AS01]
Le vaccin \zh{RTS,S/AS01} (Mosquirix) est le premier vaccin antipaludique recommandé par l'Organisation mondiale de la santé (2021). Le Cameroun est pionnier dans son déploiement en routine (janvier 2024), une première mondiale à cette échelle. De son côté, la Chine poursuit ses propres programmes de recherche vaccinale et reste le pays de l'artémisinine : cette molécule, extraite de l'armoise annuelle \zh{青蒿} (\emph{qīnghāo}) selon les recettes de la médecine traditionnelle, a valu à la chercheuse Tu Youyou \zh{屠呦呦} le prix Nobel de médecine en 2015.
\end{savaistu}

\begin{aretenir}[Bilan de la séquence]
Un article de blog argumenté en chinois repose sur quatre piliers : des faits exacts introduits par des verbes d'état précis (\zh{高发}, \zh{获得认证}), une comparaison équilibrée avec \zh{虽然……但是……} et \zh{不但……而且……}, un avis personnel avec \zh{我认为……应该……}, et une conclusion brève avec \zh{总之}. La grille sur 20 points (contenu 5, vocabulaire 5, grammaire 5, caractères 3, cohérence 2) est votre feuille de route : relisez-vous toujours critère par critère avant de rendre votre copie.
\end{aretenir}

\newpage

\coursec{Activité d'intégration}

\courssub{Objectifs de l'activité}

À l'issue de cette activité d'intégration, tu seras capable de :

\begin{itemize}
\item mobiliser le \cle{vocabulaire} des maladies endémiques (paludisme, dengue, choléra, tuberculose), des symptômes, de la prévention et du traitement, en caractères et en pinyin ;
\item comparer de façon factuelle et respectueuse la situation sanitaire du \cle{Cameroun} et de la \cle{Chine} à l'aide des structures de comparaison (\zh{比}, \zh{和……一样}, \zh{越来越}, \zh{虽然……但是……}) ;
\item argumenter et nuancer un point de vue simple en chinois, à l'écrit comme à l'oral ;
\item produire un document bilingue (chinois/français) de type infographie, support d'une présentation orale de trois minutes ;
\item travailler en binôme selon une démarche de projet : rechercher, sélectionner, rédiger, présenter, s'autoévaluer.
\end{itemize}

\courssub{Situation}

Dans le cadre de la \emph{Semaine Santé Jeunesse} organisée par ton lycée de Yaoundé, en partenariat avec un club de chinois d'un établissement de Douala et une classe partenaire de Beijing, le proviseur lance un concours : chaque classe de Terminale doit produire des affiches bilingues chinois-français de sensibilisation sur les maladies endémiques. Les meilleures productions seront exposées au hall du lycée, photographiées et envoyées à la classe partenaire chinoise, qui réalisera de son côté des affiches sur les mêmes maladies vues de Chine.

Le club santé du lycée a publié l'appel à projets suivant :

\begin{textebox}[Appel à projets du club santé — 健康周海报比赛]
同学们好！\\
Tóngxuemen hǎo !\\
Chers camarades, bonjour !\\[4pt]
为了“青年健康周”，我们学校组织一个中法双语海报比赛。\\
Wèile ``Qīngnián Jiànkāng Zhōu'', wǒmen xuéxiào zǔzhī yí ge Zhōng-Fǎ shuāngyǔ hǎibào bǐsài.\\
Pour la « Semaine Santé Jeunesse », notre école organise un concours d'affiches bilingues chinois-français.\\[4pt]
主题：非洲和亚洲的地方病（疟疾、登革热、霍乱、结核病）。\\
Zhǔtí : Fēizhōu hé Yàzhōu de dìfāngbìng (nüèji, dēnggérè, huòluàn, jiéhébìng).\\
Thème : les maladies endémiques d'Afrique et d'Asie (paludisme, dengue, choléra, tuberculose).\\[4pt]
每组两个人，选一种病，做一张A3海报，然后向全班介绍三分钟。\\
Měi zǔ liǎng ge rén, xuǎn yì zhǒng bìng, zuò yì zhāng A3 hǎibào, ránhòu xiàng quán bān jièshào sān fēnzhōng.\\
Chaque groupe de deux personnes choisit une maladie, réalise une affiche A3, puis la présente à la classe en trois minutes.\\[4pt]
海报上要有：地图、比较表、重要数字、一个“你知道吗？”小栏目，还有一句你自己的承诺。\\
Hǎibào shang yào yǒu : dìtú, bǐjiào biǎo, zhòngyào shùzì, yí ge ``Nǐ zhīdào ma ?'' xiǎo lánmù, háiyǒu yí jù nǐ zìjǐ de chéngnuò.\\
L'affiche doit contenir : une carte, un tableau comparatif, des chiffres importants, une rubrique « Le savais-tu ? », et une phrase d'engagement personnel.\\[4pt]
最好的海报会寄给我们北京的伙伴班！\\
Zuì hǎo de hǎibào huì jì gěi wǒmen Běijīng de huǒbàn bān !\\
Les meilleures affiches seront envoyées à notre classe partenaire de Beijing !
\end{textebox}

\begin{methode}[Comprendre l'appel à projets en quatre questions]
\begin{enumerate}
\item \textbf{Qui ?} — \zh{每组两个人} : le travail se fait en binôme.
\item \textbf{Quoi ?} — \zh{选一种病，做一张A3海报} : une maladie au choix, une affiche au format A3.
\item \textbf{Comment ?} — \zh{海报上要有……} : cinq éléments obligatoires sont énumérés (repère la virgule d'énumération \zh{、}).
\item \textbf{Pourquoi ?} — \zh{寄给我们北京的伙伴班} : la production a un destinataire réel, la classe partenaire ; ton chinois doit donc être soigné.
\end{enumerate}
\end{methode}

\courssub{Consignes}

En binôme, réalisez le projet en suivant les étapes ci-dessous.

\textbf{Étape 1 — Choisir et comprendre (compréhension).} Choisissez une maladie parmi : \zh{疟疾} (nüèji, paludisme), \zh{登革热} (dēnggérè, dengue), \zh{霍乱} (huòluàn, choléra), \zh{结核病} (jiéhébìng, tuberculose). Relisez le texte support de la leçon et l'appel à projets ci-dessus, puis notez en français les informations essentielles sur votre maladie : agent pathogène, mode de transmission, symptômes, prévention, traitement.

\textbf{Étape 2 — Construire le lexique (vocabulaire).} Complétez un mini-glossaire bilingue d'au moins \textbf{dix mots} utiles à votre maladie, selon le modèle : Caractères | Pinyin | Nature | Traduction. Vous pouvez partir des mots suivants et en ajouter d'autres.

\begin{center}
\begin{tabularx}{\linewidth}{|X|X|X|X|}
\hline
\rowcolor{popblueL} Caractères & Pinyin & Nature & Traduction \\
\hline
地方病 & dìfāngbìng & n. & maladie endémique \\
\hline
蚊子 & wénzi & n. & moustique \\
\hline
传染 & chuánrǎn & v. & transmettre, contaminer \\
\hline
发烧 & fāshāo & v. & avoir de la fièvre \\
\hline
咳嗽 & késou & v. & tousser \\
\hline
头疼 & tóuténg & v. & avoir mal à la tête \\
\hline
预防 & yùfáng & v. & prévenir \\
\hline
蚊帐 & wénzhàng & n. & moustiquaire \\
\hline
疫苗 & yìmiáo & n. & vaccin \\
\hline
治疗 & zhìliáo & v./n. & traiter ; traitement \\
\hline
卫生 & wèishēng & n./adj. & hygiène ; hygiénique \\
\hline
病人 & bìngrén & n. & malade, patient \\
\hline
\end{tabularx}
\end{center}

\textbf{Étape 3 — Comparer (grammaire).} Rédigez \textbf{quatre phrases comparatives Cameroun/Chine} en chinois, en utilisant au moins trois structures différentes : \zh{比}, \zh{和……一样}, \zh{越来越}, \zh{虽然……但是……}. Exemple de départ : \zh{在喀麦隆，疟疾比登革热更常见。} (\emph{Zài Kāmàilóng, nüèji bǐ dēnggérè gèng chángjiàn.} « Au Cameroun, le paludisme est plus fréquent que la dengue. »)

\textbf{Étape 4 — Concevoir l'infographie (production écrite).} Sur une feuille A3 ou un support numérique, organisez les cinq éléments obligatoires : (1) une carte de répartition simplifiée (Cameroun, Chine, monde) ; (2) un tableau comparatif à cinq lignes : agent pathogène, transmission, symptômes clés (en caractères), prévention, traitement (moderne et traditionnel) ; (3) deux ou trois chiffres clés en précisant la source et l'année (OMS, MINSANTE, NHC) — si vous n'avez pas le chiffre exact, écrivez une formulation prudente comme \zh{每年有几百万人得病} (\emph{měi nián yǒu jǐ bǎi wàn rén dé bìng}, « chaque année, plusieurs millions de personnes tombent malades ») ; (4) un encadré « 你知道吗？ » (\emph{Nǐ zhīdào ma ?}, « Le savais-tu ? ») avec une anecdote culturelle ou historique ; (5) une phrase d'engagement personnel commençant par \zh{我要……} (\emph{Wǒ yào...}, « Je veux / je vais... »).

\textbf{Étape 5 — Présenter (production orale).} Préparez une présentation orale de trois minutes : chaque membre du binôme parle environ une minute trente. Structure conseillée : salutation (\zh{大家好！}), présentation de la maladie, comparaison Cameroun/Chine, lecture de l'encadré « Le savais-tu ? », engagement personnel, remerciement (\zh{谢谢大家！}).

\textbf{Étape 6 — Autoévaluer.} Avant de rendre l'affiche, vérifiez avec la grille : caractères corrects, pinyin avec tons, au moins trois structures de comparaison, cinq éléments présents, français sans faute, oral chronométré.

\courssub{Ressources à mobiliser}

\begin{aretenir}[Boîte à outils de la leçon]
\begin{itemize}
\item \textbf{Comparaison} : A + \zh{比} + B + adjectif → \zh{霍乱比结核病传染得更快。} (\emph{Huòluàn bǐ jiéhébìng chuánrǎn de gèng kuài.} « Le choléra se transmet plus vite que la tuberculose. ») — ici \zh{得} introduit le complément de degré après le verbe \zh{传染}.
\item \textbf{Égalité} : A + \zh{和} + B + \zh{一样} (+ adjectif) → \zh{预防和治疗一样重要。} (\emph{Yùfáng hé zhìliáo yíyàng zhòngyào.} « Prévention et traitement sont aussi importants l'un que l'autre. »)
\item \textbf{Évolution} : \zh{越来越} + adjectif → \zh{用蚊帐睡觉的人越来越多。} (\emph{Yòng wénzhàng shuìjiào de rén yuè lái yuè duō.} « De plus en plus de gens dorment sous moustiquaire. »)
\item \textbf{Concession} : \zh{虽然……但是……} → \zh{虽然疟疾很危险，但是可以预防。} (\emph{Suīrán nüèji hěn wēixiǎn, dànshì kěyǐ yùfáng.} « Bien que le paludisme soit dangereux, on peut le prévenir. »)
\item \textbf{Engagement} : \zh{我要} + verbe → \zh{我要宣传用蚊帐睡觉。} (\emph{Wǒ yào xuānchuán yòng wénzhàng shuìjiào.} « Je veux promouvoir l'usage de la moustiquaire. »)
\item \textbf{Lexique} : maladies, symptômes (\zh{发烧}, \zh{咳嗽}, \zh{头疼} \emph{tóuténg}, maux de tête), prévention (\zh{洗手} \emph{xǐ shǒu}, se laver les mains ; \zh{打疫苗} \emph{dǎ yìmiáo}, se faire vacciner).
\end{itemize}
\end{aretenir}

\begin{attention}[Pièges fréquents des francophones]
\begin{itemize}
\item Ne traduis pas « plus... que » mot à mot : utilise la structure A + \zh{比} + B + adjectif, éventuellement renforcée par \zh{更} (\emph{gèng}). Ne combine jamais \zh{比} avec \zh{一样} dans la même phrase.
\item Après \zh{虽然}, n'oublie pas \zh{但是} dans la seconde partie : le chinois garde les deux termes, contrairement au français où « bien que... mais » serait fautif.
\item Attention aux tons : \zh{wénzi} (蚊子, moustique) commence par un deuxième ton montant ; ne le confonds pas avec \zh{wèn} (问, demander), quatrième ton descendant. De même, \zh{yùfáng} (预防, prévenir) porte deux tons : quatrième puis deuxième.
\item Ordre des mots : le complément de lieu introduit par \zh{在} se place avant le verbe : \zh{在喀麦隆，……} et non après le verbe comme en français.
\item On dit \zh{用蚊帐睡觉} (« dormir en utilisant une moustiquaire ») : le verbe \zh{用} (utiliser) précède l'action principale ; évite la formulation télégraphique \zh{睡蚊帐}, incorrecte en chinois soigné.
\end{itemize}
\end{attention}

\courssub{Proposition de production et corrigé commenté}

\begin{exoresolu}[Modèle d'infographie orale : le paludisme]
Énoncé : rédiger le script oral de trois minutes accompagnant une infographie sur le paludisme. \tcblower
\textbf{Production modèle (script oral) :}\\[4pt]
大家好！我们是玛丽和保罗。今天我们要介绍疟疾。\\
\emph{Dàjiā hǎo ! Wǒmen shì Mǎlì hé Bǎoluó. Jīntiān wǒmen yào jièshào nüèji.}\\
« Bonjour à tous ! Nous sommes Marie et Paul. Aujourd'hui, nous allons présenter le paludisme. »\\[4pt]
疟疾是一种地方病，蚊子传染这种病。在喀麦隆，疟疾比登革热更常见；在中国，疟疾越来越少了。\\
\emph{Nüèji shì yì zhǒng dìfāngbìng, wénzi chuánrǎn zhè zhǒng bìng. Zài Kāmàilóng, nüèji bǐ dēnggérè gèng chángjiàn ; zài Zhōngguó, nüèji yuè lái yuè shǎo le.}\\
« Le paludisme est une maladie endémique, transmise par le moustique. Au Cameroun, il est plus fréquent que la dengue ; en Chine, il devient de plus en plus rare. »\\[4pt]
病人常常发烧、头疼。虽然疟疾很危险，但是可以预防：我们要用蚊帐睡觉，也要看医生。\\
\emph{Bìngrén chángcháng fāshāo, tóuténg. Suīrán nüèji hěn wēixiǎn, dànshì kěyǐ yùfáng : wǒmen yào yòng wénzhàng shuìjiào, yě yào kàn yīshēng.}\\
« Les malades ont souvent de la fièvre et des maux de tête. Bien que le paludisme soit dangereux, on peut le prévenir : il faut dormir sous moustiquaire et consulter un médecin. »\\[4pt]
你知道吗？中国科学家屠呦呦因为研究青蒿素，得到了诺贝尔奖！\\
\emph{Nǐ zhīdào ma ? Zhōngguó kēxuéjiā Tú Yōuyōu yīnwèi yánjiū qīnghāosù, dédào le Nuòbèi'ěr jiǎng !}\\
« Le savais-tu ? La scientifique chinoise Tu Youyou a reçu le prix Nobel pour ses recherches sur l'artémisinine ! »\\[4pt]
我要宣传用蚊帐睡觉。谢谢大家！\\
\emph{Wǒ yào xuānchuán yòng wénzhàng shuìjiào. Xièxie dàjiā !}\\
« Je veux promouvoir le sommeil sous moustiquaire. Merci à tous ! »\\[6pt]
\textbf{Commentaire des choix de langue :}
\begin{itemize}
\item La structure \zh{比} est utilisée avec l'adverbe \zh{更} pour renforcer la comparaison, ce qui est naturel en chinois.
\item \zh{越来越……了} exprime l'évolution : le \zh{了} final marque le changement d'état, point délicat pour les francophones.
\item La concession \zh{虽然……但是……} est complète, avec les deux éléments, comme l'exige le chinois.
\item L'énumération des symptômes utilise la virgule chinoise \zh{、}, réservée aux énumérations de mots coordonnés.
\item L'engagement \zh{我要} + verbe est suivi d'une action concrète, conforme à la consigne ; \zh{宣传} (\emph{xuānchuán}, faire de la sensibilisation) est un mot de niveau HSK 4, valorisant en Terminale.
\item Le point-virgule chinois \zh{；} relie deux constats parallèles (Cameroun / Chine), ponctuation élégante à réutiliser dans tes productions.
\end{itemize}
\end{exoresolu}

\begin{exercice}[À toi de jouer : la tuberculose]
Rédige, sur le modèle ci-dessus, un court script oral (six à huit phrases) sur la tuberculose (\zh{结核病}). Contraintes : une phrase avec \zh{比}, une avec \zh{虽然……但是……}, un encadré « 你知道吗？ » et un engagement personnel avec \zh{我要}.
\end{exercice}

\begin{corrige}[Exercice : la tuberculose — éléments de réponse]
大家好！今天我们要介绍结核病。结核病是一种地方病，病人常常咳嗽、发烧。在喀麦隆，结核病比霍乱更常见。虽然结核病很危险，但是可以治疗：病人要看医生、吃药。你知道吗？结核病通过空气传染，所以卫生很重要！我要宣传洗手、开窗。谢谢大家！\\
\emph{Dàjiā hǎo ! Jīntiān wǒmen yào jièshào jiéhébìng. Jiéhébìng shì yì zhǒng dìfāngbìng, bìngrén chángcháng késou, fāshāo. Zài Kāmàilóng, jiéhébìng bǐ huòluàn gèng chángjiàn. Suīrán jiéhébìng hěn wēixiǎn, dànshì kěyǐ zhìliáo : bìngrén yào kàn yīshēng, chī yào. Nǐ zhīdào ma ? Jiéhébìng tōngguò kōngqì chuánrǎn, suǒyǐ wèishēng hěn zhòngyào ! Wǒ yào xuānchuán xǐ shǒu, kāi chuāng. Xièxie dàjiā !}\\
« Bonjour à tous ! Aujourd'hui, nous allons présenter la tuberculose. La tuberculose est une maladie endémique ; les malades toussent souvent et ont de la fièvre. Au Cameroun, la tuberculose est plus fréquente que le choléra. Bien que la tuberculose soit dangereuse, on peut la traiter : les malades doivent consulter un médecin et prendre des médicaments. Le savais-tu ? La tuberculose se transmet par l'air, donc l'hygiène est très importante ! Je veux promouvoir le lavage des mains et l'aération des pièces. Merci à tous ! »
\end{corrige}

\courssub{Bilan de l'activité}

Cette activité d'intégration t'a permis de réunir toutes les compétences de la leçon dans une tâche authentique et créative : comprendre un appel à projets en chinois, sélectionner et organiser un lexique thématique, comparer deux réalités sanitaires avec des structures grammaticales précises, rédiger un document bilingue et le présenter à l'oral. Tu as aussi appris à formuler des chiffres avec prudence quand la source exacte manque, compétence essentielle de l'expression responsable. Enfin, en préparant un message destiné à une classe partenaire chinoise, tu as pratiqué le chinois comme une langue vivante de communication internationale, au service de la santé publique et du dialogue entre le Cameroun et la Chine.

\newpage

\coursec{Méthodes et exercices résolus}

\courssub{Entrée en matière \& Lexique thématique : Maladies, Symptômes, Prévention}

\begin{textebox}[Situation de communication : Au centre de santé de Biyem-Assi]
\zh{护士：您好！有什么不舒服？} (Hùshì: Nín hǎo! Yǒu shénme bù shūfu?) « Infirmière : Bonjour ! Qu'est-ce qui ne va pas ? » \\
\zh{病人：我发烧、头疼，还有点咳嗽。} (Bìngrén: Wǒ fāshāo, tóuténg, hái yǒudiǎn késou.) « Patient : J'ai de la fièvre, mal à la tête et un peu de toux. » \\
\zh{护士：雨季到了，得疟疾的人多。您用蚊帐了吗？} (Hùshì: Yǔjì dào le, dé nüèjí de rén duō. Nín yòng wénzhàng le ma?) « Infirmière : La saison des pluies est là, beaucoup attrapent le paludisme. Vous utilisez une moustiquaire ? » \\
\zh{病人：用了，但我还没打疫苗。} (Bìngrén: Yòng le, dàn wǒ hái méi dǎ yìmiáo.) « Patient : Oui, mais je ne me suis pas encore fait vacciner. » \\
\zh{护士：预防很重要。多喝干净的水，注意卫生。} (Hùshì: Yùfáng hěn zhòngyào. Duō hē gānjìng de shuǐ, zhùyì wèishēng.) « Infirmière : La prévention est cruciale. Buvez de l'eau propre, soignez l'hygiène. »
\end{textebox}

\begin{definition}[Lexique thématique classé par familles (HSK 3-4)]
\begin{center}
\begin{tabularx}{\linewidth}{|X|X|X|X|}
\hline
\rowcolor{popblueL} \textbf{Caractères} & \textbf{Pinyin} & \textbf{Nature} & \textbf{Traduction} \\ \hline
\zh{疟疾} & nüèjí & nom & paludisme (maladie endémique majeure) \\ \hline
\zh{霍乱} & huòluàn & nom & choléra \\ \hline
\zh{登革热} & dēnggérè & nom & dengue \\ \hline
\zh{结核病} & jiéhébìng & nom & tuberculose \\ \hline
\zh{发烧} & fāshāo & verbe & avoir de la fièvre \\ \hline
\zh{头疼} & tóuténg & verbe & avoir mal à la tête \\ \hline
\zh{咳嗽} & késou & verbe & tousser \\ \hline
\zh{腹泻} & fùxiè & verbe/nom & diarrhée \\ \hline
\zh{预防} & yùfáng & verbe & prévenir \\ \hline
\zh{治疗} & zhìliáo & verbe & soigner, traiter \\ \hline
\zh{接种疫苗} & jiēzhòng yìmiáo & loc. verbale & se faire vacciner (formel) \\ \hline
\zh{打疫苗} & dǎ yìmiáo & loc. verbale & se faire vacciner (courant) \\ \hline
\zh{蚊帐} & wénzhàng & nom & moustiquaire \\ \hline
\zh{卫生} & wèishēng & nom & hygiène, salubrité \\ \hline
\zh{干净} & gānjìng & adj. & propre \\ \hline
\zh{免费} & miǎnfèi & adj. & gratuit \\ \hline
\zh{公共卫生} & gōnggòng wèishēng & nom & santé publique \\ \hline
\zh{雨季} & yǔjì & nom & saison des pluies \\ \hline
\zh{传播} & chuánbō & verbe & transmettre, propager \\ \hline
\end{tabularx}
\end{center}
\end{definition}

\begin{propriete}[Clés de caractères (radicaux) et aide à la mémorisation]
\begin{itemize}
\item \cle{Clé \zh{疒} (bìng, « malade »)} : présente dans \zh{疟} (nüè, paludisme), \zh{疾} (jí, maladie grave), \zh{疼} (téng, douleur), \zh{痢} (lì, dysenterie), \zh{癌} (ái, cancer). Elle signale le domaine médical/pathologique.
\item \cle{Clé \zh{氵} (shuǐ, « eau »)} : dans \zh{治} (zhì, soigner — à l'origine « gérer l'eau »), \zh{洗} (xǐ, laver), \zh{霍} (huò, choléra — agitation/soudain), \zh{溢} (yì, déborder). Lien avec l'hygiène (\zh{洗手}) et l'épidémiologie (eau contaminée).
\item \cle{Clé \zh{虫} (chóng, « insecte »)} : dans \zh{蚊} (wén, moustique), \zh{蚤} (zǎo, puce). Rappel : le paludisme et la dengue sont des \zh{虫媒疾病} (chóngméi jíbìng, maladies à transmission vectorielle).
\item \cle{Classificateurs essentiels} : \zh{一种病} (yī zhǒng bìng, une maladie), \zh{一顶蚊帐} (yī dǐng wénzhàng, une moustiquaire), \zh{一位医生} (yī wèi yīshēng, un médecin — respect), \zh{两个国家} (liǎng gè guójiā, deux pays — \zh{两} obligatoire devant classif., jamais \zh{二个}).
\end{itemize}
\end{propriete}

\courssub{Texte support : « Lutte contre les épidémies : regards croisés »}

\begin{textebox}[Article original — Niveau HSK 3/4]
\zh{喀麦隆和中国都面临传染病的挑战，但应对方式有异有同。} (Kāmàilóng hé Zhōngguó dōu miànlín chuánrǎnbìng de tiǎozhàn, dàn yìngduì fāngshì yǒu yì yǒu tóng.) « Le Cameroun et la Chine font tous deux face au défi des maladies infectieuses, mais leurs réponses présentent des différences et des similitudes. » \\
\zh{在喀麦隆，疟疾是雨季最常见的病。蚊子传播疟疾，所以预防的核心是防蚊。政府和非政府组织在学校、医院免费发放蚊帐，宣传“睡蚊帐、喝干净水、勤洗手”。霍乱有时在洪水后爆发，喝干净的水、注意饮食卫生是关键。} (Zài Kāmàilóng, nüèjí shì yǔjì zuì chángjiàn de bìng. Wénzi chuánbō nüèjí, suǒyǐ yùfáng de héxīn shì fáng wén. Zhèngfǔ hé fēi zhèngfǔ zǔzhī zài xuéxiào, yīyuàn miǎnfèi fāfàng wénzhàng, xuānchuán “shuì wénzhàng, hē gānjìng shuǐ, qín xǐ shǒu”. Huòluàn yǒushí zài hóngshuǐ hòu bàofā, hē gānjìng de shuǐ, zhùyì yǐnshí wèishēng shì guānjiàn.) « Au Cameroun, le paludisme est la maladie la plus courante en saison des pluies. Les moustiques transmettent le paludisme, donc le cœur de la prévention est la lutte antivectorielle. Le gouvernement et les ONG distribuent gratuitement des moustiquaires dans les écoles et hôpitaux, et promeuvent "dormir sous moustiquaire, boire de l'eau propre, se laver souvent les mains". Le choléra éclate parfois après les inondations ; boire de l'eau propre et surveiller l'hygiène alimentaire est crucial. » \\
\zh{在中国，政府建立了强大的公共卫生体系。城市里医院多、设备新，疫苗接种覆盖率很高。历史上中国消灭了天花、疟疾（2021年获WHO无疟疾认证），现在重点防控登革热、结核病和新发传染病。中医药在预防和康复中也发挥作用，比如用艾草驱蚊、食疗调理身体。} (Zài Zhōngguó, zhèngfǔ jiànlì le qiángdà de gōnggòng wèishēng tǐxì. Chéngshì lǐ yīyuàn duō, shèbèi xīn, yìmiáo jiēzhòng fùgàilǜ hěn gāo. Lìshǐ shàng Zhōngguó xiāomiè le tiānhuā, nüèjí (2021 nián huò WHO wú nüèjí rènzhèng), xiànzài zhòngdiǎn fángkòng dēnggérè, jiéhébìng hé xīnfā chuánrǎnbìng. Zhōngyīyào zài yùfáng hé kāngfù zhōng yě fāhuī zuòyòng, bǐrú yòng àicǎo qūwén, shíliáo tiáolǐ shēntǐ.) « En Chine, le gouvernement a bâti un système de santé publique puissant. En ville, les hôpitaux sont nombreux et bien équipés, la couverture vaccinale est très élevée. Historiquement, la Chine a éradiqué la variole et le paludisme (certification OMS "sans paludisme" en 2021), et se concentre maintenant sur la dengue, la tuberculose et les maladies émergentes. La médecine traditionnelle joue aussi un rôle en prévention et réhabilitation, ex. : utiliser l'armoise pour chasser les moustiques, la diététique pour réguler le corps. » \\
\zh{两国共识：预防优于治疗，社区参与不可少。} (Liǎng guó gòngshì: Yùfáng yōuyú zhìliáo, shèqū cānyù bùkě shǎo.) « Consensus des deux pays : la prévention vaut mieux que le traitement, la participation communautaire est indispensable. »
\end{textebox}

\courssub{Faits socioculturels comparés : Cameroun vs Chine (Tableau factuel)}

\begin{center}
\begin{tabularx}{\linewidth}{|X|X|X|}
\hline
\rowcolor{popblueL} \textbf{Thème} & \textbf{Cameroun (Contexte local)} & \textbf{Chine (Contexte national)} \\ \hline
\zh{主要流行病} (Principales maladies endémiques) & \zh{疟疾、霍乱、登革热、结核病、艾滋病} (Paludisme, choléra, dengue, tuberculose, VIH/SIDA) & \zh{登革热、结核病、病毒性肝炎、新发传染病} (Dengue, tuberculose, hépatites virales, maladies émergentes) \\ \hline
\zh{疟疾现状} (Statut paludisme) & \zh{高负荷，雨季高发，儿童/孕妇高危} (Forte charge, pic saison des pluies, enfants/femmes enceintes à risque) & \zh{已消除 (2021 WHO认证), 仅有输入性病例} (Éliminé 2021 OMS, seulement cas importés) \\ \hline
\zh{预防核心手段} (Moyens prévention clés) & \zh{长效杀虫蚊帐 (LLIN)、室内滞留喷洒 (IRS)、环境管理、疫苗试点 (RTS,S)} (Moustiquaires imprégnées, pulvérisation intra, gestion environnement, vaccin pilote) & \zh{疫苗接种国家规划、病媒生物监测、爱国卫生运动、中医药预防} (Plan national vaccination, surveillance vecteurs, mouvement patriotique hygiène, MTC prévention) \\ \hline
\zh{卫生体系特点} (Caractéristiques système santé) & \zh{金字塔结构：区医院→区域医院→中央医院；社区卫生工作者关键} (Structure pyramide : hôpital district → régional → central ; agents santé communautaire clés) & \zh{分级诊疗：社区卫生服务中心→三级医院；医保覆盖率高} (Diagnostic progressif : centre santé communautaire → hôpital niveau 3 ; couverture assurance élevée) \\ \hline
\zh{文化特色} (Spécificités culturelles) & \zh{传统草药治疗并存、社区动员靠酋长/教会、足球明星做公益广告} (Médecine traditionnelle coexiste, mobilisation via chefs/églises, stars football spots pub) & \zh{中西医结合、“爱国卫生运动”全民动员、数字化健康码/预约挂号} (Intégration médecine chinoise/occidentale, mobilisation "Mouvement patriotique hygiène", santé numérique) \\ \hline
\end{tabularx}
\end{center}

\begin{savaistu}[Le saviez-vous ? Le « Mouvement patriotique d'hygiène » 爱国卫生运动]
Lancé en 1952 en Chine, ce mouvement de masse vise à améliorer l'environnement sanitaire (élimination des « quatre nuisibles » : rats, mouches, moustiques, cafards — \zh{除四害}). Il a joué un rôle majeur dans l'éradication du paludisme et du choléra. Au Cameroun, les \zh{Journées nationales de l'assainissement} (opérations « Ville propre ») en sont l'équivalent citoyen, souvent portées par les mairies et les chefs traditionnels.
\end{savaistu}

\courssub{Outils linguistiques : Comparer, Argumenter, Nuancer, Traduire}

\begin{methode}[Comparer deux réalités sanitaires en chinois — Structures types]
Pour comparer les pratiques de santé publique au Cameroun et en Chine, suis ces étapes :
\begin{enumerate}
\item \cle{Identifier} les deux éléments comparés (A et B) : ex. \zh{喀麦隆} (Kāmàilóng) et \zh{中国} (Zhōngguó).
\item \cle{Choisir la structure} adaptée au sens visé :
\begin{itemize}
\item \textbf{Supériorité} : A + \zh{比} (bǐ) + B + \textbf{Adjectif (sans 很)} \\
\zh{疟疾比感冒严重。} (Nüèjí bǐ gǎnmào yánzhòng.) « Le paludisme est plus grave que le rhume. » \\
\emph{Renforcement} : A + \zh{比} + B + \zh{更} (gèng) / \zh{还} (hái) + Adj. → \zh{疟疾比感冒更严重。}
\item \textbf{Infériorité} : B + \zh{比} + A + Adj. OU A + \zh{没有} (méiyǒu) + B + Adj. \\
\zh{感冒没有疟疾严重。} (Gǎnmào méiyǒu nüèjí yánzhòng.) « Le rhume n'est pas aussi grave que le paludisme. »
\item \textbf{Égalité} : A + \zh{跟/和/同} (gēn/hé/tóng) + B + \zh{一样} (yīyàng) + Adj./Verbe \\
\zh{中国和喀麦隆一样重视预防。} (Zhōngguó hé Kāmàilóng yīyàng zhòngshì yùfáng.) « La Chine et le Cameroun attachent la même importance à la prévention. »
\item \textbf{Différence de degré (quantifiée)} : A + \zh{比} + B + Adj. + \zh{得多} (de duō) / \zh{多了} (duō le) / \zh{一倍} (yī bèi) \\
\zh{城市里的医院比农村多得多。} (Chéngshì lǐ de yīyuàn bǐ nóngcūn duō de duō.) « Les hôpitaux en ville sont bien plus nombreux qu'à la campagne. »
\end{itemize}
\item \cle{Placer le complément de lieu/temps AVANT le verbe} (règle d'or ordre chinois) : \zh{在喀麦隆，人们用蚊帐预防疟疾。} (Zài Kāmàilóng, rénmen yòng wénzhàng yùfáng nüèjí.)
\item \cle{Conclure} par une synthèse nuancée : \zh{虽然……，但是……} (Suīrán..., dànshì...).
\end{enumerate}
\emph{Réflexe Bac} : Après \zh{比}, \textbf{jamais de 很, 非常, 特别}. Pour intensifier : \zh{更}, \zh{还}, \zh{得多}.
\end{methode}

\begin{exoresolu}[Comparaison guidée : Application immédiate]
Énoncé — Construis trois phrases correctes avec les éléments donnés.
\begin{enumerate}
\item (a) \zh{疟疾} (nüèjí) / \zh{感冒} (gǎnmào) / \zh{严重} (yánzhòng) — Supériorité.
\item (b) \zh{中国} (Zhōngguó) / \zh{喀麦隆} (Kāmàilóng) / \zh{重视接种疫苗} (zhòngshì jiēzhòng yìmiáo) — Égalité.
\item (c) \zh{蚊帐} (wénzhàng) / \zh{药} (yào) / \zh{便宜} (piányi) — « Beaucoup moins cher ».
\end{enumerate}
\tcblower
Solution détaillée :
\begin{enumerate}
\item (a) Structure A \zh{比} B Adj. \zh{疟疾比感冒严重。} (Nüèjí bǐ gǎnmào yánzhòng.) \textbf{Pas de 很.}
\item (b) Structure A \zh{和} B \zh{一样} Verbe. \zh{中国和喀麦隆一样重视接种疫苗。} (Zhōngguó hé Kāmàilóng yīyàng zhòngshì jiēzhòng yìmiáo.) \zh{重视} est un verbe, placé après \zh{一样}.
\item (c) « Beaucoup plus » → Adj. + \zh{得多}. \zh{蚊帐比药便宜得多。} (Wénzhàng bǐ yào piányi de duō.) Attention : \zh{得} (degré), pas \zh{的}.
\end{enumerate}
\end{exoresolu}

\begin{methode}[Mobiliser le lexique : Collocations verbales + Classificateurs]
Pour parler santé avec justesse, associe chaque nom à son verbe « partenaire » et son classificateur :
\begin{center}
\begin{tabularx}{\linewidth}{|X|X|X|}
\hline
\rowcolor{poporangeL} \textbf{Nom (Objet/Maladie)} & \textbf{Verbe + Classificateur correct} & \textbf{Exemple} \\ \hline
\zh{疟疾 / 霍乱 / 病} (Maladies) & \zh{得} (dé, attraper) + \zh{种} (zhǒng) & \zh{他得了一种怪病。} (Tā dé le yī zhǒng guài bìng.) \\ \hline
\zh{疫苗} (Vaccin) & \zh{打} (dǎ) / \zh{接种} (jiēzhòng) + \zh{针} (zhēn, dose) / \zh{次} (cì, fois) & \zh{我打了两针疫苗。} (Wǒ dǎ le liǎng zhēn yìmiáo.) \\ \hline
\zh{蚊帐} (Moustiquaire) & \zh{用} (yòng) / \zh{挂} (guà, accrocher) + \zh{顶} (dǐng) & \zh{每家发一顶蚊帐。} (Měi jiā fā yī dǐng wénzhàng.) \\ \hline
\zh{医生 / 专家} (Médecin/Expert) & \zh{看} (kàn, consulter) / \zh{请} (qǐng, faire venir) + \zh{位} (wèi, respect) & \zh{请一位专家来看看。} (Qǐng yī wèi zhuānjiā lái kànkan.) \\ \hline
\zh{医院 / 诊所} (Hôpital/Clinique) & \zh{去} (qù) / \zh{住} (zhù, hospitalisé) + \zh{家} (jiā, établissement) / \zh{所} (suǒ) & \zh{去一家大医院。} (Qù yī jiā dà yīyuàn.) \\ \hline
\end{tabularx}
\end{center}
\end{methode}

\begin{methode}[Argumenter et nuancer : Le « Kit débat » en 5 étapes]
Pour ton exposé ou ton message écrit (Forum Santé Jeunesse) :
\begin{enumerate}
\item \cle{Ouvrir} : \zh{我认为……} (Wǒ rènwéi...), \zh{我觉得……} (Wǒ juéde...), \zh{在我看来，……} (Zài wǒ kànlái...).
\item \cle{Justifier (Cause → Conséquence)} : \zh{因为……，所以……} (Yīnwèi..., suǒyǐ...). \emph{NB : En chinois, garder les deux conjonctions est naturel et correct.}
\zh{因为蚊子传播疟疾，所以我们应该用蚊帐。}
\item \cle{Nuancer (Concession)} : \zh{虽然……，但是/可是……} (Suīrán..., dànshì/kěshì...).
\zh{虽然蚊帐不便宜，但是生病买药更贵。}
\item \cle{Proposer (Recommandation)} : \zh{应该} (yīnggāi), \zh{最好} (zuìhǎo), \zh{建议} (jiànyì + V).
\zh{政府应该给学校发免费蚊帐。}
\item \cle{Conclure} : \zh{总之，……} (Zǒngzhī...), \zh{所以，我觉得……} (Suǒyǐ, wǒ juéde...).
\end{enumerate}
\end{methode}

\begin{exoresolu}[Mini-argumentation rédigée : Sujet « Moustiquaires gratuites à l'école ? »]
Énoncé — Prends position en 4-5 phrases (Opinion, Cause, Nuance, Proposition, Conclusion).
\tcblower
Solution détaillée — Modèle :
\begin{enumerate}
\item \zh{我认为政府应该给所有中小学发免费的蚊帐。} (Wǒ rènwéi zhèngfǔ yīnggāi gěi suǒyǒu zhōngxiǎoxué fā miǎnfèi de wénzhàng.) \emph{Ouverture + \zh{应该} + \zh{免费的} (adj + \zh{的}).}
\item \zh{因为疟疾是喀麦隆最常见的病之一，每年很多学生都得疟疾，影响学习。} (Yīnwèi nüèjí shì Kāmàilóng zuì chángjiàn de bìng zhī yī, měi nián hěn duō xuésheng dōu dé nüèjí, yǐngxiǎng xuéxí.) \emph{\zh{最……之一} = l'un des plus... ; \zh{影响} = affecter.}
\item \zh{虽然发蚊帐要花钱，但是预防比治疗划算得多。} (Suīrán fā wénzhàng yào huāqián, dànshì yùfáng bǐ zhìliáo huásuàn de duō.) \emph{Nuance + \zh{划算} (rentable) + \zh{得多}.}
\item \zh{所以，建议学校配合卫生部门，定期检查蚊帐使用情况。} (Suǒyǐ, jiànyì xuéxiào pèihé wèishēng bùmén, dìngqī jiǎnchá wénzhàng shǐyòng qíngkuàng.) \emph{Proposition concrète : \zh{配合} (coopérer avec), \zh{定期} (régulièrement).}
\item \zh{总之，保护学生健康是最重要的。} (Zǒngzhī, bǎohù xuésheng jiànkāng shì zuì zhòngyào de.) \emph{Conclusion.}
\end{enumerate}
\end{exoresolu}

\begin{methode}[Traduire (Thème) : 3 contrôles anti-calque]
\begin{enumerate}
\item \textbf{Ordre SUJET — LIEU/TEMPS — VERBE — COMPL.} : « On se vaccine à l'hôpital » → \zh{人们在医院接种疫苗。} (Pas *\zh{人们接种疫苗在医院}).
\item \textbf{Finalité} : « Pour prévenir... » → \zh{为了} (wèile) + V (en tête) OU \zh{来} (lái) + V (devant verbe but). \zh{为了预防霍乱，要喝干净的水。}
\item \textbf{Mots-outils 的 / 得 / 地} :
\begin{itemize}
\item \zh{的} : NOM qualifié (Adj/Nom + \zh{的} + N) → \zh{干净的水}.
\item \zh{得} : DEGRÉ après Verbe/Adj → \zh{跑得快}, \zh{便宜得多}.
\item \zh{地} : MANIÈRE devant Verbe (Adv + \zh{地} + V) → \zh{认真地学习}.
\end{itemize}
\end{enumerate}
\end{methode}

\begin{exoresolu}[Thème guidé : Français → Chinois]
Énoncé — Traduis : (a) « Au Cameroun, le paludisme est plus fréquent pendant la saison des pluies. » (b) « Pour prévenir le choléra, il faut boire de l'eau propre et se laver souvent les mains. »
\tcblower
Solution détaillée :
\begin{enumerate}
\item (a) Sujet \zh{疟疾}, Lieu/Temps \zh{在喀麦隆，雨季的时候} (devant verbe). « Plus fréquent » → \zh{更常见} (gèng chángjiàn).
\zh{在喀麦隆，雨季的时候疟疾更常见。} (Zài Kāmàilóng, yǔjì de shíhou nüèjí gèng chángjiàn.)
\item (b) Finalité \zh{为了预防霍乱} (tête). « Il faut » → \zh{要} (yào) / \zh{应该}. Coordination : \zh{喝干净的水，勤洗手} (juxtaposition ou \zh{和}).
\zh{为了预防霍乱，我们要喝干净的水，勤洗手。} (Wèile yùfáng huòluàn, wǒmen yào hē gānjìng de shuǐ, qín xǐ shǒu.) \emph{\zh{勤} (souvent) devant verbe, sans \zh{很}.}
\end{enumerate}
\end{exoresolu}

\begin{attention}[Les 5 erreurs fatales au Bac (Check-list de relecture)]
\begin{enumerate}
\item \cle{« 很 » après \zh{比}} : \zh{疟疾比感冒很严重} ✗ → \zh{疟疾比感冒严重} ✓ / \zh{疟疾比感冒更严重} ✓.
\item \cle{Confusion 的 / 得 / 地} : \zh{便宜的多} ✗ → \zh{便宜得多} ✓ ; \zh{干净得水} ✗ → \zh{干净的水} ✓ ; \zh{快快的跑} ✗ → \zh{快快地跑} ✓.
\item \cle{Ordre français (Lieu/Temps APRÈS verbe)} : \zh{我接种疫苗在医院} ✗ → \zh{我在医院接种疫苗} ✓.
\item \cle{Classificateurs : \zh{二个} / \zh{一个病} / \zh{一个医生}} : \zh{两个国家} ✓ (\zh{两} devant classif.) ; \zh{一种病} ✓ ; \zh{一位医生} ✓ (respect).
\item \cle{Tons \& Caractères proches} : \zh{疟疾} (nüèjí, 4+2) ≠ \zh{痢疾} (lìjí) ; \zh{预防} (yùfáng) ≠ \zh{予防} (yúfáng, archaïque/erreur) ; \zh{疫苗} (yìmiáo, clé \zh{疒}) ; \zh{免疫} (miǎnyì) ≠ \zh{免疫力} (miǎnyìlì). Un ton faux = 0 pt à l'item.
\end{enumerate}
\end{attention}

\courssub{Atelier oral : Jeu de rôle \& Débat guidé « Forum Santé Jeunesse »}

\begin{experience}[Activité 1 : Jeu de rôle « Au dispensaire » (Binôme, 5 min)]
\textbf{Rôle A (Infirmier/ère)} : Accueille, demande symptômes (\zh{哪里不舒服？}), donne conseils prévention (\zh{应该...}, \zh{注意...}), prescrit (\zh{开点药}).
\textbf{Rôle B (Patient/Élève)} : Décrit symptômes (\zh{发烧、头疼、咳嗽}), pose questions (\zh{严重吗？}, \zh{要打针吗？}), note conseils.
\textbf{Contrainte} : Utiliser au moins 5 mots du lexique + 1 structure \zh{比} / \zh{一样} / \zh{因为……所以……}.
\textbf{Exemple réplique} : \zh{护士：你发烧三天了，比普通感冒严重，最好去医院查一下疟疾。} (Hùshì: Nǐ fāshāo sān tiān le, bǐ pǔtōng gǎnmào yánzhòng, zuìhǎo qù yīyuàn chá yīxià nüèjí.)
\end{experience}

\begin{experience}[Activité 2 : Débat structuré « Moustiquaires gratuites à l'école ? » (Groupes de 4, 10 min)]
\textbf{Phase 1 (2 min)} : Préparation individuelle — note 3 arguments POUR / CONTRE avec structures \zh{因为……所以……} / \zh{虽然……但是……}.
\textbf{Phase 2 (6 min)} : Débat oral. Un modérateur distribue la parole. Objectif : aboutir à une \zh{共识} (gòngshí, consensus) ou lister \zh{分歧} (fēnqí, désaccords).
\textbf{Phase 3 (2 min)} : Rapporteur résume en 3 phrases : \zh{我们组认为……因为……所以建议……。}
\textbf{Grille d'évaluation (Peer-assessment)} : Prononciation (tons) / Vocabulaire précis / Structures variées / Cohérence logique / Interaction.
\end{experience}

\courssub{Production écrite intégrée \& Évaluation formative}

\begin{exercice}[Évaluation formative — 30 min]
\textbf{Consigne} : Traite les 3 exercices. Écris les caractères, le pinyin (avec tons) et la traduction française.

\textbf{Ex. 1 — Vocabulaire \& Classificateurs (6 pts)} Complète avec le bon mot ET le bon classificateur : \zh{种 / 顶 / 位 / 针 / 家}.
(a) 我买了两 \_\_\_\_\_ \zh{蚊帐}。 (b) 这 \_\_\_\_\_ \zh{病} 很少见。 (c) 请一 \_\_\_\_\_ \zh{专家} 来看看。 (d) 孩子打了三 \_\_\_\_\_ \zh{疫苗}。 (e) 附近有两 \_\_\_\_\_ \zh{医院}。

\textbf{Ex. 2 — Grammaire : Comparaison \& Argumentation (8 pts)}
(a) Réécris avec \zh{比} : « Le choléra est plus dangereux que le rhume. » (\zh{霍乱} huòluàn, \zh{危险} wēixiǎn).
(b) Réécris avec \zh{一样} : « La prévention est aussi importante que le traitement. » (\zh{预防} yùfáng, \zh{治疗} zhìliáo, \zh{重要} zhòngyào).
(c) Traduis : « Bien que les vaccins soient gratuits, certains parents ne font pas vacciner leurs enfants. » (\zh{虽然……但是……}, \zh{免费} miǎnfèi, \zh{家长} jiāzhǎng, \zh{孩子} háizi).
(d) Complète : \zh{为了预防登革热，我们 \_\_\_\_\_ 清理积水 (jīshuǐ, eau stagnante)。} (\zh{应该} / \zh{要} / \zh{最好} — choisis et justifie).

\textbf{Ex. 3 — Expression écrite guidée (12 pts)}
Sujet : \zh{“如何在校园预防传染病”} (Comment prévenir les maladies infectieuses à l'école).
Rédige un court message (60-80 caractères) pour l'affiche du « Forum Santé Jeunesse ». Structure imposée :
1. Titre accrocheur. 2. Un danger identifié (\zh{比如……}). 3. Deux mesures concrètes (\zh{第一，……；第二，……}). 4. Appel à l'action (\zh{大家一起行动吧！}).
\end{exercice}

\begin{corrige}[Évaluation formative]
\textbf{Ex. 1} (a) \zh{顶} (dǐng) — \zh{两顶蚊帐}. (b) \zh{种} (zhǒng) — \zh{这种病}. (c) \zh{位} (wèi) — \zh{一位专家}. (d) \zh{针} (zhēn) — \zh{三针疫苗}. (e) \zh{家} (jiā) / \zh{所} (suǒ) — \zh{两家医院} / \zh{两所医院}.

\textbf{Ex. 2}
(a) \zh{霍乱比感冒危险。} (Huòluàn bǐ gǎnmào wēixiǎn.) — Pas de \zh{很}.
(b) \zh{预防和治疗一样重要。} (Yùfáng hé zhìliáo yīyàng zhòngyào.) — Verbe/Adj après \zh{一样}.
(c) \zh{虽然疫苗免费，但是有些家长不给孩子打疫苗。} (Suīrán yìmiáo miǎnfèi, dànshì yǒuxiē jiāzhǎng bù gěi háizi dǎ yìmiáo.) — \zh{有些} (certains), \zh{给} (à/pour) introduit le bénéficiaire.
(d) \zh{应该} (yīnggāi) ou \zh{要} (yào). \zh{最好} possible mais plus conseillé qu'obligation. \zh{为了预防登革热，我们应该清理积水。}

\textbf{Ex. 3 — Modèle de production (≈ 75 caractères)}
\zh{校园防病，人人有责！} (Xiàoyuán fángbìng, rénrén yǒu zé!) \\
\zh{比如流感、结核病容易在宿舍传播。} (Bǐrú liúgǎn, jiéhébìng róngyì zài sùshè chuánbō.) \\
\zh{第一，保持教室通风，勤洗手、戴口罩。} (Dì yī, bǎochí jiàoshì tōngfēng, qín xǐ shǒu, dài kǒuzhào.) \\
\zh{第二，生病及时去医务室，不带病上课。} (Dì èr, shēngbìng jíshí qù yīwùshì, bù dàibìng shàngkè.) \\
\zh{大家一起行动，健康校园靠我们！} (Dàjiā yīqǐ xíngdòng, jiànkāng xiàoyuán kào wǒmen!)
\end{corrige}

\begin{aretenir}[Bilan de la leçon ZH16 — Check-list compétences]
\begin{itemize}
\item \textbf{Socio-culturel} : Je sais citer 3 maladies endémiques Cameroun/Chine, expliquer 2 différences majeures (statut paludisme, rôle MTC/mouvement patriotique) et 1 point commun (prévention prioritaire).
\item \textbf{Vocabulaire} : Je maîtrise 20 mots-clés (maladies, symptômes, prévention), leurs classificateurs (\zh{种/顶/位/针/家}) et collocations (\zh{得病, 打疫苗, 用蚊帐, 注意卫生}).
\item \textbf{Grammaire} : J'utilise correctement \zh{比} (sans \zh{很}), \zh{一样}, \zh{得多}, \zh{因为……所以……}, \zh{虽然……但是……}, \zh{为了……}, ordre Lieu/Temps-Verbe, \zh{的/得/地}.
\item \textbf{Écrit} : Je traduis 2 phrases thème/version sans calque ; j'écris un message d'affiche structuré (60-80 car.).
\item \textbf{Oral} : Je joue un dialogue dispensaire ; je débats 3 min avec arguments structurés.
\end{itemize}
\end{aretenir}

\newpage

\coursec{Exercices}

\courssub{Je vérifie mes connaissances}

\begin{exercice}[QCM : vocabulaire de la santé]
Pour chaque question, entoure la bonne réponse.

\begin{enumerate}
\item Le paludisme se dit en chinois :
\begin{enumerate}
\item \zh{肺结核} \item \zh{疟疾} \item \zh{登革热}
\end{enumerate}
\item \zh{预防} (yùfáng) signifie :
\begin{enumerate}
\item soigner \item prévenir \item guérir
\end{enumerate}
\item Un moustique se dit :
\begin{enumerate}
\item \zh{蚊子} \item \zh{苍蝇} \item \zh{虫子}
\end{enumerate}
\item \zh{医疗队} (yīliáo duì) désigne :
\begin{enumerate}
\item un hôpital \item une équipe médicale \item une pharmacie
\end{enumerate}
\end{enumerate}
\end{exercice}

\begin{exercice}[Vrai ou faux — justifie ta réponse]
Indique si chaque affirmation est vraie (V) ou fausse (F), puis justifie en une phrase en français.

\begin{enumerate}
\item Le paludisme est transmis par l'eau sale.
\item La moustiquaire imprégnée est un moyen de prévention du paludisme.
\item La Chine a complètement éliminé le paludisme sur son territoire.
\item La tuberculose se soigne sans traitement médical.
\end{enumerate}
\end{exercice}

\begin{exercice}[Complète le vocabulaire]
Complète chaque phrase avec le mot chinois qui convient : \zh{发烧}、\zh{医院}、\zh{疫苗}、\zh{症状}、\zh{传染}。

\begin{enumerate}
\item 他头疼，还\_\_\_\_\_\_，应该去看医生。(Il a mal à la tête et de la fièvre.)
\item 疟疾是一种会\_\_\_\_\_\_的病。(Le paludisme est une maladie qui se transmet.)
\item 打\_\_\_\_\_\_可以预防很多疾病。(Les vaccins préviennent beaucoup de maladies.)
\item 发烧和咳嗽是这种病的\_\_\_\_\_\_。(Fièvre et toux sont les symptômes de cette maladie.)
\item 病人应该马上去\_\_\_\_\_\_。(Le malade doit aller tout de suite à l'hôpital.)
\end{enumerate}
\end{exercice}

\begin{attention}[Le point de langue : « contagieux »]
Le paludisme ne se transmet \emph{pas} directement d'homme à homme : il passe par le moustique. En chinois, on peut dire \zh{疟疾是一种会传染的病} au sens large (maladie transmissible), mais pour être précis on dira : \zh{疟疾通过蚊子传染。} (Nüèjí tōngguò wénzi chuánrǎn.) « Le paludisme se transmet par les moustiques. » Retiens la structure : \cle{通过 + moyen + verbe} = « par le biais de ».
\end{attention}

\begin{exercice}[Associe caractères et pinyin]
Relie chaque mot en caractères à son pinyin.

\begin{center}
\begin{tabularx}{\linewidth}{|X|X|}
\hline
\rowcolor{popblueL} Caractères & Pinyin \\
\hline
1. 疟疾 & a. fèijiéhé \\
\hline
2. 登革热 & b. yùfáng \\
\hline
3. 肺结核 & c. nüèjí \\
\hline
4. 预防 & d. dēnggérè \\
\hline
5. 治疗 & e. zhìliáo \\
\hline
\end{tabularx}
\end{center}
\end{exercice}

\begin{attention}[Piège de prononciation : 结核]
Le caractère \zh{结} se prononce \textbf{jié} (2\up{e} ton) dans \zh{肺结核} (fèijiéhé, tuberculose) et \zh{结果} (jiéguǒ, résultat), mais \textbf{jiē} (1\up{er} ton) dans \zh{结实} (jiēshi, solide). Attention aussi à \zh{疟} : \textbf{nüè} avec ü + 4\up{e} ton, et non « nuè » ni « nüě ».
\end{attention}

\courssub{Je m'entraîne}

\begin{exercice}[Traduis en français]
Traduis les phrases suivantes en français correct.

\begin{enumerate}
\item \zh{喀麦隆的疟疾比中国多得多。} (Kāmàilóng de nüèjí bǐ Zhōngguó duō de duō.)
\item \zh{由于蚊子很多，因此我们应该用蚊帐。} (Yóuyú wénzi hěn duō, yīncǐ wǒmen yīnggāi yòng wénzhàng.)
\item \zh{中国不但派了医疗队，而且送了医疗设备。} (Zhōngguó búdàn pài le yīliáo duì, érqiě sòng le yīliáo shèbèi.)
\item \zh{虽然预防很重要，但是治疗也不能忘记。} (Suīrán yùfáng hěn zhòngyào, dànshì zhìliáo yě bù néng wàngjì.)
\end{enumerate}
\end{exercice}

\begin{exercice}[Traduis en chinois (thème)]
Traduis les phrases suivantes en chinois (caractères obligatoires).

\begin{enumerate}
\item Le paludisme est plus dangereux que la grippe.
\item Parce qu'il pleut beaucoup à Douala, il y a beaucoup de moustiques.
\item Non seulement il faut dormir sous la moustiquaire, mais il faut aussi se faire vacciner.
\item Bien que la Chine soit loin, elle aide le Cameroun dans la santé.
\end{enumerate}
\end{exercice}

\begin{exercice}[Texte à trous : mots-outils et structures]
Complète le texte avec les mots de la liste : \zh{比}、\zh{由于}、\zh{因此}、\zh{不但}、\zh{而且}、\zh{了}。

\begin{textebox}[La santé dans mon quartier]
\_\_\_\_\_\_ 我们区的蚊子很多，\_\_\_\_\_\_ 很多人得了疟疾。\\
(Yóuyú wǒmen qū de wénzi hěn duō, yīncǐ hěn duō rén dé le nüèjí.)\\
疟疾\_\_\_\_\_\_感冒危险得多。\\
(Nüèjí bǐ gǎnmào wēixiǎn de duō.)\\
政府\_\_\_\_\_\_发了蚊帐，\_\_\_\_\_\_组织了宣传活动。\\
(Zhèngfǔ búdàn fā le wénzhàng, érqiě zǔzhī le xuānchuán huódòng.)\\
去年，我们区建\_\_\_\_\_\_一个新的卫生中心。\\
(Qùnián, wǒmen qū jiàn le yí gè xīn de wèishēng zhōngxīn.)
\end{textebox}
\end{exercice}

\begin{exercice}[Remets les mots dans l'ordre]
Remets les mots dans l'ordre pour former des phrases correctes, puis traduis-les en français.

\begin{enumerate}
\item 疟疾 / 比 / 登革热 / 在喀麦隆 / 常见 / 更
\item 由于 / 下雨 / 经常 / 因此 / 很多 / 蚊子 / 有
\item 医疗队 / 中国 / 不但 / 来了 / 而且 / 带来 / 药品 / 了
\item 应该 / 我们 / 每天 / 蚊帐 / 用 / 睡觉 / 的时候
\end{enumerate}
\end{exercice}

\courssub{Je me prépare au Bac}

\begin{exercice}[Compréhension écrite — type Bac]
Lis le texte suivant, puis réponds aux questions.

\begin{textebox}[La coopération médicale entre la Chine et l'Afrique]
中国医疗队来喀麦隆已经很多年了。\\
(Zhōngguó yīliáo duì lái Kāmàilóng yǐjīng hěn duō nián le.)\\
他们不但在雅温得和杜阿拉的大医院工作，而且也去农村帮助病人。\\
(Tāmen búdàn zài Yǎwēndé hé Dù'ālā de dà yīyuàn gōngzuò, érqiě yě qù nóngcūn bāngzhù bìngrén.)\\
由于疟疾是喀麦隆最常见的疾病之一，因此中国医生带来了药品和设备，还培训了当地医生。\\
(Yóuyú nüèjí shì Kāmàilóng zuì chángjiàn de jíbìng zhī yī, yīncǐ Zhōngguó yīshēng dài lái le yàopǐn hé shèbèi, hái péixùn le dāngdì yīshēng.)\\
虽然两个国家很远，但是在健康方面，他们的合作越来越多。\\
(Suīrán liǎng ge guójiā hěn yuǎn, dànshì zài jiànkāng fāngmiàn, tāmen de hézuò yuè lái yuè duō.)
\end{textebox}

\textbf{Questions :}

\begin{enumerate}
\item \textbf{QCM vocabulaire.} Choisis la bonne traduction : a) \zh{医疗队} : hôpital / équipe médicale / pharmacie ; b) \zh{培训} : punir / former / payer ; c) \zh{合作} : coopération / compétition / construction.
\item \textbf{Vrai ou faux} (justifie en citant le texte en chinois) : a) Les médecins chinois travaillent seulement dans les grandes villes. b) Le paludisme est une maladie fréquente au Cameroun. c) La coopération sino-camerounaise diminue.
\item \textbf{Questions ouvertes} (réponds en français) : a) Où les équipes médicales chinoises travaillent-elles ? b) Quelles sont les trois formes d'aide mentionnées dans le texte ?
\item \textbf{Traduction.} Traduis en français les deux phrases suivantes : a) \zh{由于疟疾是喀麦隆最常见的疾病之一，因此中国医生带来了药品和设备。} b) \zh{虽然两个国家很远，但是在健康方面，他们的合作越来越多。}
\end{enumerate}
\end{exercice}

\begin{exercice}[Thème et version — type Bac]
\textbf{A. Thème.} Traduis en chinois (caractères obligatoires) :

\begin{enumerate}
\item La tuberculose est plus difficile à soigner que le paludisme.
\item Parce que les malades ne finissent pas leur traitement, la maladie revient.
\item Non seulement la Chine envoie des médecins, mais elle forme aussi des infirmiers camerounais.
\end{enumerate}

\textbf{B. Version.} Traduis en français :

\begin{enumerate}
\item \zh{预防比治疗便宜，也比治疗容易。} (Yùfáng bǐ zhìliáo piányi, yě bǐ zhìliáo róngyì.)
\item \zh{由于政府发了免费的蚊帐，因此得疟疾的人少了。} (Yóuyú zhèngfǔ fā le miǎnfèi de wénzhàng, yīncǐ dé nüèjí de rén shǎo le.)
\end{enumerate}
\end{exercice}

\begin{exercice}[Expression écrite — type Bac]
\textbf{Sujet.} Tu es délégué(e) santé de ton lycée. Rédige en chinois une lettre ouverte au Ministre de la Santé Publique (environ \cle{150 caractères}) pour proposer \cle{trois mesures} inspirées de l'expérience chinoise afin de lutter contre le paludisme au Cameroun.

\textbf{Contraintes obligatoires :}
\begin{itemize}
\item Utilise au moins \cle{5 mots} du glossaire de la leçon (par exemple : \zh{疟疾}、\zh{预防}、\zh{蚊帐}、\zh{疫苗}、\zh{医疗队}、\zh{卫生}).
\item Emploie au moins \cle{3 structures} de comparaison ou d'argumentation vues en classe : \zh{比}、\zh{由于……因此……}、\zh{不但……而且……}、\zh{虽然……但是……}.
\item Respecte la forme de la lettre : formule d'appel (\zh{尊敬的部长：}), corps du texte, formule de politesse, signature et date.
\end{itemize}

\textbf{Expression orale (préparation au Bac blanc).} Tire au sort une fiche « Maladie + Angle » (exemples : Dengue / prévention vectorielle ; Tuberculose / observance du traitement ; Paludisme / moustiquaires). Prépare 5 minutes, puis présente 2 minutes : définition de la maladie, situation comparée Cameroun / Chine, mesure prioritaire, avis personnel. Évaluation sur grille de 20 points (contenu 6, vocabulaire 4, structures 4, prononciation et tons 4, aisance 2).
\end{exercice}

\begin{methode}[Réussir la lettre ouverte en chinois]
\begin{enumerate}
\item \textbf{Forme d'abord.} Écris l'appel \zh{尊敬的部长：} en haut, termine par une formule de politesse (\zh{谢谢您！} ou \zh{此致敬礼！}), puis la signature (\zh{学生代表} + ton prénom) et la date à droite.
\item \textbf{Plan en trois mesures.} Annonce-les avec \zh{第一}、\zh{第二}、\zh{第三} : par exemple distribution de moustiquaires, campagnes de sensibilisation, formation de personnels de santé avec l'aide des équipes médicales chinoises.
\item \textbf{Structures imposées.} Place un comparatif (\zh{预防比治疗便宜}), une cause (\zh{由于……因此……}) et une addition (\zh{不但……而且……}). Une seule occurrence bien construite de chaque suffit.
\item \textbf{Vérifie.} Compte tes caractères (environ 150), souligne tes 5 mots du glossaire, relis les tons et l'ordre des mots (complément de temps et de lieu avant le verbe).
\end{enumerate}
\end{methode}

\begin{exemplebox}[Modèle de lettre (début)]
\zh{尊敬的部长：} (Zūnjìng de bùzhǎng :) « Monsieur le Ministre, »\\
\zh{我是中学的学生代表。由于疟疾是喀麦隆最常见的疾病之一，因此我给您写这封信。} (Wǒ shì zhōngxué de xuéshēng dàibiǎo. Yóuyú nüèjí shì Kāmàilóng zuì chángjiàn de jíbìng zhī yī, yīncǐ wǒ gěi nín xiě zhè fēng xìn.) « Je suis le délégué des élèves du lycée. Comme le paludisme est l'une des maladies les plus fréquentes au Cameroun, je vous écris cette lettre. »\\
\zh{第一，政府不但应该发免费的蚊帐，而且应该组织宣传活动。} (Dì yī, zhèngfǔ búdàn yīnggāi fā miǎnfèi de wénzhàng, érqiě yīnggāi zǔzhī xuānchuán huódòng.) « Premièrement, le gouvernement devrait non seulement distribuer des moustiquaires gratuites, mais aussi organiser des campagnes de sensibilisation. »
\end{exemplebox}

\newpage

\coursec{Corrigés des exercices}

\begin{corrige}[Exercice 1 — QCM : vocabulaire de la santé]
\begin{enumerate}
\item \textbf{b)} \zh{疟疾} (nüèjí) = le paludisme. Les distracteurs : \zh{肺结核} (fèijiéhé) = la tuberculose ; \zh{登革热} (dēnggérè) = la dengue. Astuce : le caractère \zh{疟} contient la clé de la maladie \zh{疒}.
\item \textbf{b)} \zh{预防} (yùfáng) = prévenir. « Soigner » se dit \zh{治疗} (zhìliáo) ; « guérir » \zh{治好} (zhìhǎo). Ne pas confondre \zh{预防} (prévention, avant la maladie) et \zh{治疗} (traitement, après).
\item \textbf{a)} \zh{蚊子} (wénzi) = le moustique. \zh{苍蝇} (cāngying) = la mouche ; \zh{虫子} (chóngzi) = l'insecte en général.
\item \textbf{b)} \zh{医疗队} (yīliáo duì) = une équipe médicale. « Hôpital » = \zh{医院} (yīyuàn) ; « pharmacie » = \zh{药店} (yàodiàn). Le mot \zh{队} (duì) signifie « équipe, escouade ».
\end{enumerate}
\textbf{Point de vigilance :} le suffixe \zh{子} (zi) de \zh{蚊子} est atone ; ne pas lui donner de ton.
\end{corrige}

\begin{corrige}[Exercice 2 — Vrai ou faux]
\begin{enumerate}
\item \textbf{Faux.} Le paludisme est transmis par la piqûre d'un moustique infecté (l'anophèle femelle), pas par l'eau sale. Les maladies liées à l'eau insalubre sont plutôt le choléra (\zh{霍乱}, huòluàn) ou la typhoïde (\zh{伤寒}, shānghán).
\item \textbf{Vrai.} Dormir sous une moustiquaire imprégnée d'insecticide (\zh{蚊帐}, wénzhàng) est l'un des moyens de prévention les plus efficaces contre le paludisme.
\item \textbf{Vrai.} Après des décennies de lutte (assainissement, traitements, moustiquaires), la Chine a été certifiée exempte de paludisme par l'OMS en 2021 : il n'y a plus de transmission locale sur son territoire.
\item \textbf{Faux.} La tuberculose exige un traitement médical long (plusieurs mois d'antibiotiques) et suivi ; sans traitement complet, la maladie persiste, se transmet et peut devenir résistante.
\end{enumerate}
\textbf{Point de vigilance :} une justification en une phrase complète est exigée ; « vrai » ou « faux » seul ne vaut pas tous les points.
\end{corrige}

\begin{corrige}[Exercice 3 — Complète le vocabulaire]
\begin{enumerate}
\item 他头疼，还\textbf{发烧}，应该去看医生。(Tā tóu téng, hái fāshāo, yīnggāi qù kàn yīshēng.) — \zh{发烧} (fāshāo) = avoir de la fièvre ; \zh{还} (hái) = « en plus ».
\item 疟疾是一种会\textbf{传染}的病。(Nüèjí shì yì zhǒng huì chuánrǎn de bìng.) — \zh{传染} (chuánrǎn) = transmettre, être contagieux ; \zh{会} + verbe = « susceptible de ».
\item 打\textbf{疫苗}可以预防很多疾病。(Dǎ yìmiáo kěyǐ yùfáng hěn duō jíbìng.) — \zh{打疫苗} (dǎ yìmiáo) = se faire vacciner ; le verbe \zh{打} s'emploie aussi pour les injections.
\item 发烧和咳嗽是这种病的\textbf{症状}。(Fāshāo hé késòu shì zhè zhǒng bìng de zhèngzhuàng.) — \zh{症状} (zhèngzhuàng) = symptôme.
\item 病人应该马上去\textbf{医院}。(Bìngrén yīnggāi mǎshàng qù yīyuàn.) — \zh{医院} (yīyuàn) = hôpital ; \zh{马上} (mǎshàng) = immédiatement.
\end{enumerate}
\textbf{Point de vigilance :} chaque mot de la liste ne sert qu'une fois ; vérifier la cohérence du sens avant de recopier les caractères.
\end{corrige}

\begin{corrige}[Exercice 4 — Associe caractères et pinyin]
1 -- c (\zh{疟疾} nüèjí, paludisme) ; 2 -- d (\zh{登革热} dēnggérè, dengue) ; 3 -- a (\zh{肺结核} fèijiéhé, tuberculose pulmonaire) ; 4 -- b (\zh{预防} yùfáng, prévenir) ; 5 -- e (\zh{治疗} zhìliáo, soigner, traitement).

\textbf{Point de vigilance :} \zh{疟} se prononce \textbf{nüè} (ü + 4\up{e} ton) ; \zh{结} dans \zh{结核} se prononce \textbf{jié} (2\up{e} ton), jamais jiē ici. Le mot \zh{登革热} est une transcription phonétique de l'anglais \emph{dengue} : on le reconnaît au caractère \zh{热} (rè, chaleur, fièvre).
\end{corrige}

\begin{corrige}[Exercice 5 — Traduis en français]
\begin{enumerate}
\item « Le paludisme est beaucoup plus fréquent au Cameroun qu'en Chine. » Structure : A \zh{比} B + adjectif + \zh{得多} = « beaucoup plus … que ». On ne traduit pas \zh{的} entre \zh{喀麦隆} et \zh{疟疾} par « de » : c'est « le paludisme au Cameroun ».
\item « Comme il y a beaucoup de moustiques, nous devons utiliser des moustiquaires. » \zh{由于……因此……} = « puisque… donc… » ; en français, un seul connecteur suffit.
\item « La Chine a non seulement envoyé des équipes médicales, mais aussi fourni du matériel médical. » \zh{不但……而且……} = « non seulement… mais aussi… » ; \zh{了} marque des actions accomplies.
\item « Bien que la prévention soit très importante, il ne faut pas non plus oublier le traitement. » \zh{虽然……但是……} = « bien que… mais… » ; en français on ne garde qu'un seul des deux connecteurs.
\end{enumerate}
\textbf{Point de vigilance :} éviter le calque « bien que… mais… » en français : faute classique sanctionnée en version.
\end{corrige}

\begin{corrige}[Exercice 6 — Traduis en chinois (thème)]
\begin{enumerate}
\item \zh{疟疾比感冒更危险。} (Nüèjí bǐ gǎnmào gèng wēixiǎn.) — « La grippe » courante se dit \zh{感冒} (gǎnmào) ; \zh{更} renforce le comparatif \zh{比}. Ne jamais dire \zh{疟疾比感冒很危险} : avec \zh{比}, l'adjectif ne prend pas \zh{很}.
\item \zh{由于杜阿拉经常下雨，因此蚊子很多。} (Yóuyú Dù'ālā jīngcháng xiàyǔ, yīncǐ wénzi hěn duō.) — La cause \zh{由于} se place en tête ; \zh{因此} introduit la conséquence.
\item \zh{不但要睡在蚊帐里，而且还要打疫苗。} (Búdàn yào shuì zài wénzhàng lǐ, érqiě hái yào dǎ yìmiáo.) — « Dormir sous la moustiquaire » = \zh{睡在蚊帐里} ; « se faire vacciner » = \zh{打疫苗}.
\item \zh{虽然中国很远，但是它在健康方面帮助喀麦隆。} (Suīrán Zhōngguó hěn yuǎn, dànshì tā zài jiànkāng fāngmiàn bāngzhù Kāmàilóng.) — « Dans le domaine de la santé » = \zh{在健康方面} ; le complément \zh{在……方面} précède le verbe.
\end{enumerate}
\textbf{Point de vigilance :} en chinois, contrairement au français, on peut (et on préfère) garder les deux connecteurs \zh{虽然……但是……} et \zh{由于……因此……} dans la même phrase.
\end{corrige}

\begin{corrige}[Exercice 7 — Texte à trous : mots-outils et structures]
\begin{enumerate}
\item \zh{由于}我们区的蚊子很多，\zh{因此}很多人得了疟疾。 (Yóuyú … yīncǐ …) — cause + conséquence ; les deux mots vont obligatoirement ensemble ici.
\item 疟疾\zh{比}感冒危险得多。 (… bǐ …) — comparatif : A \zh{比} B + adjectif (+ \zh{得多}).
\item 政府\zh{不但}发了蚊帐，\zh{而且}组织了宣传活动。 (… búdàn … érqiě …) — addition de deux actions du même sujet.
\item 去年，我们区建\zh{了}一个新的卫生中心。 (… jiàn le …) — \zh{了} d'accomplissement placé juste après le verbe, action datée (\zh{去年}) et achevée.
\end{enumerate}
\textbf{Point de vigilance :} \zh{了} se place \emph{après le verbe}, pas en fin de phrase, quand l'objet est déterminé (\zh{一个新的卫生中心}).
\end{corrige}

\begin{corrige}[Exercice 8 — Remets les mots dans l'ordre]
\begin{enumerate}
\item \zh{在喀麦隆，疟疾比登革热更常见。} (Zài Kāmàilóng, nüèjí bǐ dēnggérè gèng chángjiàn.) « Au Cameroun, le paludisme est plus fréquent que la dengue. » Le complément de lieu \zh{在喀麦隆} se place en tête de phrase ; \zh{更} précède l'adjectif.
\item \zh{由于经常下雨，因此有很多蚊子。} (Yóuyú jīngcháng xiàyǔ, yīncǐ yǒu hěn duō wénzi.) « Comme il pleut souvent, il y a beaucoup de moustiques. » L'adverbe de fréquence \zh{经常} précède le verbe \zh{下雨}.
\item \zh{中国医疗队不但来了，而且带来了药品。} (Zhōngguó yīliáo duì búdàn lái le, érqiě dài lái le yàopǐn.) « L'équipe médicale chinoise est non seulement venue, mais a aussi apporté des médicaments. » \zh{不但} se place après le sujet commun.
\item \zh{我们每天睡觉的时候应该用蚊帐。} (Wǒmen měitiān shuìjiào de shíhou yīnggāi yòng wénzhàng.) « Chaque jour, au moment de dormir, nous devons utiliser une moustiquaire. » Ordre : sujet + temps (\zh{每天}) + complément \zh{睡觉的时候} + modal \zh{应该} + verbe.
\end{enumerate}
\textbf{Point de vigilance :} en chinois, les compléments de temps et de lieu se placent \emph{avant} le verbe principal ; c'est l'erreur n°1 des francophones.
\end{corrige}

\begin{corrige}[Exercice 9 — Compréhension écrite — type Bac]
\textbf{Barème indicatif : 20 points} (QCM 3 pts ; Vrai/Faux 6 pts ; questions ouvertes 5 pts ; traduction 6 pts).

\begin{enumerate}
\item a) équipe médicale ; b) former ; c) coopération. (1 pt par bonne réponse.)
\item a) \textbf{Faux} : le texte dit \zh{他们也去农村帮助病人} — ils vont aussi dans les campagnes aider les malades. b) \textbf{Vrai} : \zh{疟疾是喀麦隆最常见的疾病之一} — le paludisme est l'une des maladies les plus fréquentes au Cameroun. c) \textbf{Faux} : \zh{他们的合作越来越多} — leur coopération est de plus en plus importante. (1 pt par réponse + 1 pt par citation chinoise justificative.)
\item a) Elles travaillent dans les grands hôpitaux de Yaoundé et de Douala, et aussi dans les zones rurales. b) Les trois formes d'aide : l'envoi de médicaments (\zh{药品}), de matériel médical (\zh{设备}) et la formation des médecins locaux (\zh{培训当地医生}).
\item a) « Comme le paludisme est l'une des maladies les plus fréquentes au Cameroun, les médecins chinois ont apporté des médicaments et du matériel. » b) « Bien que les deux pays soient très éloignés, leur coopération dans le domaine de la santé est de plus en plus développée. »
\end{enumerate}
\textbf{Points de vigilance :} pour le Vrai/Faux, la justification doit \emph{citer le texte en chinois} ; en traduction, rendre \zh{之一} par « l'un(e) des… » et \zh{越来越} par « de plus en plus ».
\end{corrige}

\begin{corrige}[Exercice 10 — Thème et version — type Bac]
\textbf{Barème indicatif : 10 points} (1 pt par phrase de thème pour la structure, 1 pt pour l'exactitude des caractères ; 1,5 pt par phrase de version).

\textbf{A. Thème.}
\begin{enumerate}
\item \zh{肺结核比疟疾更难治疗。} (Fèijiéhé bǐ nüèjí gèng nán zhìliáo.) — « Difficile à soigner » = \zh{难治疗} ; le comparatif exige \zh{比} et interdit \zh{很} devant l'adjectif.
\item \zh{由于病人没有完成治疗，因此病又回来了。} (Yóuyú bìngrén méiyǒu wánchéng zhìliáo, yīncǐ bìng yòu huílái le.) — « Finir le traitement » = \zh{完成治疗} ; la négation d'une action non accomplie au passé se fait avec \zh{没有} (jamais \zh{不} devant un accompli) ; \zh{又} = « de nouveau ».
\item \zh{中国不但派医生，而且还培训喀麦隆的护士。} (Zhōngguó búdàn pài yīshēng, érqiě hái péixùn Kāmàilóng de hùshi.) — « Infirmier » = \zh{护士} (hùshi) ; \zh{不但} suit le sujet unique \zh{中国}.
\end{enumerate}

\textbf{B. Version.}
\begin{enumerate}
\item « La prévention est moins chère que le traitement, et aussi plus facile que le traitement. » Double comparatif avec \zh{比} ; \zh{也} = « aussi » ; \zh{便宜} (piányi) = bon marché.
\item « Comme le gouvernement a distribué des moustiquaires gratuites, les cas de paludisme ont diminué. » \zh{得疟疾的人少了} : littéralement « les personnes qui attrapent le paludisme sont devenues moins nombreuses » ; le \zh{了} final marque le changement d'état.
\end{enumerate}
\textbf{Points de vigilance :} en thème, toute phrase sans caractères (pinyin seul) n'est pas comptée ; en version, ne pas traduire mot à mot \zh{人少了} par « les gens sont petits » — \zh{少} (shǎo) = « peu nombreux ».
\end{corrige}

\begin{corrige}[Exercice 11 — Expression écrite — type Bac]
\textbf{Barème indicatif : 20 points} — respect de la forme épistolaire 4 pts ; trois mesures claires et argumentées 6 pts ; 5 mots du glossaire 3 pts ; 3 structures imposées correctes 4 pts ; exactitude des caractères et de la grammaire 3 pts.

\textbf{Proposition de rédaction modèle (environ 150 caractères) :}

\begin{textebox}[Lettre ouverte au Ministre de la Santé]
尊敬的部长：\\
(Zūnjìng de bùzhǎng :)\\
我是中学的学生代表。由于疟疾是喀麦隆最常见的疾病之一，因此我给您写这封信。\\
(Wǒ shì zhōngxué de xuéshēng dàibiǎo. Yóuyú nüèjí shì Kāmàilóng zuì chángjiàn de jíbìng zhī yī, yīncǐ wǒ gěi nín xiě zhè fēng xìn.)\\
第一，政府不但应该发免费的蚊帐，而且应该组织宣传活动，因为预防比治疗便宜。\\
(Dì yī, zhèngfǔ búdàn yīnggāi fā miǎnfèi de wénzhàng, érqiě yīnggāi zǔzhī xuānchuán huódòng, yīnwèi yùfáng bǐ zhìliáo piányi.)\\
第二，每个区都应该有卫生中心。\\
(Dì èr, měi ge qū dōu yīnggāi yǒu wèishēng zhōngxīn.)\\
第三，我们应该学习中国的经验，请中国医疗队培训我们的医生。\\
(Dì sān, wǒmen yīnggāi xuéxí Zhōngguó de jīngyàn, qǐng Zhōngguó yīliáo duì péixùn wǒmen de yīshēng.)\\
虽然这个问题很难，但是我们一起努力，一定会成功。\\
(Suīrán zhè ge wèntí hěn nán, dànshì wǒmen yìqǐ nǔlì, yídìng huì chénggōng.)\\
谢谢您！\\
(Xièxie nín !)\\
学生代表 玛丽\\
(Xuéshēng dàibiǎo Mǎlì)\\
二〇二五年三月\\
(Èr líng èr wǔ nián sān yuè)
\end{textebox}

\textbf{Traduction du modèle :} « Monsieur le Ministre, je suis le délégué des élèves du lycée. Comme le paludisme est l'une des maladies les plus fréquentes au Cameroun, je vous écris cette lettre. Premièrement, le gouvernement devrait non seulement distribuer des moustiquaires gratuites, mais aussi organiser des campagnes de sensibilisation, car la prévention coûte moins cher que le traitement. Deuxièmement, chaque quartier devrait avoir un centre de santé. Troisièmement, nous devrions tirer profit de l'expérience chinoise et demander aux équipes médicales chinoises de former nos médecins. Bien que ce problème soit difficile, si nous travaillons ensemble, nous réussirons certainement. Merci ! Le délégué des élèves, Marie. Mars 2025. »

\textbf{Points de vigilance :}
\begin{itemize}
\item L'appel \zh{尊敬的部长：} prend la \cle{ponctuation chinoise} « ：» et non le deux-points français.
\item Annoncer les mesures avec \zh{第一}、\zh{第二}、\zh{第三} : c'est la marque d'un texte structuré.
\item Vérifier que les trois structures imposées apparaissent au moins une fois chacune et sont \emph{complètes} (\zh{不但} sans \zh{而且} est une faute).
\item Compter les caractères : en dessous de 120 ou au-dessus de 180, on sort du contrat.
\item À l'oral, soigner les tons de \zh{疟疾} (nüèjí) et \zh{预防} (yùfáng) : ce sont les mots les plus souvent mal prononcés.
\end{itemize}
\end{corrige}

\newpage

\coursec{Fiche bilan}

\begin{popfigure}
\centering
\begin{tikzpicture}[mindmap, grow cyclic, every node/.style={concept, align=center, font=\small}, concept color=popblue!60, root concept/.append style={concept color=popblue, font=\small\bfseries}, level 1/.append style={level distance=3.2cm, sibling angle=51.4}, level 2/.append style={level distance=1.9cm, font=\scriptsize}]
\node[root concept, align=center]{Maladies endémiques\\Afrique -- Asie\\ZH16}
  child[concept color=poporange!70]{ node[align=center]{Lexique\\maladies} child { node[align=center]{疟疾 nüèji\\paludisme} } child { node[align=center]{霍乱 huòluàn\\choléra} } }
  child[concept color=popgreen!70]{ node {Symptômes} child { node[align=center]{发烧 fāshāo\\fièvre} } child { node[align=center]{咳嗽 késou\\toux} } }
  child[concept color=poppurple!70]{ node {Prévention} child { node[align=center]{预防 yùfáng\\prévenir} } child { node[align=center]{疫苗 yìmiáo\\vaccin} } }
  child[concept color=popgold!80]{ node {Comparer} child { node {比 bǐ} } child { node {和……一样} } }
  child[concept color=poppink!70]{ node {Argumenter} child { node {我认为 wǒ rènwéi} } child { node {因为……所以} } }
  child[concept color=popgreen!50]{ node {Cameroun} child { node {moustiquaires} } }
  child[concept color=poporange!50]{ node {Chine} child { node[align=center]{médecine\\traditionnelle} } };
\end{tikzpicture}
\legende{Carte mentale de la leçon ZH16 : maladies endémiques en Afrique et en Asie}
\end{popfigure}

\begin{bilan}
\textbf{1. Lexique clé à connaître par cœur}
\begin{center}
\begin{tabularx}{\linewidth}{|X|X|X|X|}
\hline
\rowcolor{popblueL} Caractères & Pinyin & Nature & Traduction \\
\hline
疾病 & jíbìng & n. & maladie \\
\hline
疟疾 & nüèji & n. & paludisme \\
\hline
霍乱 & huòluàn & n. & choléra \\
\hline
伤寒 & shānghán & n. & typhoïde \\
\hline
症状 & zhèngzhuàng & n. & symptôme \\
\hline
发烧 & fāshāo & v. & avoir de la fièvre \\
\hline
腹泻 & fùxiè & v. & avoir la diarrhée \\
\hline
传染 & chuánrǎn & v. & transmettre, contaminer \\
\hline
蚊子 & wénzi & n. & moustique \\
\hline
蚊帐 & wénzhàng & n. & moustiquaire \\
\hline
预防 & yùfáng & v. & prévenir \\
\hline
疫苗 & yìmiáo & n. & vaccin \\
\hline
卫生 & wèishēng & n./adj. & hygiène ; hygiénique \\
\hline
流行病 & liúxíngbìng & n. & épidémie \\
\hline
治疗 & zhìliáo & v. & soigner, traiter \\
\hline
\end{tabularx}
\end{center}

\textbf{2. Grammaire à connaître par cœur}
\begin{center}
\begin{tabularx}{\linewidth}{|X|X|X|}
\hline
\rowcolor{popblueL} Structure & Exemple & Traduction \\
\hline
A 比 B + adj. & 喀麦隆比中国热。Kāmàilóng bǐ Zhōngguó rè. & Le Cameroun est plus chaud que la Chine. \\
\hline
A 和 B 一样 + adj. & 预防和治疗一样重要。Yùfáng hé zhìliáo yíyàng zhòngyào. & Prévenir est aussi important que soigner. \\
\hline
因为……所以…… & 因为蚊子传染疟疾，所以我们要用蚊帐。 & Parce que les moustiques transmettent le paludisme, nous utilisons des moustiquaires. \\
\hline
我认为 / 我觉得 + phrase & 我认为预防比治疗更重要。Wǒ rènwéi yùfáng bǐ zhìliáo gèng zhòngyào. & Je pense que la prévention est plus importante que le traitement. \\
\hline
应该 / 得 (děi) + verbe & 我们应该喝干净的水。Wǒmen yīnggāi hē gānjìng de shuǐ. & Nous devons boire de l'eau propre. \\
\hline
\end{tabularx}
\end{center}

\textbf{3. Faits socioculturels à retenir}
\begin{itemize}
\item Au Cameroun, le \cle{paludisme} (\zh{疟疾}) est la maladie endémique la plus répandue ; la prévention passe par les moustiquaires imprégnées et l'assainissement.
\item En Chine, les grandes campagnes de santé publique et la vaccination ont fortement réduit les épidémies ; la médecine traditionnelle chinoise (\zh{中医} Zhōngyī) coexiste avec la médecine moderne.
\item Dans les deux pays, l'eau potable et l'hygiène (\zh{卫生}) sont les clés de la prévention du choléra et de la typhoïde.
\end{itemize}

\textbf{4. Méthodes express}
\begin{itemize}
\item \cle{Comparer} : repérer deux éléments (Cameroun / Chine), choisir un adjectif, appliquer 比 ou 和……一样.
\item \cle{Argumenter} : opinion (我认为) + raison (因为……所以……) + exemple concret.
\item \cle{Nuancer} : utiliser 但是 (mais), 不过 (cependant), 一方面……另一方面…… (d'un côté… de l'autre…).
\item \cle{Mini-exposé} : introduction (主题), 2 arguments, exemple Cameroun, exemple Chine, conclusion personnelle.
\end{itemize}
\end{bilan}

\begin{aretenir}[Checklist avant le Bac]
\begin{itemize}
\item[$\square$] Je connais le lexique des maladies : 疟疾, 霍乱, 伤寒, 流行病, 疾病.
\item[$\square$] Je connais le lexique des symptômes : 发烧, 咳嗽, 腹泻, 头疼.
\item[$\square$] Je connais le lexique de la prévention : 预防, 疫苗, 蚊帐, 卫生, 干净的水.
\item[$\square$] Je sais comparer avec 比 : A 比 B + adjectif (+ 多了 / 一点儿).
\item[$\square$] Je sais exprimer l'égalité avec 和……一样 / 不一样.
\item[$\square$] Je sais donner mon opinion avec 我认为 / 我觉得 sans oublier le sujet.
\item[$\square$] Je sais justifier avec 因为……所以…… (je peux n'utiliser qu'un seul des deux, jamais « parce que… donc » en français calqué).
\item[$\square$] Je sais nuancer avec 但是, 不过, 一方面……另一方面…….
\item[$\square$] Je peux citer un fait sanitaire du Cameroun et un fait de la Chine, de façon factuelle et respectueuse.
\item[$\square$] Je fais attention aux tons : yùfáng (预防) ne se prononce pas comme yìmiáo (疫苗).
\item[$\square$] Je place correctement les compléments : Sujet + 应该 / 得 + verbe + complément.
\item[$\square$] Je sais rédiger un paragraphe de 5 à 8 phrases : opinion, comparaison, exemple, conclusion.
\end{itemize}
\end{aretenir}

\findecours

\end{document}
